% production_exp.tex
% Stand-alone description of the production experiment.
% Can be embedded in main.tex via % production_exp.tex
% Stand-alone description of the production experiment.
% Can be embedded in main.tex via % production_exp.tex
% Stand-alone description of the production experiment.
% Can be embedded in main.tex via % production_exp.tex
% Stand-alone description of the production experiment.
% Can be embedded in main.tex via \input{production_exp}.
%
% Compile stand-alone:
%   pdflatex production_exp.tex
%
% Embed in main.tex:
%   \input{production_exp}   % (inside a \section{Experiment} environment)

\documentclass[12pt]{article}
\usepackage{CSPArticleStyle}
\usepackage{csquotes}
\usepackage[margin=1in]{geometry}
\usepackage{booktabs}
\usepackage{parskip}
\usepackage{amsmath}
\usepackage{times}
\usepackage[T1]{fontenc}
\usepackage[utf8]{inputenc}
\usepackage{xcolor}
\usepackage{tabularray}   % robust tables; fall back to tabular if unavailable
\usepackage{gb4e}         % linguistic examples
\noautomath

\title{Production Experiment: Eliciting Consensus-Marker Choice\\
       Across Perceived-Controversy and Goal-Strength Conditions}
\author{}
\date{}

\begin{document}

%% ============================================================
%% When embedding in main.tex, delete from \documentclass to
%% \begin{document} and delete \end{document} at the bottom.
%% ============================================================

\section{Experiment: Consensus-Marker Production}

\subsection{Overview}

The production experiment empirically estimates the decision thresholds
that govern the use of five German consensus markers by placing
participants in the speaker role.
Participants read short vignettes that independently varied
\emph{pragmatic perceived controversy} (\textsc{pc\textsubscript{prag}})
and the speaker's \emph{persuasive-goal strength} ($g$), then selected
which consensus marker they would most naturally use in the described
situation.
\emph{Propositional perceived controversy} (\textsc{pc\textsubscript{prop}})
is held constant within the experiment and operationalised as a continuous
topic-level covariate obtained from a separate norming study: for each
of the eight topic items, participants rated the degree to which the
proposition $p$ is considered established/accepted (0--100 scale);
the resulting means enter the regression analyses as a fixed covariate.
The resulting distribution of marker choices across the
$\textsc{pc}_\text{prag} \times g$ design, combined with the
continuous \textsc{pc\textsubscript{prop}} covariate, provides the
empirical basis for estimating the ordinal thresholds
$\mu_1 < \cdots < \mu_4$ of the noisy-threshold RSA model.

\subsection{Participants}

Participants will be recruited via Prolific and will be required to be
native speakers of German (language-at-home filter).
Data collection is ongoing; a target of $N = 80$ participants is planned
($n = 20$ per list across 4 lists), each completing 16 trials (8 critical + 8 filler).

\subsection{Materials}

\subsubsection{Markers}

The five consensus markers under investigation are, from weakest to
strongest presumed common-ground commitment:

\begin{enumerate}
  \item \textit{sofern ich wei\ss{}} (`as far as I know')
  \item \textit{wie du ja wei\ss t} (`as you (prt) know')
  \item \textit{wie wir wissen} (`as we know')
  \item \textit{ja} (modal particle; roughly `as you know / of course')
  \item \textit{bekanntlich} (`as is well known / notoriously')
\end{enumerate}

These markers were embedded in the sentential \emph{Mittelfeld}
(middle field), immediately after the finite verb, as illustrated in
example~(\ref{ex:frame}):

\begin{exe}
  \ex \label{ex:frame}
  Digitale Kompetenzen sind \textbf{[MARKER]}
  in der heutigen Arbeitswelt unverzichtbar.\\
  `Digital skills are \textbf{[MARKER]} indispensable in today's
  working world.'
\end{exe}

This syntactic position is licit for all five markers: modal particles
such as \textit{ja} canonically occupy the left edge of the Mittelfeld
\citep{thurmair1989modalpartikeln}, while the adverb \textit{bekanntlich}
and parenthetical clauses (\textit{wie wir wissen}, etc.) can appear
there as well (see \citealt{jacobs_semantics_1991} for discussion).
Each target sentence was followed by a consequent clause with the
recommended action $q$: \textit{Deshalb solltest du [q].}
(`Therefore, you should [q].')

\subsubsection{Topic Items}

Eight topic domains were used, each associated with a fixed proposition
$p$ (sentence frame) and consequent action $q$ (Table~\ref{tab:items}).
The first four items were used in earlier pilot work; items 5--8 were
added to reach the minimum of eight required for a complete Latin square.

\begin{table}[ht]
\centering
\caption{Topic items: sentence frame for $p$ (with marker slot [M]) and consequent $q$.}
\label{tab:items}
\small
\begin{tabular}{clll}
\toprule
\textbf{\#} & \textbf{Topic} & \textbf{Sentence frame} & \textbf{Consequent $q$} \\
\midrule
1 & Klimawandel &
  \textit{Der Klimawandel wird [M]} &
  auf Flugreisen verzichten \\
  & & \textit{durch menschliche Aktivit\"aten verursacht.} & \\[3pt]
2 & Ern\"ahrung &
  \textit{Eine \"uberwiegend pflanzliche Ern\"ahrung} &
  mehr pflanzliche Lebens- \\
  & & \textit{f\"ordert [M] die Gesundheit.} &
  mittel einbauen \\[3pt]
3 & Stadtverkehr &
  \textit{Ein ausgebautes Nahverkehrsnetz} &
  h\"aufiger auf \"offentliche \\
  & & \textit{entlastet [M] den Stadtverkehr deutlich.} &
  Verkehrsmittel umsteigen \\[3pt]
4 & Digitalisierung &
  \textit{Digitale Kompetenzen sind [M] in der} &
  regelm\"a\ss{}ig Weiterbildungs- \\
  & & \textit{heutigen Arbeitswelt unverzichtbar.} &
  angebote nutzen \\[3pt]
5 & Schlaf &
  \textit{Ausreichend Schlaf ist [M] entscheidend} &
  auf feste Schlafzeiten \\
  & & \textit{f\"ur die kognitive Leistungsf\"ahigkeit.} &
  achten \\[3pt]
6 & Plastik &
  \textit{Einwegplastik schadet [M] den} &
  konsequent auf Einweg- \\
  & & \textit{Meeres\"okosystemen erheblich.} &
  plastik verzichten \\[3pt]
7 & Sport &
  \textit{Regelm\"a\ss{}ige k\"orperliche Bewegung} &
  regelm\"a\ss{}ig Sport in \\
  & & \textit{senkt [M] das Risiko f\"ur} &
  den Alltag integrieren \\
  & & \textit{Herz-Kreislauf-Erkrankungen deutlich.} & \\[3pt]
8 & Lokalkauf &
  \textit{Das Kaufen bei lokalen H\"andlern} &
  \"ofter bei lokalen \\
  & & \textit{st\"arkt [M] die regionale Wirtschaft} &
  Gesch\"aften einkaufen \\
  & & \textit{nachhaltig.} & \\
\bottomrule
\end{tabular}
\end{table}

\noindent [M] denotes the marker slot (dropdown).

\subsubsection{Vignettes}

\paragraph{Propositional perceived controversy (\textsc{pc\textsubscript{prop}}).}
\textsc{pc\textsubscript{prop}} was \emph{not} manipulated within the
production experiment.
Instead, it was operationalised as a continuous, topic-level covariate
obtained from a separate norming study conducted prior to the production
experiment.
In that norming study, a different sample of native German speakers rated,
for each topic proposition $p$, the degree to which they considered $p$
accepted/established in (German) public discourse, on a 0--100 scale
($0 = $ `completely contested', $100 = $ `universally accepted').
The resulting per-topic means (e.g., \textit{Digitalisierung}: 85/100;
\textit{Lokalkauf}: 58/100) enter the regression analyses as a
continuous, centred predictor \textsc{pc\textsubscript{prop}}$_c$
(mean-centred, SD $\approx$ 1).
Context paragraphs in the production experiment therefore present
the proposition $p$ without any editorial evaluation, so that
participants' experience of propositional controversy is shaped
solely by their own prior knowledge and the norming-derived covariate
captures this variation in analysis.

\paragraph{Pragmatic perceived controversy (\textsc{pc\textsubscript{prag}}).}
\textsc{pc\textsubscript{prag}} was manipulated between instances of each item
by varying whether the speaker's social circle \emph{largely agreed} (low)
or \emph{held divergent views} (high) about $p$.
The two context versions for the \textit{Digitalisierung} item are:

\begin{itemize}
  \item $\textsc{pc}_\text{prag} = \text{low}$: \textit{Du sprichst mit einer
        Kollegin {\"u}ber Berufsausbildung und Weiterbildung.
        In eurem Team teilen die meisten die Ansicht,
        dass digitale Kompetenzen in der heutigen Arbeitswelt unverzichtbar sind.}
  \item $\textsc{pc}_\text{prag} = \text{high}$: \textit{Du sprichst mit einer
        Kollegin {\"u}ber Berufsausbildung und Weiterbildung.
        In eurem Team gehen die Meinungen dazu stark auseinander.}
\end{itemize}

\paragraph{Speaker goal strength ($g$).}
Goal strength was manipulated via a standardised goal instruction
appended after the context paragraph:

\begin{itemize}
  \item $g_\text{low}$: \textit{Du m\"ochtest das nur kurz anmerken.}
        (`You just want to make a brief note of this.')
  \item $g_\text{high}$: \textit{Es ist dir sehr wichtig, dass dein
        Gespr\"achspartner das akzeptiert.}
        (`It is very important to you that your interlocutor accepts this.')
\end{itemize}

\paragraph{Complete sample item (all four conditions).}
Examples (\ref{ex:cond1})--(\ref{ex:cond4}) show the full vignette for
the topic \textit{Digitalisierung}
(norming $\textsc{pc}_\text{prop}$ rating: 85/100;
sentence frame: \textit{Digitale Kompetenzen sind [M] in der heutigen
Arbeitswelt unverzichtbar.
Deshalb solltest du regelm\"a\ss{}ig digitale Weiterbildungsangebote nutzen.})
across all four conditions of the $2\times2$ design.

\begin{exe}
  \ex \label{ex:cond1}
  \textbf{Condition 0: $\textsc{pc}_\text{prag} = \text{low}$, $g = \text{low}$}\\
  \textit{Du sprichst mit einer Kollegin {\"u}ber Berufsausbildung und Weiterbildung.
  In eurem Team teilen die meisten die Ansicht,
  dass digitale Kompetenzen in der heutigen Arbeitswelt unverzichtbar sind.}\\
  \textit{Du m\"ochtest das nur kurz anmerken.}\\
  Digitale Kompetenzen sind \underline{\hspace{2.2cm}}
  in der heutigen Arbeitswelt unverzichtbar.

  \ex \label{ex:cond2}
  \textbf{Condition 1: $\textsc{pc}_\text{prag} = \text{low}$, $g = \text{high}$}\\
  \textit{Du sprichst mit einer Kollegin {\"u}ber Berufsausbildung und Weiterbildung.
  In eurem Team teilen die meisten die Ansicht,
  dass digitale Kompetenzen in der heutigen Arbeitswelt unverzichtbar sind.}\\
  \textit{Es ist dir sehr wichtig, dass dein Gespr\"achspartner das akzeptiert.}\\
  Digitale Kompetenzen sind \underline{\hspace{2.2cm}}
  in der heutigen Arbeitswelt unverzichtbar.

  \ex \label{ex:cond3}
  \textbf{Condition 2: $\textsc{pc}_\text{prag} = \text{high}$, $g = \text{low}$}\\
  \textit{Du sprichst mit einer Kollegin {\"u}ber Berufsausbildung und Weiterbildung.
  In eurem Team gehen die Meinungen dazu stark auseinander.}\\
  \textit{Du m\"ochtest das nur kurz anmerken.}\\
  Digitale Kompetenzen sind \underline{\hspace{2.2cm}}
  in der heutigen Arbeitswelt unverzichtbar.

  \ex \label{ex:cond4}
  \textbf{Condition 3: $\textsc{pc}_\text{prag} = \text{high}$, $g = \text{high}$}\\
  \textit{Du sprichst mit einer Kollegin {\"u}ber Berufsausbildung und Weiterbildung.
  In eurem Team gehen die Meinungen dazu stark auseinander.}\\
  \textit{Es ist dir sehr wichtig, dass dein Gespr\"achspartner das akzeptiert.}\\
  Digitale Kompetenzen sind \underline{\hspace{2.2cm}}
  in der heutigen Arbeitswelt unverzichtbar.
\end{exe}

\noindent In all four vignettes the sentence frame and consequent clause
(\textit{Deshalb solltest du regelm\"a\ss{}ig digitale
Weiterbildungsangebote nutzen.}) are identical; only the context
paragraph and goal instruction vary across conditions.
The marker slot [M] / \underline{\hspace{2.2cm}} was presented as a
drop-down menu in the actual experiment.

\subsection{Design}

\subsubsection{Factorial structure}

The experiment employed a \textbf{2 (pc\textsubscript{prag}: low vs.\ high)
$\times$ 2 ($g$: low vs.\ high) fully crossed design},
yielding 4 conditions.
\textsc{pc\textsubscript{prop}} is not a within-experiment factor; it enters
as a continuous topic-level covariate from the norming study.
With 8 topic items and a Latin square rotation across 4 lists, each
participant saw exactly \textbf{8 critical trials}---two items per condition---plus
8 filler trials, for \textbf{16 trials} in total
(Table~\ref{tab:design}).

\begin{table}[ht]
\centering
\caption{%
  The $2\times2$ design (4 conditions).
  \textsc{pc\textsubscript{prop}} is a continuous norming covariate, not a
  within-experiment factor.
  In each list, the 8 items are assigned to conditions via the Latin square
  rotation described in Section~\ref{subsub:counterbalancing}.}
\label{tab:design}
\small
\begin{tabular}{ccl}
\toprule
\textbf{pc\textsubscript{prag}} &
\textbf{$g$} &
\textbf{Condition label} \\
\midrule
low  & low  & (0) community agrees, casual \\
low  & high & (1) community agrees, pressing \\
high & low  & (2) community divided, casual \\
high & high & (3) community divided, pressing \\
\bottomrule
\end{tabular}
\end{table}

\subsubsection{Counterbalancing}
\label{subsub:counterbalancing}

The experiment used a \textbf{Latin square} to counterbalance item--condition
pairings across participants.
With $k = 4$ conditions and 8 items, a balanced assignment is achieved:
each of the 4 lists assigns item $i$ ($i = 0,\ldots,7$) to condition
$(i + l)\bmod 4$, where $l \in \{0,1,2,3\}$ is the list number.
The resulting assignment matrix (Table~\ref{tab:ls}) guarantees that
\begin{itemize}
  \item each participant sees every item exactly once;
  \item each participant sees each condition exactly twice (two items per
        condition), ensuring balanced coverage of the 2$\times$2 design;
  \item across all 4 lists, each item appears equally often in each
        condition (twice per condition over the full participant sample).
\end{itemize}

\begin{table}[ht]
\centering
\caption{Latin square assignment of conditions to items across 4 lists.
  Cell entries are condition indices (0--3) as defined in
  Table~\ref{tab:design}.
  Rows = lists ($l$); columns = items ($i$).}
\label{tab:ls}
\small
\begin{tabular}{ccccccccc}
\toprule
\textbf{List} &
\textbf{Item 1} & \textbf{Item 2} & \textbf{Item 3} & \textbf{Item 4} &
\textbf{Item 5} & \textbf{Item 6} & \textbf{Item 7} & \textbf{Item 8} \\
\midrule
0 & 0 & 1 & 2 & 3 & 0 & 1 & 2 & 3 \\
1 & 1 & 2 & 3 & 0 & 1 & 2 & 3 & 0 \\
2 & 2 & 3 & 0 & 1 & 2 & 3 & 0 & 1 \\
3 & 3 & 0 & 1 & 2 & 3 & 0 & 1 & 2 \\
\bottomrule
\end{tabular}
\end{table}

\noindent List assignment is determined automatically from the URL query
parameter \texttt{?list=N} (0--3); if absent, a random list is selected.
Each participant is thus within-subjects across items and sees all four
conditions (twice each). Condition effects are estimated from within-participant
variation across the 2$\times$2 design, controlling for by-participant
random intercepts.

In addition to the 8 critical trials, each participant completed 8
\textbf{filler trials} drawn from separate topic domains (train travel,
reading light, cooking, coffee, vacation, hydration, fresh air, noise).
Filler items used the same trial format but were excluded from the
main analysis; they serve to dilute the critical items and reduce
potential demand characteristics.

\subsubsection{Dependent variable}

The dependent variable was the marker chosen from a drop-down menu
displaying all five options in fixed order (as listed in
Section~\ref{sub:markers}).
No response time threshold was imposed, but response times were logged.
Recorded per trial: \texttt{trial\_id}, \texttt{topic}, \texttt{is\_filler},
\texttt{condition\_index}, \texttt{pc\_prag}, \texttt{g},
\texttt{selected\_marker}, \texttt{rt} (ms); and once per participant:
\texttt{list\_num}.
The continuous \texttt{pc\_prop\_rating} covariate (topic-level norming mean)
is joined to the trial data during analysis.

\subsection{Procedure}

Participants completed the experiment online via the
magpie platform \citep{kochetkova2024magpie} in a browser.
After reading instructions in German, participants were assigned to one
of 4 counterbalancing lists (via URL parameter or randomly).
They then proceeded through 16 trials---8 critical and 8 fillers---in a
fully randomised order, followed by a post-test questionnaire eliciting
native language, age, gender, education level, and optional comments.
Each trial displayed (a) a context paragraph, (b) a goal instruction,
and (c) the sentence frame with the marker drop-down embedded at the
Mittelfeld position.
The ``Weiter'' (Next) button was disabled until a marker had been
selected, ensuring a response on every trial.

\end{document}
.
%
% Compile stand-alone:
%   pdflatex production_exp.tex
%
% Embed in main.tex:
%   % production_exp.tex
% Stand-alone description of the production experiment.
% Can be embedded in main.tex via \input{production_exp}.
%
% Compile stand-alone:
%   pdflatex production_exp.tex
%
% Embed in main.tex:
%   \input{production_exp}   % (inside a \section{Experiment} environment)

\documentclass[12pt]{article}
\usepackage{CSPArticleStyle}
\usepackage{csquotes}
\usepackage[margin=1in]{geometry}
\usepackage{booktabs}
\usepackage{parskip}
\usepackage{amsmath}
\usepackage{times}
\usepackage[T1]{fontenc}
\usepackage[utf8]{inputenc}
\usepackage{xcolor}
\usepackage{tabularray}   % robust tables; fall back to tabular if unavailable
\usepackage{gb4e}         % linguistic examples
\noautomath

\title{Production Experiment: Eliciting Consensus-Marker Choice\\
       Across Perceived-Controversy and Goal-Strength Conditions}
\author{}
\date{}

\begin{document}

%% ============================================================
%% When embedding in main.tex, delete from \documentclass to
%% \begin{document} and delete \end{document} at the bottom.
%% ============================================================

\section{Experiment: Consensus-Marker Production}

\subsection{Overview}

The production experiment empirically estimates the decision thresholds
that govern the use of five German consensus markers by placing
participants in the speaker role.
Participants read short vignettes that independently varied
\emph{pragmatic perceived controversy} (\textsc{pc\textsubscript{prag}})
and the speaker's \emph{persuasive-goal strength} ($g$), then selected
which consensus marker they would most naturally use in the described
situation.
\emph{Propositional perceived controversy} (\textsc{pc\textsubscript{prop}})
is held constant within the experiment and operationalised as a continuous
topic-level covariate obtained from a separate norming study: for each
of the eight topic items, participants rated the degree to which the
proposition $p$ is considered established/accepted (0--100 scale);
the resulting means enter the regression analyses as a fixed covariate.
The resulting distribution of marker choices across the
$\textsc{pc}_\text{prag} \times g$ design, combined with the
continuous \textsc{pc\textsubscript{prop}} covariate, provides the
empirical basis for estimating the ordinal thresholds
$\mu_1 < \cdots < \mu_4$ of the noisy-threshold RSA model.

\subsection{Participants}

Participants will be recruited via Prolific and will be required to be
native speakers of German (language-at-home filter).
Data collection is ongoing; a target of $N = 80$ participants is planned
($n = 20$ per list across 4 lists), each completing 16 trials (8 critical + 8 filler).

\subsection{Materials}

\subsubsection{Markers}

The five consensus markers under investigation are, from weakest to
strongest presumed common-ground commitment:

\begin{enumerate}
  \item \textit{sofern ich wei\ss{}} (`as far as I know')
  \item \textit{wie du ja wei\ss t} (`as you (prt) know')
  \item \textit{wie wir wissen} (`as we know')
  \item \textit{ja} (modal particle; roughly `as you know / of course')
  \item \textit{bekanntlich} (`as is well known / notoriously')
\end{enumerate}

These markers were embedded in the sentential \emph{Mittelfeld}
(middle field), immediately after the finite verb, as illustrated in
example~(\ref{ex:frame}):

\begin{exe}
  \ex \label{ex:frame}
  Digitale Kompetenzen sind \textbf{[MARKER]}
  in der heutigen Arbeitswelt unverzichtbar.\\
  `Digital skills are \textbf{[MARKER]} indispensable in today's
  working world.'
\end{exe}

This syntactic position is licit for all five markers: modal particles
such as \textit{ja} canonically occupy the left edge of the Mittelfeld
\citep{thurmair1989modalpartikeln}, while the adverb \textit{bekanntlich}
and parenthetical clauses (\textit{wie wir wissen}, etc.) can appear
there as well (see \citealt{jacobs_semantics_1991} for discussion).
Each target sentence was followed by a consequent clause with the
recommended action $q$: \textit{Deshalb solltest du [q].}
(`Therefore, you should [q].')

\subsubsection{Topic Items}

Eight topic domains were used, each associated with a fixed proposition
$p$ (sentence frame) and consequent action $q$ (Table~\ref{tab:items}).
The first four items were used in earlier pilot work; items 5--8 were
added to reach the minimum of eight required for a complete Latin square.

\begin{table}[ht]
\centering
\caption{Topic items: sentence frame for $p$ (with marker slot [M]) and consequent $q$.}
\label{tab:items}
\small
\begin{tabular}{clll}
\toprule
\textbf{\#} & \textbf{Topic} & \textbf{Sentence frame} & \textbf{Consequent $q$} \\
\midrule
1 & Klimawandel &
  \textit{Der Klimawandel wird [M]} &
  auf Flugreisen verzichten \\
  & & \textit{durch menschliche Aktivit\"aten verursacht.} & \\[3pt]
2 & Ern\"ahrung &
  \textit{Eine \"uberwiegend pflanzliche Ern\"ahrung} &
  mehr pflanzliche Lebens- \\
  & & \textit{f\"ordert [M] die Gesundheit.} &
  mittel einbauen \\[3pt]
3 & Stadtverkehr &
  \textit{Ein ausgebautes Nahverkehrsnetz} &
  h\"aufiger auf \"offentliche \\
  & & \textit{entlastet [M] den Stadtverkehr deutlich.} &
  Verkehrsmittel umsteigen \\[3pt]
4 & Digitalisierung &
  \textit{Digitale Kompetenzen sind [M] in der} &
  regelm\"a\ss{}ig Weiterbildungs- \\
  & & \textit{heutigen Arbeitswelt unverzichtbar.} &
  angebote nutzen \\[3pt]
5 & Schlaf &
  \textit{Ausreichend Schlaf ist [M] entscheidend} &
  auf feste Schlafzeiten \\
  & & \textit{f\"ur die kognitive Leistungsf\"ahigkeit.} &
  achten \\[3pt]
6 & Plastik &
  \textit{Einwegplastik schadet [M] den} &
  konsequent auf Einweg- \\
  & & \textit{Meeres\"okosystemen erheblich.} &
  plastik verzichten \\[3pt]
7 & Sport &
  \textit{Regelm\"a\ss{}ige k\"orperliche Bewegung} &
  regelm\"a\ss{}ig Sport in \\
  & & \textit{senkt [M] das Risiko f\"ur} &
  den Alltag integrieren \\
  & & \textit{Herz-Kreislauf-Erkrankungen deutlich.} & \\[3pt]
8 & Lokalkauf &
  \textit{Das Kaufen bei lokalen H\"andlern} &
  \"ofter bei lokalen \\
  & & \textit{st\"arkt [M] die regionale Wirtschaft} &
  Gesch\"aften einkaufen \\
  & & \textit{nachhaltig.} & \\
\bottomrule
\end{tabular}
\end{table}

\noindent [M] denotes the marker slot (dropdown).

\subsubsection{Vignettes}

\paragraph{Propositional perceived controversy (\textsc{pc\textsubscript{prop}}).}
\textsc{pc\textsubscript{prop}} was \emph{not} manipulated within the
production experiment.
Instead, it was operationalised as a continuous, topic-level covariate
obtained from a separate norming study conducted prior to the production
experiment.
In that norming study, a different sample of native German speakers rated,
for each topic proposition $p$, the degree to which they considered $p$
accepted/established in (German) public discourse, on a 0--100 scale
($0 = $ `completely contested', $100 = $ `universally accepted').
The resulting per-topic means (e.g., \textit{Digitalisierung}: 85/100;
\textit{Lokalkauf}: 58/100) enter the regression analyses as a
continuous, centred predictor \textsc{pc\textsubscript{prop}}$_c$
(mean-centred, SD $\approx$ 1).
Context paragraphs in the production experiment therefore present
the proposition $p$ without any editorial evaluation, so that
participants' experience of propositional controversy is shaped
solely by their own prior knowledge and the norming-derived covariate
captures this variation in analysis.

\paragraph{Pragmatic perceived controversy (\textsc{pc\textsubscript{prag}}).}
\textsc{pc\textsubscript{prag}} was manipulated between instances of each item
by varying whether the speaker's social circle \emph{largely agreed} (low)
or \emph{held divergent views} (high) about $p$.
The two context versions for the \textit{Digitalisierung} item are:

\begin{itemize}
  \item $\textsc{pc}_\text{prag} = \text{low}$: \textit{Du sprichst mit einer
        Kollegin {\"u}ber Berufsausbildung und Weiterbildung.
        In eurem Team teilen die meisten die Ansicht,
        dass digitale Kompetenzen in der heutigen Arbeitswelt unverzichtbar sind.}
  \item $\textsc{pc}_\text{prag} = \text{high}$: \textit{Du sprichst mit einer
        Kollegin {\"u}ber Berufsausbildung und Weiterbildung.
        In eurem Team gehen die Meinungen dazu stark auseinander.}
\end{itemize}

\paragraph{Speaker goal strength ($g$).}
Goal strength was manipulated via a standardised goal instruction
appended after the context paragraph:

\begin{itemize}
  \item $g_\text{low}$: \textit{Du m\"ochtest das nur kurz anmerken.}
        (`You just want to make a brief note of this.')
  \item $g_\text{high}$: \textit{Es ist dir sehr wichtig, dass dein
        Gespr\"achspartner das akzeptiert.}
        (`It is very important to you that your interlocutor accepts this.')
\end{itemize}

\paragraph{Complete sample item (all four conditions).}
Examples (\ref{ex:cond1})--(\ref{ex:cond4}) show the full vignette for
the topic \textit{Digitalisierung}
(norming $\textsc{pc}_\text{prop}$ rating: 85/100;
sentence frame: \textit{Digitale Kompetenzen sind [M] in der heutigen
Arbeitswelt unverzichtbar.
Deshalb solltest du regelm\"a\ss{}ig digitale Weiterbildungsangebote nutzen.})
across all four conditions of the $2\times2$ design.

\begin{exe}
  \ex \label{ex:cond1}
  \textbf{Condition 0: $\textsc{pc}_\text{prag} = \text{low}$, $g = \text{low}$}\\
  \textit{Du sprichst mit einer Kollegin {\"u}ber Berufsausbildung und Weiterbildung.
  In eurem Team teilen die meisten die Ansicht,
  dass digitale Kompetenzen in der heutigen Arbeitswelt unverzichtbar sind.}\\
  \textit{Du m\"ochtest das nur kurz anmerken.}\\
  Digitale Kompetenzen sind \underline{\hspace{2.2cm}}
  in der heutigen Arbeitswelt unverzichtbar.

  \ex \label{ex:cond2}
  \textbf{Condition 1: $\textsc{pc}_\text{prag} = \text{low}$, $g = \text{high}$}\\
  \textit{Du sprichst mit einer Kollegin {\"u}ber Berufsausbildung und Weiterbildung.
  In eurem Team teilen die meisten die Ansicht,
  dass digitale Kompetenzen in der heutigen Arbeitswelt unverzichtbar sind.}\\
  \textit{Es ist dir sehr wichtig, dass dein Gespr\"achspartner das akzeptiert.}\\
  Digitale Kompetenzen sind \underline{\hspace{2.2cm}}
  in der heutigen Arbeitswelt unverzichtbar.

  \ex \label{ex:cond3}
  \textbf{Condition 2: $\textsc{pc}_\text{prag} = \text{high}$, $g = \text{low}$}\\
  \textit{Du sprichst mit einer Kollegin {\"u}ber Berufsausbildung und Weiterbildung.
  In eurem Team gehen die Meinungen dazu stark auseinander.}\\
  \textit{Du m\"ochtest das nur kurz anmerken.}\\
  Digitale Kompetenzen sind \underline{\hspace{2.2cm}}
  in der heutigen Arbeitswelt unverzichtbar.

  \ex \label{ex:cond4}
  \textbf{Condition 3: $\textsc{pc}_\text{prag} = \text{high}$, $g = \text{high}$}\\
  \textit{Du sprichst mit einer Kollegin {\"u}ber Berufsausbildung und Weiterbildung.
  In eurem Team gehen die Meinungen dazu stark auseinander.}\\
  \textit{Es ist dir sehr wichtig, dass dein Gespr\"achspartner das akzeptiert.}\\
  Digitale Kompetenzen sind \underline{\hspace{2.2cm}}
  in der heutigen Arbeitswelt unverzichtbar.
\end{exe}

\noindent In all four vignettes the sentence frame and consequent clause
(\textit{Deshalb solltest du regelm\"a\ss{}ig digitale
Weiterbildungsangebote nutzen.}) are identical; only the context
paragraph and goal instruction vary across conditions.
The marker slot [M] / \underline{\hspace{2.2cm}} was presented as a
drop-down menu in the actual experiment.

\subsection{Design}

\subsubsection{Factorial structure}

The experiment employed a \textbf{2 (pc\textsubscript{prag}: low vs.\ high)
$\times$ 2 ($g$: low vs.\ high) fully crossed design},
yielding 4 conditions.
\textsc{pc\textsubscript{prop}} is not a within-experiment factor; it enters
as a continuous topic-level covariate from the norming study.
With 8 topic items and a Latin square rotation across 4 lists, each
participant saw exactly \textbf{8 critical trials}---two items per condition---plus
8 filler trials, for \textbf{16 trials} in total
(Table~\ref{tab:design}).

\begin{table}[ht]
\centering
\caption{%
  The $2\times2$ design (4 conditions).
  \textsc{pc\textsubscript{prop}} is a continuous norming covariate, not a
  within-experiment factor.
  In each list, the 8 items are assigned to conditions via the Latin square
  rotation described in Section~\ref{subsub:counterbalancing}.}
\label{tab:design}
\small
\begin{tabular}{ccl}
\toprule
\textbf{pc\textsubscript{prag}} &
\textbf{$g$} &
\textbf{Condition label} \\
\midrule
low  & low  & (0) community agrees, casual \\
low  & high & (1) community agrees, pressing \\
high & low  & (2) community divided, casual \\
high & high & (3) community divided, pressing \\
\bottomrule
\end{tabular}
\end{table}

\subsubsection{Counterbalancing}
\label{subsub:counterbalancing}

The experiment used a \textbf{Latin square} to counterbalance item--condition
pairings across participants.
With $k = 4$ conditions and 8 items, a balanced assignment is achieved:
each of the 4 lists assigns item $i$ ($i = 0,\ldots,7$) to condition
$(i + l)\bmod 4$, where $l \in \{0,1,2,3\}$ is the list number.
The resulting assignment matrix (Table~\ref{tab:ls}) guarantees that
\begin{itemize}
  \item each participant sees every item exactly once;
  \item each participant sees each condition exactly twice (two items per
        condition), ensuring balanced coverage of the 2$\times$2 design;
  \item across all 4 lists, each item appears equally often in each
        condition (twice per condition over the full participant sample).
\end{itemize}

\begin{table}[ht]
\centering
\caption{Latin square assignment of conditions to items across 4 lists.
  Cell entries are condition indices (0--3) as defined in
  Table~\ref{tab:design}.
  Rows = lists ($l$); columns = items ($i$).}
\label{tab:ls}
\small
\begin{tabular}{ccccccccc}
\toprule
\textbf{List} &
\textbf{Item 1} & \textbf{Item 2} & \textbf{Item 3} & \textbf{Item 4} &
\textbf{Item 5} & \textbf{Item 6} & \textbf{Item 7} & \textbf{Item 8} \\
\midrule
0 & 0 & 1 & 2 & 3 & 0 & 1 & 2 & 3 \\
1 & 1 & 2 & 3 & 0 & 1 & 2 & 3 & 0 \\
2 & 2 & 3 & 0 & 1 & 2 & 3 & 0 & 1 \\
3 & 3 & 0 & 1 & 2 & 3 & 0 & 1 & 2 \\
\bottomrule
\end{tabular}
\end{table}

\noindent List assignment is determined automatically from the URL query
parameter \texttt{?list=N} (0--3); if absent, a random list is selected.
Each participant is thus within-subjects across items and sees all four
conditions (twice each). Condition effects are estimated from within-participant
variation across the 2$\times$2 design, controlling for by-participant
random intercepts.

In addition to the 8 critical trials, each participant completed 8
\textbf{filler trials} drawn from separate topic domains (train travel,
reading light, cooking, coffee, vacation, hydration, fresh air, noise).
Filler items used the same trial format but were excluded from the
main analysis; they serve to dilute the critical items and reduce
potential demand characteristics.

\subsubsection{Dependent variable}

The dependent variable was the marker chosen from a drop-down menu
displaying all five options in fixed order (as listed in
Section~\ref{sub:markers}).
No response time threshold was imposed, but response times were logged.
Recorded per trial: \texttt{trial\_id}, \texttt{topic}, \texttt{is\_filler},
\texttt{condition\_index}, \texttt{pc\_prag}, \texttt{g},
\texttt{selected\_marker}, \texttt{rt} (ms); and once per participant:
\texttt{list\_num}.
The continuous \texttt{pc\_prop\_rating} covariate (topic-level norming mean)
is joined to the trial data during analysis.

\subsection{Procedure}

Participants completed the experiment online via the
magpie platform \citep{kochetkova2024magpie} in a browser.
After reading instructions in German, participants were assigned to one
of 4 counterbalancing lists (via URL parameter or randomly).
They then proceeded through 16 trials---8 critical and 8 fillers---in a
fully randomised order, followed by a post-test questionnaire eliciting
native language, age, gender, education level, and optional comments.
Each trial displayed (a) a context paragraph, (b) a goal instruction,
and (c) the sentence frame with the marker drop-down embedded at the
Mittelfeld position.
The ``Weiter'' (Next) button was disabled until a marker had been
selected, ensuring a response on every trial.

\end{document}
   % (inside a \section{Experiment} environment)

\documentclass[12pt]{article}
\usepackage{CSPArticleStyle}
\usepackage{csquotes}
\usepackage[margin=1in]{geometry}
\usepackage{booktabs}
\usepackage{parskip}
\usepackage{amsmath}
\usepackage{times}
\usepackage[T1]{fontenc}
\usepackage[utf8]{inputenc}
\usepackage{xcolor}
\usepackage{tabularray}   % robust tables; fall back to tabular if unavailable
\usepackage{gb4e}         % linguistic examples
\noautomath

\title{Production Experiment: Eliciting Consensus-Marker Choice\\
       Across Perceived-Controversy and Goal-Strength Conditions}
\author{}
\date{}

\begin{document}

%% ============================================================
%% When embedding in main.tex, delete from \documentclass to
%% \begin{document} and delete \end{document} at the bottom.
%% ============================================================

\section{Experiment: Consensus-Marker Production}

\subsection{Overview}

The production experiment empirically estimates the decision thresholds
that govern the use of five German consensus markers by placing
participants in the speaker role.
Participants read short vignettes that independently varied
\emph{pragmatic perceived controversy} (\textsc{pc\textsubscript{prag}})
and the speaker's \emph{persuasive-goal strength} ($g$), then selected
which consensus marker they would most naturally use in the described
situation.
\emph{Propositional perceived controversy} (\textsc{pc\textsubscript{prop}})
is held constant within the experiment and operationalised as a continuous
topic-level covariate obtained from a separate norming study: for each
of the eight topic items, participants rated the degree to which the
proposition $p$ is considered established/accepted (0--100 scale);
the resulting means enter the regression analyses as a fixed covariate.
The resulting distribution of marker choices across the
$\textsc{pc}_\text{prag} \times g$ design, combined with the
continuous \textsc{pc\textsubscript{prop}} covariate, provides the
empirical basis for estimating the ordinal thresholds
$\mu_1 < \cdots < \mu_4$ of the noisy-threshold RSA model.

\subsection{Participants}

Participants will be recruited via Prolific and will be required to be
native speakers of German (language-at-home filter).
Data collection is ongoing; a target of $N = 80$ participants is planned
($n = 20$ per list across 4 lists), each completing 16 trials (8 critical + 8 filler).

\subsection{Materials}

\subsubsection{Markers}

The five consensus markers under investigation are, from weakest to
strongest presumed common-ground commitment:

\begin{enumerate}
  \item \textit{sofern ich wei\ss{}} (`as far as I know')
  \item \textit{wie du ja wei\ss t} (`as you (prt) know')
  \item \textit{wie wir wissen} (`as we know')
  \item \textit{ja} (modal particle; roughly `as you know / of course')
  \item \textit{bekanntlich} (`as is well known / notoriously')
\end{enumerate}

These markers were embedded in the sentential \emph{Mittelfeld}
(middle field), immediately after the finite verb, as illustrated in
example~(\ref{ex:frame}):

\begin{exe}
  \ex \label{ex:frame}
  Digitale Kompetenzen sind \textbf{[MARKER]}
  in der heutigen Arbeitswelt unverzichtbar.\\
  `Digital skills are \textbf{[MARKER]} indispensable in today's
  working world.'
\end{exe}

This syntactic position is licit for all five markers: modal particles
such as \textit{ja} canonically occupy the left edge of the Mittelfeld
\citep{thurmair1989modalpartikeln}, while the adverb \textit{bekanntlich}
and parenthetical clauses (\textit{wie wir wissen}, etc.) can appear
there as well (see \citealt{jacobs_semantics_1991} for discussion).
Each target sentence was followed by a consequent clause with the
recommended action $q$: \textit{Deshalb solltest du [q].}
(`Therefore, you should [q].')

\subsubsection{Topic Items}

Eight topic domains were used, each associated with a fixed proposition
$p$ (sentence frame) and consequent action $q$ (Table~\ref{tab:items}).
The first four items were used in earlier pilot work; items 5--8 were
added to reach the minimum of eight required for a complete Latin square.

\begin{table}[ht]
\centering
\caption{Topic items: sentence frame for $p$ (with marker slot [M]) and consequent $q$.}
\label{tab:items}
\small
\begin{tabular}{clll}
\toprule
\textbf{\#} & \textbf{Topic} & \textbf{Sentence frame} & \textbf{Consequent $q$} \\
\midrule
1 & Klimawandel &
  \textit{Der Klimawandel wird [M]} &
  auf Flugreisen verzichten \\
  & & \textit{durch menschliche Aktivit\"aten verursacht.} & \\[3pt]
2 & Ern\"ahrung &
  \textit{Eine \"uberwiegend pflanzliche Ern\"ahrung} &
  mehr pflanzliche Lebens- \\
  & & \textit{f\"ordert [M] die Gesundheit.} &
  mittel einbauen \\[3pt]
3 & Stadtverkehr &
  \textit{Ein ausgebautes Nahverkehrsnetz} &
  h\"aufiger auf \"offentliche \\
  & & \textit{entlastet [M] den Stadtverkehr deutlich.} &
  Verkehrsmittel umsteigen \\[3pt]
4 & Digitalisierung &
  \textit{Digitale Kompetenzen sind [M] in der} &
  regelm\"a\ss{}ig Weiterbildungs- \\
  & & \textit{heutigen Arbeitswelt unverzichtbar.} &
  angebote nutzen \\[3pt]
5 & Schlaf &
  \textit{Ausreichend Schlaf ist [M] entscheidend} &
  auf feste Schlafzeiten \\
  & & \textit{f\"ur die kognitive Leistungsf\"ahigkeit.} &
  achten \\[3pt]
6 & Plastik &
  \textit{Einwegplastik schadet [M] den} &
  konsequent auf Einweg- \\
  & & \textit{Meeres\"okosystemen erheblich.} &
  plastik verzichten \\[3pt]
7 & Sport &
  \textit{Regelm\"a\ss{}ige k\"orperliche Bewegung} &
  regelm\"a\ss{}ig Sport in \\
  & & \textit{senkt [M] das Risiko f\"ur} &
  den Alltag integrieren \\
  & & \textit{Herz-Kreislauf-Erkrankungen deutlich.} & \\[3pt]
8 & Lokalkauf &
  \textit{Das Kaufen bei lokalen H\"andlern} &
  \"ofter bei lokalen \\
  & & \textit{st\"arkt [M] die regionale Wirtschaft} &
  Gesch\"aften einkaufen \\
  & & \textit{nachhaltig.} & \\
\bottomrule
\end{tabular}
\end{table}

\noindent [M] denotes the marker slot (dropdown).

\subsubsection{Vignettes}

\paragraph{Propositional perceived controversy (\textsc{pc\textsubscript{prop}}).}
\textsc{pc\textsubscript{prop}} was \emph{not} manipulated within the
production experiment.
Instead, it was operationalised as a continuous, topic-level covariate
obtained from a separate norming study conducted prior to the production
experiment.
In that norming study, a different sample of native German speakers rated,
for each topic proposition $p$, the degree to which they considered $p$
accepted/established in (German) public discourse, on a 0--100 scale
($0 = $ `completely contested', $100 = $ `universally accepted').
The resulting per-topic means (e.g., \textit{Digitalisierung}: 85/100;
\textit{Lokalkauf}: 58/100) enter the regression analyses as a
continuous, centred predictor \textsc{pc\textsubscript{prop}}$_c$
(mean-centred, SD $\approx$ 1).
Context paragraphs in the production experiment therefore present
the proposition $p$ without any editorial evaluation, so that
participants' experience of propositional controversy is shaped
solely by their own prior knowledge and the norming-derived covariate
captures this variation in analysis.

\paragraph{Pragmatic perceived controversy (\textsc{pc\textsubscript{prag}}).}
\textsc{pc\textsubscript{prag}} was manipulated between instances of each item
by varying whether the speaker's social circle \emph{largely agreed} (low)
or \emph{held divergent views} (high) about $p$.
The two context versions for the \textit{Digitalisierung} item are:

\begin{itemize}
  \item $\textsc{pc}_\text{prag} = \text{low}$: \textit{Du sprichst mit einer
        Kollegin {\"u}ber Berufsausbildung und Weiterbildung.
        In eurem Team teilen die meisten die Ansicht,
        dass digitale Kompetenzen in der heutigen Arbeitswelt unverzichtbar sind.}
  \item $\textsc{pc}_\text{prag} = \text{high}$: \textit{Du sprichst mit einer
        Kollegin {\"u}ber Berufsausbildung und Weiterbildung.
        In eurem Team gehen die Meinungen dazu stark auseinander.}
\end{itemize}

\paragraph{Speaker goal strength ($g$).}
Goal strength was manipulated via a standardised goal instruction
appended after the context paragraph:

\begin{itemize}
  \item $g_\text{low}$: \textit{Du m\"ochtest das nur kurz anmerken.}
        (`You just want to make a brief note of this.')
  \item $g_\text{high}$: \textit{Es ist dir sehr wichtig, dass dein
        Gespr\"achspartner das akzeptiert.}
        (`It is very important to you that your interlocutor accepts this.')
\end{itemize}

\paragraph{Complete sample item (all four conditions).}
Examples (\ref{ex:cond1})--(\ref{ex:cond4}) show the full vignette for
the topic \textit{Digitalisierung}
(norming $\textsc{pc}_\text{prop}$ rating: 85/100;
sentence frame: \textit{Digitale Kompetenzen sind [M] in der heutigen
Arbeitswelt unverzichtbar.
Deshalb solltest du regelm\"a\ss{}ig digitale Weiterbildungsangebote nutzen.})
across all four conditions of the $2\times2$ design.

\begin{exe}
  \ex \label{ex:cond1}
  \textbf{Condition 0: $\textsc{pc}_\text{prag} = \text{low}$, $g = \text{low}$}\\
  \textit{Du sprichst mit einer Kollegin {\"u}ber Berufsausbildung und Weiterbildung.
  In eurem Team teilen die meisten die Ansicht,
  dass digitale Kompetenzen in der heutigen Arbeitswelt unverzichtbar sind.}\\
  \textit{Du m\"ochtest das nur kurz anmerken.}\\
  Digitale Kompetenzen sind \underline{\hspace{2.2cm}}
  in der heutigen Arbeitswelt unverzichtbar.

  \ex \label{ex:cond2}
  \textbf{Condition 1: $\textsc{pc}_\text{prag} = \text{low}$, $g = \text{high}$}\\
  \textit{Du sprichst mit einer Kollegin {\"u}ber Berufsausbildung und Weiterbildung.
  In eurem Team teilen die meisten die Ansicht,
  dass digitale Kompetenzen in der heutigen Arbeitswelt unverzichtbar sind.}\\
  \textit{Es ist dir sehr wichtig, dass dein Gespr\"achspartner das akzeptiert.}\\
  Digitale Kompetenzen sind \underline{\hspace{2.2cm}}
  in der heutigen Arbeitswelt unverzichtbar.

  \ex \label{ex:cond3}
  \textbf{Condition 2: $\textsc{pc}_\text{prag} = \text{high}$, $g = \text{low}$}\\
  \textit{Du sprichst mit einer Kollegin {\"u}ber Berufsausbildung und Weiterbildung.
  In eurem Team gehen die Meinungen dazu stark auseinander.}\\
  \textit{Du m\"ochtest das nur kurz anmerken.}\\
  Digitale Kompetenzen sind \underline{\hspace{2.2cm}}
  in der heutigen Arbeitswelt unverzichtbar.

  \ex \label{ex:cond4}
  \textbf{Condition 3: $\textsc{pc}_\text{prag} = \text{high}$, $g = \text{high}$}\\
  \textit{Du sprichst mit einer Kollegin {\"u}ber Berufsausbildung und Weiterbildung.
  In eurem Team gehen die Meinungen dazu stark auseinander.}\\
  \textit{Es ist dir sehr wichtig, dass dein Gespr\"achspartner das akzeptiert.}\\
  Digitale Kompetenzen sind \underline{\hspace{2.2cm}}
  in der heutigen Arbeitswelt unverzichtbar.
\end{exe}

\noindent In all four vignettes the sentence frame and consequent clause
(\textit{Deshalb solltest du regelm\"a\ss{}ig digitale
Weiterbildungsangebote nutzen.}) are identical; only the context
paragraph and goal instruction vary across conditions.
The marker slot [M] / \underline{\hspace{2.2cm}} was presented as a
drop-down menu in the actual experiment.

\subsection{Design}

\subsubsection{Factorial structure}

The experiment employed a \textbf{2 (pc\textsubscript{prag}: low vs.\ high)
$\times$ 2 ($g$: low vs.\ high) fully crossed design},
yielding 4 conditions.
\textsc{pc\textsubscript{prop}} is not a within-experiment factor; it enters
as a continuous topic-level covariate from the norming study.
With 8 topic items and a Latin square rotation across 4 lists, each
participant saw exactly \textbf{8 critical trials}---two items per condition---plus
8 filler trials, for \textbf{16 trials} in total
(Table~\ref{tab:design}).

\begin{table}[ht]
\centering
\caption{%
  The $2\times2$ design (4 conditions).
  \textsc{pc\textsubscript{prop}} is a continuous norming covariate, not a
  within-experiment factor.
  In each list, the 8 items are assigned to conditions via the Latin square
  rotation described in Section~\ref{subsub:counterbalancing}.}
\label{tab:design}
\small
\begin{tabular}{ccl}
\toprule
\textbf{pc\textsubscript{prag}} &
\textbf{$g$} &
\textbf{Condition label} \\
\midrule
low  & low  & (0) community agrees, casual \\
low  & high & (1) community agrees, pressing \\
high & low  & (2) community divided, casual \\
high & high & (3) community divided, pressing \\
\bottomrule
\end{tabular}
\end{table}

\subsubsection{Counterbalancing}
\label{subsub:counterbalancing}

The experiment used a \textbf{Latin square} to counterbalance item--condition
pairings across participants.
With $k = 4$ conditions and 8 items, a balanced assignment is achieved:
each of the 4 lists assigns item $i$ ($i = 0,\ldots,7$) to condition
$(i + l)\bmod 4$, where $l \in \{0,1,2,3\}$ is the list number.
The resulting assignment matrix (Table~\ref{tab:ls}) guarantees that
\begin{itemize}
  \item each participant sees every item exactly once;
  \item each participant sees each condition exactly twice (two items per
        condition), ensuring balanced coverage of the 2$\times$2 design;
  \item across all 4 lists, each item appears equally often in each
        condition (twice per condition over the full participant sample).
\end{itemize}

\begin{table}[ht]
\centering
\caption{Latin square assignment of conditions to items across 4 lists.
  Cell entries are condition indices (0--3) as defined in
  Table~\ref{tab:design}.
  Rows = lists ($l$); columns = items ($i$).}
\label{tab:ls}
\small
\begin{tabular}{ccccccccc}
\toprule
\textbf{List} &
\textbf{Item 1} & \textbf{Item 2} & \textbf{Item 3} & \textbf{Item 4} &
\textbf{Item 5} & \textbf{Item 6} & \textbf{Item 7} & \textbf{Item 8} \\
\midrule
0 & 0 & 1 & 2 & 3 & 0 & 1 & 2 & 3 \\
1 & 1 & 2 & 3 & 0 & 1 & 2 & 3 & 0 \\
2 & 2 & 3 & 0 & 1 & 2 & 3 & 0 & 1 \\
3 & 3 & 0 & 1 & 2 & 3 & 0 & 1 & 2 \\
\bottomrule
\end{tabular}
\end{table}

\noindent List assignment is determined automatically from the URL query
parameter \texttt{?list=N} (0--3); if absent, a random list is selected.
Each participant is thus within-subjects across items and sees all four
conditions (twice each). Condition effects are estimated from within-participant
variation across the 2$\times$2 design, controlling for by-participant
random intercepts.

In addition to the 8 critical trials, each participant completed 8
\textbf{filler trials} drawn from separate topic domains (train travel,
reading light, cooking, coffee, vacation, hydration, fresh air, noise).
Filler items used the same trial format but were excluded from the
main analysis; they serve to dilute the critical items and reduce
potential demand characteristics.

\subsubsection{Dependent variable}

The dependent variable was the marker chosen from a drop-down menu
displaying all five options in fixed order (as listed in
Section~\ref{sub:markers}).
No response time threshold was imposed, but response times were logged.
Recorded per trial: \texttt{trial\_id}, \texttt{topic}, \texttt{is\_filler},
\texttt{condition\_index}, \texttt{pc\_prag}, \texttt{g},
\texttt{selected\_marker}, \texttt{rt} (ms); and once per participant:
\texttt{list\_num}.
The continuous \texttt{pc\_prop\_rating} covariate (topic-level norming mean)
is joined to the trial data during analysis.

\subsection{Procedure}

Participants completed the experiment online via the
magpie platform \citep{kochetkova2024magpie} in a browser.
After reading instructions in German, participants were assigned to one
of 4 counterbalancing lists (via URL parameter or randomly).
They then proceeded through 16 trials---8 critical and 8 fillers---in a
fully randomised order, followed by a post-test questionnaire eliciting
native language, age, gender, education level, and optional comments.
Each trial displayed (a) a context paragraph, (b) a goal instruction,
and (c) the sentence frame with the marker drop-down embedded at the
Mittelfeld position.
The ``Weiter'' (Next) button was disabled until a marker had been
selected, ensuring a response on every trial.

\end{document}
.
%
% Compile stand-alone:
%   pdflatex production_exp.tex
%
% Embed in main.tex:
%   % production_exp.tex
% Stand-alone description of the production experiment.
% Can be embedded in main.tex via % production_exp.tex
% Stand-alone description of the production experiment.
% Can be embedded in main.tex via \input{production_exp}.
%
% Compile stand-alone:
%   pdflatex production_exp.tex
%
% Embed in main.tex:
%   \input{production_exp}   % (inside a \section{Experiment} environment)

\documentclass[12pt]{article}
\usepackage{CSPArticleStyle}
\usepackage{csquotes}
\usepackage[margin=1in]{geometry}
\usepackage{booktabs}
\usepackage{parskip}
\usepackage{amsmath}
\usepackage{times}
\usepackage[T1]{fontenc}
\usepackage[utf8]{inputenc}
\usepackage{xcolor}
\usepackage{tabularray}   % robust tables; fall back to tabular if unavailable
\usepackage{gb4e}         % linguistic examples
\noautomath

\title{Production Experiment: Eliciting Consensus-Marker Choice\\
       Across Perceived-Controversy and Goal-Strength Conditions}
\author{}
\date{}

\begin{document}

%% ============================================================
%% When embedding in main.tex, delete from \documentclass to
%% \begin{document} and delete \end{document} at the bottom.
%% ============================================================

\section{Experiment: Consensus-Marker Production}

\subsection{Overview}

The production experiment empirically estimates the decision thresholds
that govern the use of five German consensus markers by placing
participants in the speaker role.
Participants read short vignettes that independently varied
\emph{pragmatic perceived controversy} (\textsc{pc\textsubscript{prag}})
and the speaker's \emph{persuasive-goal strength} ($g$), then selected
which consensus marker they would most naturally use in the described
situation.
\emph{Propositional perceived controversy} (\textsc{pc\textsubscript{prop}})
is held constant within the experiment and operationalised as a continuous
topic-level covariate obtained from a separate norming study: for each
of the eight topic items, participants rated the degree to which the
proposition $p$ is considered established/accepted (0--100 scale);
the resulting means enter the regression analyses as a fixed covariate.
The resulting distribution of marker choices across the
$\textsc{pc}_\text{prag} \times g$ design, combined with the
continuous \textsc{pc\textsubscript{prop}} covariate, provides the
empirical basis for estimating the ordinal thresholds
$\mu_1 < \cdots < \mu_4$ of the noisy-threshold RSA model.

\subsection{Participants}

Participants will be recruited via Prolific and will be required to be
native speakers of German (language-at-home filter).
Data collection is ongoing; a target of $N = 80$ participants is planned
($n = 20$ per list across 4 lists), each completing 16 trials (8 critical + 8 filler).

\subsection{Materials}

\subsubsection{Markers}

The five consensus markers under investigation are, from weakest to
strongest presumed common-ground commitment:

\begin{enumerate}
  \item \textit{sofern ich wei\ss{}} (`as far as I know')
  \item \textit{wie du ja wei\ss t} (`as you (prt) know')
  \item \textit{wie wir wissen} (`as we know')
  \item \textit{ja} (modal particle; roughly `as you know / of course')
  \item \textit{bekanntlich} (`as is well known / notoriously')
\end{enumerate}

These markers were embedded in the sentential \emph{Mittelfeld}
(middle field), immediately after the finite verb, as illustrated in
example~(\ref{ex:frame}):

\begin{exe}
  \ex \label{ex:frame}
  Digitale Kompetenzen sind \textbf{[MARKER]}
  in der heutigen Arbeitswelt unverzichtbar.\\
  `Digital skills are \textbf{[MARKER]} indispensable in today's
  working world.'
\end{exe}

This syntactic position is licit for all five markers: modal particles
such as \textit{ja} canonically occupy the left edge of the Mittelfeld
\citep{thurmair1989modalpartikeln}, while the adverb \textit{bekanntlich}
and parenthetical clauses (\textit{wie wir wissen}, etc.) can appear
there as well (see \citealt{jacobs_semantics_1991} for discussion).
Each target sentence was followed by a consequent clause with the
recommended action $q$: \textit{Deshalb solltest du [q].}
(`Therefore, you should [q].')

\subsubsection{Topic Items}

Eight topic domains were used, each associated with a fixed proposition
$p$ (sentence frame) and consequent action $q$ (Table~\ref{tab:items}).
The first four items were used in earlier pilot work; items 5--8 were
added to reach the minimum of eight required for a complete Latin square.

\begin{table}[ht]
\centering
\caption{Topic items: sentence frame for $p$ (with marker slot [M]) and consequent $q$.}
\label{tab:items}
\small
\begin{tabular}{clll}
\toprule
\textbf{\#} & \textbf{Topic} & \textbf{Sentence frame} & \textbf{Consequent $q$} \\
\midrule
1 & Klimawandel &
  \textit{Der Klimawandel wird [M]} &
  auf Flugreisen verzichten \\
  & & \textit{durch menschliche Aktivit\"aten verursacht.} & \\[3pt]
2 & Ern\"ahrung &
  \textit{Eine \"uberwiegend pflanzliche Ern\"ahrung} &
  mehr pflanzliche Lebens- \\
  & & \textit{f\"ordert [M] die Gesundheit.} &
  mittel einbauen \\[3pt]
3 & Stadtverkehr &
  \textit{Ein ausgebautes Nahverkehrsnetz} &
  h\"aufiger auf \"offentliche \\
  & & \textit{entlastet [M] den Stadtverkehr deutlich.} &
  Verkehrsmittel umsteigen \\[3pt]
4 & Digitalisierung &
  \textit{Digitale Kompetenzen sind [M] in der} &
  regelm\"a\ss{}ig Weiterbildungs- \\
  & & \textit{heutigen Arbeitswelt unverzichtbar.} &
  angebote nutzen \\[3pt]
5 & Schlaf &
  \textit{Ausreichend Schlaf ist [M] entscheidend} &
  auf feste Schlafzeiten \\
  & & \textit{f\"ur die kognitive Leistungsf\"ahigkeit.} &
  achten \\[3pt]
6 & Plastik &
  \textit{Einwegplastik schadet [M] den} &
  konsequent auf Einweg- \\
  & & \textit{Meeres\"okosystemen erheblich.} &
  plastik verzichten \\[3pt]
7 & Sport &
  \textit{Regelm\"a\ss{}ige k\"orperliche Bewegung} &
  regelm\"a\ss{}ig Sport in \\
  & & \textit{senkt [M] das Risiko f\"ur} &
  den Alltag integrieren \\
  & & \textit{Herz-Kreislauf-Erkrankungen deutlich.} & \\[3pt]
8 & Lokalkauf &
  \textit{Das Kaufen bei lokalen H\"andlern} &
  \"ofter bei lokalen \\
  & & \textit{st\"arkt [M] die regionale Wirtschaft} &
  Gesch\"aften einkaufen \\
  & & \textit{nachhaltig.} & \\
\bottomrule
\end{tabular}
\end{table}

\noindent [M] denotes the marker slot (dropdown).

\subsubsection{Vignettes}

\paragraph{Propositional perceived controversy (\textsc{pc\textsubscript{prop}}).}
\textsc{pc\textsubscript{prop}} was \emph{not} manipulated within the
production experiment.
Instead, it was operationalised as a continuous, topic-level covariate
obtained from a separate norming study conducted prior to the production
experiment.
In that norming study, a different sample of native German speakers rated,
for each topic proposition $p$, the degree to which they considered $p$
accepted/established in (German) public discourse, on a 0--100 scale
($0 = $ `completely contested', $100 = $ `universally accepted').
The resulting per-topic means (e.g., \textit{Digitalisierung}: 85/100;
\textit{Lokalkauf}: 58/100) enter the regression analyses as a
continuous, centred predictor \textsc{pc\textsubscript{prop}}$_c$
(mean-centred, SD $\approx$ 1).
Context paragraphs in the production experiment therefore present
the proposition $p$ without any editorial evaluation, so that
participants' experience of propositional controversy is shaped
solely by their own prior knowledge and the norming-derived covariate
captures this variation in analysis.

\paragraph{Pragmatic perceived controversy (\textsc{pc\textsubscript{prag}}).}
\textsc{pc\textsubscript{prag}} was manipulated between instances of each item
by varying whether the speaker's social circle \emph{largely agreed} (low)
or \emph{held divergent views} (high) about $p$.
The two context versions for the \textit{Digitalisierung} item are:

\begin{itemize}
  \item $\textsc{pc}_\text{prag} = \text{low}$: \textit{Du sprichst mit einer
        Kollegin {\"u}ber Berufsausbildung und Weiterbildung.
        In eurem Team teilen die meisten die Ansicht,
        dass digitale Kompetenzen in der heutigen Arbeitswelt unverzichtbar sind.}
  \item $\textsc{pc}_\text{prag} = \text{high}$: \textit{Du sprichst mit einer
        Kollegin {\"u}ber Berufsausbildung und Weiterbildung.
        In eurem Team gehen die Meinungen dazu stark auseinander.}
\end{itemize}

\paragraph{Speaker goal strength ($g$).}
Goal strength was manipulated via a standardised goal instruction
appended after the context paragraph:

\begin{itemize}
  \item $g_\text{low}$: \textit{Du m\"ochtest das nur kurz anmerken.}
        (`You just want to make a brief note of this.')
  \item $g_\text{high}$: \textit{Es ist dir sehr wichtig, dass dein
        Gespr\"achspartner das akzeptiert.}
        (`It is very important to you that your interlocutor accepts this.')
\end{itemize}

\paragraph{Complete sample item (all four conditions).}
Examples (\ref{ex:cond1})--(\ref{ex:cond4}) show the full vignette for
the topic \textit{Digitalisierung}
(norming $\textsc{pc}_\text{prop}$ rating: 85/100;
sentence frame: \textit{Digitale Kompetenzen sind [M] in der heutigen
Arbeitswelt unverzichtbar.
Deshalb solltest du regelm\"a\ss{}ig digitale Weiterbildungsangebote nutzen.})
across all four conditions of the $2\times2$ design.

\begin{exe}
  \ex \label{ex:cond1}
  \textbf{Condition 0: $\textsc{pc}_\text{prag} = \text{low}$, $g = \text{low}$}\\
  \textit{Du sprichst mit einer Kollegin {\"u}ber Berufsausbildung und Weiterbildung.
  In eurem Team teilen die meisten die Ansicht,
  dass digitale Kompetenzen in der heutigen Arbeitswelt unverzichtbar sind.}\\
  \textit{Du m\"ochtest das nur kurz anmerken.}\\
  Digitale Kompetenzen sind \underline{\hspace{2.2cm}}
  in der heutigen Arbeitswelt unverzichtbar.

  \ex \label{ex:cond2}
  \textbf{Condition 1: $\textsc{pc}_\text{prag} = \text{low}$, $g = \text{high}$}\\
  \textit{Du sprichst mit einer Kollegin {\"u}ber Berufsausbildung und Weiterbildung.
  In eurem Team teilen die meisten die Ansicht,
  dass digitale Kompetenzen in der heutigen Arbeitswelt unverzichtbar sind.}\\
  \textit{Es ist dir sehr wichtig, dass dein Gespr\"achspartner das akzeptiert.}\\
  Digitale Kompetenzen sind \underline{\hspace{2.2cm}}
  in der heutigen Arbeitswelt unverzichtbar.

  \ex \label{ex:cond3}
  \textbf{Condition 2: $\textsc{pc}_\text{prag} = \text{high}$, $g = \text{low}$}\\
  \textit{Du sprichst mit einer Kollegin {\"u}ber Berufsausbildung und Weiterbildung.
  In eurem Team gehen die Meinungen dazu stark auseinander.}\\
  \textit{Du m\"ochtest das nur kurz anmerken.}\\
  Digitale Kompetenzen sind \underline{\hspace{2.2cm}}
  in der heutigen Arbeitswelt unverzichtbar.

  \ex \label{ex:cond4}
  \textbf{Condition 3: $\textsc{pc}_\text{prag} = \text{high}$, $g = \text{high}$}\\
  \textit{Du sprichst mit einer Kollegin {\"u}ber Berufsausbildung und Weiterbildung.
  In eurem Team gehen die Meinungen dazu stark auseinander.}\\
  \textit{Es ist dir sehr wichtig, dass dein Gespr\"achspartner das akzeptiert.}\\
  Digitale Kompetenzen sind \underline{\hspace{2.2cm}}
  in der heutigen Arbeitswelt unverzichtbar.
\end{exe}

\noindent In all four vignettes the sentence frame and consequent clause
(\textit{Deshalb solltest du regelm\"a\ss{}ig digitale
Weiterbildungsangebote nutzen.}) are identical; only the context
paragraph and goal instruction vary across conditions.
The marker slot [M] / \underline{\hspace{2.2cm}} was presented as a
drop-down menu in the actual experiment.

\subsection{Design}

\subsubsection{Factorial structure}

The experiment employed a \textbf{2 (pc\textsubscript{prag}: low vs.\ high)
$\times$ 2 ($g$: low vs.\ high) fully crossed design},
yielding 4 conditions.
\textsc{pc\textsubscript{prop}} is not a within-experiment factor; it enters
as a continuous topic-level covariate from the norming study.
With 8 topic items and a Latin square rotation across 4 lists, each
participant saw exactly \textbf{8 critical trials}---two items per condition---plus
8 filler trials, for \textbf{16 trials} in total
(Table~\ref{tab:design}).

\begin{table}[ht]
\centering
\caption{%
  The $2\times2$ design (4 conditions).
  \textsc{pc\textsubscript{prop}} is a continuous norming covariate, not a
  within-experiment factor.
  In each list, the 8 items are assigned to conditions via the Latin square
  rotation described in Section~\ref{subsub:counterbalancing}.}
\label{tab:design}
\small
\begin{tabular}{ccl}
\toprule
\textbf{pc\textsubscript{prag}} &
\textbf{$g$} &
\textbf{Condition label} \\
\midrule
low  & low  & (0) community agrees, casual \\
low  & high & (1) community agrees, pressing \\
high & low  & (2) community divided, casual \\
high & high & (3) community divided, pressing \\
\bottomrule
\end{tabular}
\end{table}

\subsubsection{Counterbalancing}
\label{subsub:counterbalancing}

The experiment used a \textbf{Latin square} to counterbalance item--condition
pairings across participants.
With $k = 4$ conditions and 8 items, a balanced assignment is achieved:
each of the 4 lists assigns item $i$ ($i = 0,\ldots,7$) to condition
$(i + l)\bmod 4$, where $l \in \{0,1,2,3\}$ is the list number.
The resulting assignment matrix (Table~\ref{tab:ls}) guarantees that
\begin{itemize}
  \item each participant sees every item exactly once;
  \item each participant sees each condition exactly twice (two items per
        condition), ensuring balanced coverage of the 2$\times$2 design;
  \item across all 4 lists, each item appears equally often in each
        condition (twice per condition over the full participant sample).
\end{itemize}

\begin{table}[ht]
\centering
\caption{Latin square assignment of conditions to items across 4 lists.
  Cell entries are condition indices (0--3) as defined in
  Table~\ref{tab:design}.
  Rows = lists ($l$); columns = items ($i$).}
\label{tab:ls}
\small
\begin{tabular}{ccccccccc}
\toprule
\textbf{List} &
\textbf{Item 1} & \textbf{Item 2} & \textbf{Item 3} & \textbf{Item 4} &
\textbf{Item 5} & \textbf{Item 6} & \textbf{Item 7} & \textbf{Item 8} \\
\midrule
0 & 0 & 1 & 2 & 3 & 0 & 1 & 2 & 3 \\
1 & 1 & 2 & 3 & 0 & 1 & 2 & 3 & 0 \\
2 & 2 & 3 & 0 & 1 & 2 & 3 & 0 & 1 \\
3 & 3 & 0 & 1 & 2 & 3 & 0 & 1 & 2 \\
\bottomrule
\end{tabular}
\end{table}

\noindent List assignment is determined automatically from the URL query
parameter \texttt{?list=N} (0--3); if absent, a random list is selected.
Each participant is thus within-subjects across items and sees all four
conditions (twice each). Condition effects are estimated from within-participant
variation across the 2$\times$2 design, controlling for by-participant
random intercepts.

In addition to the 8 critical trials, each participant completed 8
\textbf{filler trials} drawn from separate topic domains (train travel,
reading light, cooking, coffee, vacation, hydration, fresh air, noise).
Filler items used the same trial format but were excluded from the
main analysis; they serve to dilute the critical items and reduce
potential demand characteristics.

\subsubsection{Dependent variable}

The dependent variable was the marker chosen from a drop-down menu
displaying all five options in fixed order (as listed in
Section~\ref{sub:markers}).
No response time threshold was imposed, but response times were logged.
Recorded per trial: \texttt{trial\_id}, \texttt{topic}, \texttt{is\_filler},
\texttt{condition\_index}, \texttt{pc\_prag}, \texttt{g},
\texttt{selected\_marker}, \texttt{rt} (ms); and once per participant:
\texttt{list\_num}.
The continuous \texttt{pc\_prop\_rating} covariate (topic-level norming mean)
is joined to the trial data during analysis.

\subsection{Procedure}

Participants completed the experiment online via the
magpie platform \citep{kochetkova2024magpie} in a browser.
After reading instructions in German, participants were assigned to one
of 4 counterbalancing lists (via URL parameter or randomly).
They then proceeded through 16 trials---8 critical and 8 fillers---in a
fully randomised order, followed by a post-test questionnaire eliciting
native language, age, gender, education level, and optional comments.
Each trial displayed (a) a context paragraph, (b) a goal instruction,
and (c) the sentence frame with the marker drop-down embedded at the
Mittelfeld position.
The ``Weiter'' (Next) button was disabled until a marker had been
selected, ensuring a response on every trial.

\end{document}
.
%
% Compile stand-alone:
%   pdflatex production_exp.tex
%
% Embed in main.tex:
%   % production_exp.tex
% Stand-alone description of the production experiment.
% Can be embedded in main.tex via \input{production_exp}.
%
% Compile stand-alone:
%   pdflatex production_exp.tex
%
% Embed in main.tex:
%   \input{production_exp}   % (inside a \section{Experiment} environment)

\documentclass[12pt]{article}
\usepackage{CSPArticleStyle}
\usepackage{csquotes}
\usepackage[margin=1in]{geometry}
\usepackage{booktabs}
\usepackage{parskip}
\usepackage{amsmath}
\usepackage{times}
\usepackage[T1]{fontenc}
\usepackage[utf8]{inputenc}
\usepackage{xcolor}
\usepackage{tabularray}   % robust tables; fall back to tabular if unavailable
\usepackage{gb4e}         % linguistic examples
\noautomath

\title{Production Experiment: Eliciting Consensus-Marker Choice\\
       Across Perceived-Controversy and Goal-Strength Conditions}
\author{}
\date{}

\begin{document}

%% ============================================================
%% When embedding in main.tex, delete from \documentclass to
%% \begin{document} and delete \end{document} at the bottom.
%% ============================================================

\section{Experiment: Consensus-Marker Production}

\subsection{Overview}

The production experiment empirically estimates the decision thresholds
that govern the use of five German consensus markers by placing
participants in the speaker role.
Participants read short vignettes that independently varied
\emph{pragmatic perceived controversy} (\textsc{pc\textsubscript{prag}})
and the speaker's \emph{persuasive-goal strength} ($g$), then selected
which consensus marker they would most naturally use in the described
situation.
\emph{Propositional perceived controversy} (\textsc{pc\textsubscript{prop}})
is held constant within the experiment and operationalised as a continuous
topic-level covariate obtained from a separate norming study: for each
of the eight topic items, participants rated the degree to which the
proposition $p$ is considered established/accepted (0--100 scale);
the resulting means enter the regression analyses as a fixed covariate.
The resulting distribution of marker choices across the
$\textsc{pc}_\text{prag} \times g$ design, combined with the
continuous \textsc{pc\textsubscript{prop}} covariate, provides the
empirical basis for estimating the ordinal thresholds
$\mu_1 < \cdots < \mu_4$ of the noisy-threshold RSA model.

\subsection{Participants}

Participants will be recruited via Prolific and will be required to be
native speakers of German (language-at-home filter).
Data collection is ongoing; a target of $N = 80$ participants is planned
($n = 20$ per list across 4 lists), each completing 16 trials (8 critical + 8 filler).

\subsection{Materials}

\subsubsection{Markers}

The five consensus markers under investigation are, from weakest to
strongest presumed common-ground commitment:

\begin{enumerate}
  \item \textit{sofern ich wei\ss{}} (`as far as I know')
  \item \textit{wie du ja wei\ss t} (`as you (prt) know')
  \item \textit{wie wir wissen} (`as we know')
  \item \textit{ja} (modal particle; roughly `as you know / of course')
  \item \textit{bekanntlich} (`as is well known / notoriously')
\end{enumerate}

These markers were embedded in the sentential \emph{Mittelfeld}
(middle field), immediately after the finite verb, as illustrated in
example~(\ref{ex:frame}):

\begin{exe}
  \ex \label{ex:frame}
  Digitale Kompetenzen sind \textbf{[MARKER]}
  in der heutigen Arbeitswelt unverzichtbar.\\
  `Digital skills are \textbf{[MARKER]} indispensable in today's
  working world.'
\end{exe}

This syntactic position is licit for all five markers: modal particles
such as \textit{ja} canonically occupy the left edge of the Mittelfeld
\citep{thurmair1989modalpartikeln}, while the adverb \textit{bekanntlich}
and parenthetical clauses (\textit{wie wir wissen}, etc.) can appear
there as well (see \citealt{jacobs_semantics_1991} for discussion).
Each target sentence was followed by a consequent clause with the
recommended action $q$: \textit{Deshalb solltest du [q].}
(`Therefore, you should [q].')

\subsubsection{Topic Items}

Eight topic domains were used, each associated with a fixed proposition
$p$ (sentence frame) and consequent action $q$ (Table~\ref{tab:items}).
The first four items were used in earlier pilot work; items 5--8 were
added to reach the minimum of eight required for a complete Latin square.

\begin{table}[ht]
\centering
\caption{Topic items: sentence frame for $p$ (with marker slot [M]) and consequent $q$.}
\label{tab:items}
\small
\begin{tabular}{clll}
\toprule
\textbf{\#} & \textbf{Topic} & \textbf{Sentence frame} & \textbf{Consequent $q$} \\
\midrule
1 & Klimawandel &
  \textit{Der Klimawandel wird [M]} &
  auf Flugreisen verzichten \\
  & & \textit{durch menschliche Aktivit\"aten verursacht.} & \\[3pt]
2 & Ern\"ahrung &
  \textit{Eine \"uberwiegend pflanzliche Ern\"ahrung} &
  mehr pflanzliche Lebens- \\
  & & \textit{f\"ordert [M] die Gesundheit.} &
  mittel einbauen \\[3pt]
3 & Stadtverkehr &
  \textit{Ein ausgebautes Nahverkehrsnetz} &
  h\"aufiger auf \"offentliche \\
  & & \textit{entlastet [M] den Stadtverkehr deutlich.} &
  Verkehrsmittel umsteigen \\[3pt]
4 & Digitalisierung &
  \textit{Digitale Kompetenzen sind [M] in der} &
  regelm\"a\ss{}ig Weiterbildungs- \\
  & & \textit{heutigen Arbeitswelt unverzichtbar.} &
  angebote nutzen \\[3pt]
5 & Schlaf &
  \textit{Ausreichend Schlaf ist [M] entscheidend} &
  auf feste Schlafzeiten \\
  & & \textit{f\"ur die kognitive Leistungsf\"ahigkeit.} &
  achten \\[3pt]
6 & Plastik &
  \textit{Einwegplastik schadet [M] den} &
  konsequent auf Einweg- \\
  & & \textit{Meeres\"okosystemen erheblich.} &
  plastik verzichten \\[3pt]
7 & Sport &
  \textit{Regelm\"a\ss{}ige k\"orperliche Bewegung} &
  regelm\"a\ss{}ig Sport in \\
  & & \textit{senkt [M] das Risiko f\"ur} &
  den Alltag integrieren \\
  & & \textit{Herz-Kreislauf-Erkrankungen deutlich.} & \\[3pt]
8 & Lokalkauf &
  \textit{Das Kaufen bei lokalen H\"andlern} &
  \"ofter bei lokalen \\
  & & \textit{st\"arkt [M] die regionale Wirtschaft} &
  Gesch\"aften einkaufen \\
  & & \textit{nachhaltig.} & \\
\bottomrule
\end{tabular}
\end{table}

\noindent [M] denotes the marker slot (dropdown).

\subsubsection{Vignettes}

\paragraph{Propositional perceived controversy (\textsc{pc\textsubscript{prop}}).}
\textsc{pc\textsubscript{prop}} was \emph{not} manipulated within the
production experiment.
Instead, it was operationalised as a continuous, topic-level covariate
obtained from a separate norming study conducted prior to the production
experiment.
In that norming study, a different sample of native German speakers rated,
for each topic proposition $p$, the degree to which they considered $p$
accepted/established in (German) public discourse, on a 0--100 scale
($0 = $ `completely contested', $100 = $ `universally accepted').
The resulting per-topic means (e.g., \textit{Digitalisierung}: 85/100;
\textit{Lokalkauf}: 58/100) enter the regression analyses as a
continuous, centred predictor \textsc{pc\textsubscript{prop}}$_c$
(mean-centred, SD $\approx$ 1).
Context paragraphs in the production experiment therefore present
the proposition $p$ without any editorial evaluation, so that
participants' experience of propositional controversy is shaped
solely by their own prior knowledge and the norming-derived covariate
captures this variation in analysis.

\paragraph{Pragmatic perceived controversy (\textsc{pc\textsubscript{prag}}).}
\textsc{pc\textsubscript{prag}} was manipulated between instances of each item
by varying whether the speaker's social circle \emph{largely agreed} (low)
or \emph{held divergent views} (high) about $p$.
The two context versions for the \textit{Digitalisierung} item are:

\begin{itemize}
  \item $\textsc{pc}_\text{prag} = \text{low}$: \textit{Du sprichst mit einer
        Kollegin {\"u}ber Berufsausbildung und Weiterbildung.
        In eurem Team teilen die meisten die Ansicht,
        dass digitale Kompetenzen in der heutigen Arbeitswelt unverzichtbar sind.}
  \item $\textsc{pc}_\text{prag} = \text{high}$: \textit{Du sprichst mit einer
        Kollegin {\"u}ber Berufsausbildung und Weiterbildung.
        In eurem Team gehen die Meinungen dazu stark auseinander.}
\end{itemize}

\paragraph{Speaker goal strength ($g$).}
Goal strength was manipulated via a standardised goal instruction
appended after the context paragraph:

\begin{itemize}
  \item $g_\text{low}$: \textit{Du m\"ochtest das nur kurz anmerken.}
        (`You just want to make a brief note of this.')
  \item $g_\text{high}$: \textit{Es ist dir sehr wichtig, dass dein
        Gespr\"achspartner das akzeptiert.}
        (`It is very important to you that your interlocutor accepts this.')
\end{itemize}

\paragraph{Complete sample item (all four conditions).}
Examples (\ref{ex:cond1})--(\ref{ex:cond4}) show the full vignette for
the topic \textit{Digitalisierung}
(norming $\textsc{pc}_\text{prop}$ rating: 85/100;
sentence frame: \textit{Digitale Kompetenzen sind [M] in der heutigen
Arbeitswelt unverzichtbar.
Deshalb solltest du regelm\"a\ss{}ig digitale Weiterbildungsangebote nutzen.})
across all four conditions of the $2\times2$ design.

\begin{exe}
  \ex \label{ex:cond1}
  \textbf{Condition 0: $\textsc{pc}_\text{prag} = \text{low}$, $g = \text{low}$}\\
  \textit{Du sprichst mit einer Kollegin {\"u}ber Berufsausbildung und Weiterbildung.
  In eurem Team teilen die meisten die Ansicht,
  dass digitale Kompetenzen in der heutigen Arbeitswelt unverzichtbar sind.}\\
  \textit{Du m\"ochtest das nur kurz anmerken.}\\
  Digitale Kompetenzen sind \underline{\hspace{2.2cm}}
  in der heutigen Arbeitswelt unverzichtbar.

  \ex \label{ex:cond2}
  \textbf{Condition 1: $\textsc{pc}_\text{prag} = \text{low}$, $g = \text{high}$}\\
  \textit{Du sprichst mit einer Kollegin {\"u}ber Berufsausbildung und Weiterbildung.
  In eurem Team teilen die meisten die Ansicht,
  dass digitale Kompetenzen in der heutigen Arbeitswelt unverzichtbar sind.}\\
  \textit{Es ist dir sehr wichtig, dass dein Gespr\"achspartner das akzeptiert.}\\
  Digitale Kompetenzen sind \underline{\hspace{2.2cm}}
  in der heutigen Arbeitswelt unverzichtbar.

  \ex \label{ex:cond3}
  \textbf{Condition 2: $\textsc{pc}_\text{prag} = \text{high}$, $g = \text{low}$}\\
  \textit{Du sprichst mit einer Kollegin {\"u}ber Berufsausbildung und Weiterbildung.
  In eurem Team gehen die Meinungen dazu stark auseinander.}\\
  \textit{Du m\"ochtest das nur kurz anmerken.}\\
  Digitale Kompetenzen sind \underline{\hspace{2.2cm}}
  in der heutigen Arbeitswelt unverzichtbar.

  \ex \label{ex:cond4}
  \textbf{Condition 3: $\textsc{pc}_\text{prag} = \text{high}$, $g = \text{high}$}\\
  \textit{Du sprichst mit einer Kollegin {\"u}ber Berufsausbildung und Weiterbildung.
  In eurem Team gehen die Meinungen dazu stark auseinander.}\\
  \textit{Es ist dir sehr wichtig, dass dein Gespr\"achspartner das akzeptiert.}\\
  Digitale Kompetenzen sind \underline{\hspace{2.2cm}}
  in der heutigen Arbeitswelt unverzichtbar.
\end{exe}

\noindent In all four vignettes the sentence frame and consequent clause
(\textit{Deshalb solltest du regelm\"a\ss{}ig digitale
Weiterbildungsangebote nutzen.}) are identical; only the context
paragraph and goal instruction vary across conditions.
The marker slot [M] / \underline{\hspace{2.2cm}} was presented as a
drop-down menu in the actual experiment.

\subsection{Design}

\subsubsection{Factorial structure}

The experiment employed a \textbf{2 (pc\textsubscript{prag}: low vs.\ high)
$\times$ 2 ($g$: low vs.\ high) fully crossed design},
yielding 4 conditions.
\textsc{pc\textsubscript{prop}} is not a within-experiment factor; it enters
as a continuous topic-level covariate from the norming study.
With 8 topic items and a Latin square rotation across 4 lists, each
participant saw exactly \textbf{8 critical trials}---two items per condition---plus
8 filler trials, for \textbf{16 trials} in total
(Table~\ref{tab:design}).

\begin{table}[ht]
\centering
\caption{%
  The $2\times2$ design (4 conditions).
  \textsc{pc\textsubscript{prop}} is a continuous norming covariate, not a
  within-experiment factor.
  In each list, the 8 items are assigned to conditions via the Latin square
  rotation described in Section~\ref{subsub:counterbalancing}.}
\label{tab:design}
\small
\begin{tabular}{ccl}
\toprule
\textbf{pc\textsubscript{prag}} &
\textbf{$g$} &
\textbf{Condition label} \\
\midrule
low  & low  & (0) community agrees, casual \\
low  & high & (1) community agrees, pressing \\
high & low  & (2) community divided, casual \\
high & high & (3) community divided, pressing \\
\bottomrule
\end{tabular}
\end{table}

\subsubsection{Counterbalancing}
\label{subsub:counterbalancing}

The experiment used a \textbf{Latin square} to counterbalance item--condition
pairings across participants.
With $k = 4$ conditions and 8 items, a balanced assignment is achieved:
each of the 4 lists assigns item $i$ ($i = 0,\ldots,7$) to condition
$(i + l)\bmod 4$, where $l \in \{0,1,2,3\}$ is the list number.
The resulting assignment matrix (Table~\ref{tab:ls}) guarantees that
\begin{itemize}
  \item each participant sees every item exactly once;
  \item each participant sees each condition exactly twice (two items per
        condition), ensuring balanced coverage of the 2$\times$2 design;
  \item across all 4 lists, each item appears equally often in each
        condition (twice per condition over the full participant sample).
\end{itemize}

\begin{table}[ht]
\centering
\caption{Latin square assignment of conditions to items across 4 lists.
  Cell entries are condition indices (0--3) as defined in
  Table~\ref{tab:design}.
  Rows = lists ($l$); columns = items ($i$).}
\label{tab:ls}
\small
\begin{tabular}{ccccccccc}
\toprule
\textbf{List} &
\textbf{Item 1} & \textbf{Item 2} & \textbf{Item 3} & \textbf{Item 4} &
\textbf{Item 5} & \textbf{Item 6} & \textbf{Item 7} & \textbf{Item 8} \\
\midrule
0 & 0 & 1 & 2 & 3 & 0 & 1 & 2 & 3 \\
1 & 1 & 2 & 3 & 0 & 1 & 2 & 3 & 0 \\
2 & 2 & 3 & 0 & 1 & 2 & 3 & 0 & 1 \\
3 & 3 & 0 & 1 & 2 & 3 & 0 & 1 & 2 \\
\bottomrule
\end{tabular}
\end{table}

\noindent List assignment is determined automatically from the URL query
parameter \texttt{?list=N} (0--3); if absent, a random list is selected.
Each participant is thus within-subjects across items and sees all four
conditions (twice each). Condition effects are estimated from within-participant
variation across the 2$\times$2 design, controlling for by-participant
random intercepts.

In addition to the 8 critical trials, each participant completed 8
\textbf{filler trials} drawn from separate topic domains (train travel,
reading light, cooking, coffee, vacation, hydration, fresh air, noise).
Filler items used the same trial format but were excluded from the
main analysis; they serve to dilute the critical items and reduce
potential demand characteristics.

\subsubsection{Dependent variable}

The dependent variable was the marker chosen from a drop-down menu
displaying all five options in fixed order (as listed in
Section~\ref{sub:markers}).
No response time threshold was imposed, but response times were logged.
Recorded per trial: \texttt{trial\_id}, \texttt{topic}, \texttt{is\_filler},
\texttt{condition\_index}, \texttt{pc\_prag}, \texttt{g},
\texttt{selected\_marker}, \texttt{rt} (ms); and once per participant:
\texttt{list\_num}.
The continuous \texttt{pc\_prop\_rating} covariate (topic-level norming mean)
is joined to the trial data during analysis.

\subsection{Procedure}

Participants completed the experiment online via the
magpie platform \citep{kochetkova2024magpie} in a browser.
After reading instructions in German, participants were assigned to one
of 4 counterbalancing lists (via URL parameter or randomly).
They then proceeded through 16 trials---8 critical and 8 fillers---in a
fully randomised order, followed by a post-test questionnaire eliciting
native language, age, gender, education level, and optional comments.
Each trial displayed (a) a context paragraph, (b) a goal instruction,
and (c) the sentence frame with the marker drop-down embedded at the
Mittelfeld position.
The ``Weiter'' (Next) button was disabled until a marker had been
selected, ensuring a response on every trial.

\end{document}
   % (inside a \section{Experiment} environment)

\documentclass[12pt]{article}
\usepackage{CSPArticleStyle}
\usepackage{csquotes}
\usepackage[margin=1in]{geometry}
\usepackage{booktabs}
\usepackage{parskip}
\usepackage{amsmath}
\usepackage{times}
\usepackage[T1]{fontenc}
\usepackage[utf8]{inputenc}
\usepackage{xcolor}
\usepackage{tabularray}   % robust tables; fall back to tabular if unavailable
\usepackage{gb4e}         % linguistic examples
\noautomath

\title{Production Experiment: Eliciting Consensus-Marker Choice\\
       Across Perceived-Controversy and Goal-Strength Conditions}
\author{}
\date{}

\begin{document}

%% ============================================================
%% When embedding in main.tex, delete from \documentclass to
%% \begin{document} and delete \end{document} at the bottom.
%% ============================================================

\section{Experiment: Consensus-Marker Production}

\subsection{Overview}

The production experiment empirically estimates the decision thresholds
that govern the use of five German consensus markers by placing
participants in the speaker role.
Participants read short vignettes that independently varied
\emph{pragmatic perceived controversy} (\textsc{pc\textsubscript{prag}})
and the speaker's \emph{persuasive-goal strength} ($g$), then selected
which consensus marker they would most naturally use in the described
situation.
\emph{Propositional perceived controversy} (\textsc{pc\textsubscript{prop}})
is held constant within the experiment and operationalised as a continuous
topic-level covariate obtained from a separate norming study: for each
of the eight topic items, participants rated the degree to which the
proposition $p$ is considered established/accepted (0--100 scale);
the resulting means enter the regression analyses as a fixed covariate.
The resulting distribution of marker choices across the
$\textsc{pc}_\text{prag} \times g$ design, combined with the
continuous \textsc{pc\textsubscript{prop}} covariate, provides the
empirical basis for estimating the ordinal thresholds
$\mu_1 < \cdots < \mu_4$ of the noisy-threshold RSA model.

\subsection{Participants}

Participants will be recruited via Prolific and will be required to be
native speakers of German (language-at-home filter).
Data collection is ongoing; a target of $N = 80$ participants is planned
($n = 20$ per list across 4 lists), each completing 16 trials (8 critical + 8 filler).

\subsection{Materials}

\subsubsection{Markers}

The five consensus markers under investigation are, from weakest to
strongest presumed common-ground commitment:

\begin{enumerate}
  \item \textit{sofern ich wei\ss{}} (`as far as I know')
  \item \textit{wie du ja wei\ss t} (`as you (prt) know')
  \item \textit{wie wir wissen} (`as we know')
  \item \textit{ja} (modal particle; roughly `as you know / of course')
  \item \textit{bekanntlich} (`as is well known / notoriously')
\end{enumerate}

These markers were embedded in the sentential \emph{Mittelfeld}
(middle field), immediately after the finite verb, as illustrated in
example~(\ref{ex:frame}):

\begin{exe}
  \ex \label{ex:frame}
  Digitale Kompetenzen sind \textbf{[MARKER]}
  in der heutigen Arbeitswelt unverzichtbar.\\
  `Digital skills are \textbf{[MARKER]} indispensable in today's
  working world.'
\end{exe}

This syntactic position is licit for all five markers: modal particles
such as \textit{ja} canonically occupy the left edge of the Mittelfeld
\citep{thurmair1989modalpartikeln}, while the adverb \textit{bekanntlich}
and parenthetical clauses (\textit{wie wir wissen}, etc.) can appear
there as well (see \citealt{jacobs_semantics_1991} for discussion).
Each target sentence was followed by a consequent clause with the
recommended action $q$: \textit{Deshalb solltest du [q].}
(`Therefore, you should [q].')

\subsubsection{Topic Items}

Eight topic domains were used, each associated with a fixed proposition
$p$ (sentence frame) and consequent action $q$ (Table~\ref{tab:items}).
The first four items were used in earlier pilot work; items 5--8 were
added to reach the minimum of eight required for a complete Latin square.

\begin{table}[ht]
\centering
\caption{Topic items: sentence frame for $p$ (with marker slot [M]) and consequent $q$.}
\label{tab:items}
\small
\begin{tabular}{clll}
\toprule
\textbf{\#} & \textbf{Topic} & \textbf{Sentence frame} & \textbf{Consequent $q$} \\
\midrule
1 & Klimawandel &
  \textit{Der Klimawandel wird [M]} &
  auf Flugreisen verzichten \\
  & & \textit{durch menschliche Aktivit\"aten verursacht.} & \\[3pt]
2 & Ern\"ahrung &
  \textit{Eine \"uberwiegend pflanzliche Ern\"ahrung} &
  mehr pflanzliche Lebens- \\
  & & \textit{f\"ordert [M] die Gesundheit.} &
  mittel einbauen \\[3pt]
3 & Stadtverkehr &
  \textit{Ein ausgebautes Nahverkehrsnetz} &
  h\"aufiger auf \"offentliche \\
  & & \textit{entlastet [M] den Stadtverkehr deutlich.} &
  Verkehrsmittel umsteigen \\[3pt]
4 & Digitalisierung &
  \textit{Digitale Kompetenzen sind [M] in der} &
  regelm\"a\ss{}ig Weiterbildungs- \\
  & & \textit{heutigen Arbeitswelt unverzichtbar.} &
  angebote nutzen \\[3pt]
5 & Schlaf &
  \textit{Ausreichend Schlaf ist [M] entscheidend} &
  auf feste Schlafzeiten \\
  & & \textit{f\"ur die kognitive Leistungsf\"ahigkeit.} &
  achten \\[3pt]
6 & Plastik &
  \textit{Einwegplastik schadet [M] den} &
  konsequent auf Einweg- \\
  & & \textit{Meeres\"okosystemen erheblich.} &
  plastik verzichten \\[3pt]
7 & Sport &
  \textit{Regelm\"a\ss{}ige k\"orperliche Bewegung} &
  regelm\"a\ss{}ig Sport in \\
  & & \textit{senkt [M] das Risiko f\"ur} &
  den Alltag integrieren \\
  & & \textit{Herz-Kreislauf-Erkrankungen deutlich.} & \\[3pt]
8 & Lokalkauf &
  \textit{Das Kaufen bei lokalen H\"andlern} &
  \"ofter bei lokalen \\
  & & \textit{st\"arkt [M] die regionale Wirtschaft} &
  Gesch\"aften einkaufen \\
  & & \textit{nachhaltig.} & \\
\bottomrule
\end{tabular}
\end{table}

\noindent [M] denotes the marker slot (dropdown).

\subsubsection{Vignettes}

\paragraph{Propositional perceived controversy (\textsc{pc\textsubscript{prop}}).}
\textsc{pc\textsubscript{prop}} was \emph{not} manipulated within the
production experiment.
Instead, it was operationalised as a continuous, topic-level covariate
obtained from a separate norming study conducted prior to the production
experiment.
In that norming study, a different sample of native German speakers rated,
for each topic proposition $p$, the degree to which they considered $p$
accepted/established in (German) public discourse, on a 0--100 scale
($0 = $ `completely contested', $100 = $ `universally accepted').
The resulting per-topic means (e.g., \textit{Digitalisierung}: 85/100;
\textit{Lokalkauf}: 58/100) enter the regression analyses as a
continuous, centred predictor \textsc{pc\textsubscript{prop}}$_c$
(mean-centred, SD $\approx$ 1).
Context paragraphs in the production experiment therefore present
the proposition $p$ without any editorial evaluation, so that
participants' experience of propositional controversy is shaped
solely by their own prior knowledge and the norming-derived covariate
captures this variation in analysis.

\paragraph{Pragmatic perceived controversy (\textsc{pc\textsubscript{prag}}).}
\textsc{pc\textsubscript{prag}} was manipulated between instances of each item
by varying whether the speaker's social circle \emph{largely agreed} (low)
or \emph{held divergent views} (high) about $p$.
The two context versions for the \textit{Digitalisierung} item are:

\begin{itemize}
  \item $\textsc{pc}_\text{prag} = \text{low}$: \textit{Du sprichst mit einer
        Kollegin {\"u}ber Berufsausbildung und Weiterbildung.
        In eurem Team teilen die meisten die Ansicht,
        dass digitale Kompetenzen in der heutigen Arbeitswelt unverzichtbar sind.}
  \item $\textsc{pc}_\text{prag} = \text{high}$: \textit{Du sprichst mit einer
        Kollegin {\"u}ber Berufsausbildung und Weiterbildung.
        In eurem Team gehen die Meinungen dazu stark auseinander.}
\end{itemize}

\paragraph{Speaker goal strength ($g$).}
Goal strength was manipulated via a standardised goal instruction
appended after the context paragraph:

\begin{itemize}
  \item $g_\text{low}$: \textit{Du m\"ochtest das nur kurz anmerken.}
        (`You just want to make a brief note of this.')
  \item $g_\text{high}$: \textit{Es ist dir sehr wichtig, dass dein
        Gespr\"achspartner das akzeptiert.}
        (`It is very important to you that your interlocutor accepts this.')
\end{itemize}

\paragraph{Complete sample item (all four conditions).}
Examples (\ref{ex:cond1})--(\ref{ex:cond4}) show the full vignette for
the topic \textit{Digitalisierung}
(norming $\textsc{pc}_\text{prop}$ rating: 85/100;
sentence frame: \textit{Digitale Kompetenzen sind [M] in der heutigen
Arbeitswelt unverzichtbar.
Deshalb solltest du regelm\"a\ss{}ig digitale Weiterbildungsangebote nutzen.})
across all four conditions of the $2\times2$ design.

\begin{exe}
  \ex \label{ex:cond1}
  \textbf{Condition 0: $\textsc{pc}_\text{prag} = \text{low}$, $g = \text{low}$}\\
  \textit{Du sprichst mit einer Kollegin {\"u}ber Berufsausbildung und Weiterbildung.
  In eurem Team teilen die meisten die Ansicht,
  dass digitale Kompetenzen in der heutigen Arbeitswelt unverzichtbar sind.}\\
  \textit{Du m\"ochtest das nur kurz anmerken.}\\
  Digitale Kompetenzen sind \underline{\hspace{2.2cm}}
  in der heutigen Arbeitswelt unverzichtbar.

  \ex \label{ex:cond2}
  \textbf{Condition 1: $\textsc{pc}_\text{prag} = \text{low}$, $g = \text{high}$}\\
  \textit{Du sprichst mit einer Kollegin {\"u}ber Berufsausbildung und Weiterbildung.
  In eurem Team teilen die meisten die Ansicht,
  dass digitale Kompetenzen in der heutigen Arbeitswelt unverzichtbar sind.}\\
  \textit{Es ist dir sehr wichtig, dass dein Gespr\"achspartner das akzeptiert.}\\
  Digitale Kompetenzen sind \underline{\hspace{2.2cm}}
  in der heutigen Arbeitswelt unverzichtbar.

  \ex \label{ex:cond3}
  \textbf{Condition 2: $\textsc{pc}_\text{prag} = \text{high}$, $g = \text{low}$}\\
  \textit{Du sprichst mit einer Kollegin {\"u}ber Berufsausbildung und Weiterbildung.
  In eurem Team gehen die Meinungen dazu stark auseinander.}\\
  \textit{Du m\"ochtest das nur kurz anmerken.}\\
  Digitale Kompetenzen sind \underline{\hspace{2.2cm}}
  in der heutigen Arbeitswelt unverzichtbar.

  \ex \label{ex:cond4}
  \textbf{Condition 3: $\textsc{pc}_\text{prag} = \text{high}$, $g = \text{high}$}\\
  \textit{Du sprichst mit einer Kollegin {\"u}ber Berufsausbildung und Weiterbildung.
  In eurem Team gehen die Meinungen dazu stark auseinander.}\\
  \textit{Es ist dir sehr wichtig, dass dein Gespr\"achspartner das akzeptiert.}\\
  Digitale Kompetenzen sind \underline{\hspace{2.2cm}}
  in der heutigen Arbeitswelt unverzichtbar.
\end{exe}

\noindent In all four vignettes the sentence frame and consequent clause
(\textit{Deshalb solltest du regelm\"a\ss{}ig digitale
Weiterbildungsangebote nutzen.}) are identical; only the context
paragraph and goal instruction vary across conditions.
The marker slot [M] / \underline{\hspace{2.2cm}} was presented as a
drop-down menu in the actual experiment.

\subsection{Design}

\subsubsection{Factorial structure}

The experiment employed a \textbf{2 (pc\textsubscript{prag}: low vs.\ high)
$\times$ 2 ($g$: low vs.\ high) fully crossed design},
yielding 4 conditions.
\textsc{pc\textsubscript{prop}} is not a within-experiment factor; it enters
as a continuous topic-level covariate from the norming study.
With 8 topic items and a Latin square rotation across 4 lists, each
participant saw exactly \textbf{8 critical trials}---two items per condition---plus
8 filler trials, for \textbf{16 trials} in total
(Table~\ref{tab:design}).

\begin{table}[ht]
\centering
\caption{%
  The $2\times2$ design (4 conditions).
  \textsc{pc\textsubscript{prop}} is a continuous norming covariate, not a
  within-experiment factor.
  In each list, the 8 items are assigned to conditions via the Latin square
  rotation described in Section~\ref{subsub:counterbalancing}.}
\label{tab:design}
\small
\begin{tabular}{ccl}
\toprule
\textbf{pc\textsubscript{prag}} &
\textbf{$g$} &
\textbf{Condition label} \\
\midrule
low  & low  & (0) community agrees, casual \\
low  & high & (1) community agrees, pressing \\
high & low  & (2) community divided, casual \\
high & high & (3) community divided, pressing \\
\bottomrule
\end{tabular}
\end{table}

\subsubsection{Counterbalancing}
\label{subsub:counterbalancing}

The experiment used a \textbf{Latin square} to counterbalance item--condition
pairings across participants.
With $k = 4$ conditions and 8 items, a balanced assignment is achieved:
each of the 4 lists assigns item $i$ ($i = 0,\ldots,7$) to condition
$(i + l)\bmod 4$, where $l \in \{0,1,2,3\}$ is the list number.
The resulting assignment matrix (Table~\ref{tab:ls}) guarantees that
\begin{itemize}
  \item each participant sees every item exactly once;
  \item each participant sees each condition exactly twice (two items per
        condition), ensuring balanced coverage of the 2$\times$2 design;
  \item across all 4 lists, each item appears equally often in each
        condition (twice per condition over the full participant sample).
\end{itemize}

\begin{table}[ht]
\centering
\caption{Latin square assignment of conditions to items across 4 lists.
  Cell entries are condition indices (0--3) as defined in
  Table~\ref{tab:design}.
  Rows = lists ($l$); columns = items ($i$).}
\label{tab:ls}
\small
\begin{tabular}{ccccccccc}
\toprule
\textbf{List} &
\textbf{Item 1} & \textbf{Item 2} & \textbf{Item 3} & \textbf{Item 4} &
\textbf{Item 5} & \textbf{Item 6} & \textbf{Item 7} & \textbf{Item 8} \\
\midrule
0 & 0 & 1 & 2 & 3 & 0 & 1 & 2 & 3 \\
1 & 1 & 2 & 3 & 0 & 1 & 2 & 3 & 0 \\
2 & 2 & 3 & 0 & 1 & 2 & 3 & 0 & 1 \\
3 & 3 & 0 & 1 & 2 & 3 & 0 & 1 & 2 \\
\bottomrule
\end{tabular}
\end{table}

\noindent List assignment is determined automatically from the URL query
parameter \texttt{?list=N} (0--3); if absent, a random list is selected.
Each participant is thus within-subjects across items and sees all four
conditions (twice each). Condition effects are estimated from within-participant
variation across the 2$\times$2 design, controlling for by-participant
random intercepts.

In addition to the 8 critical trials, each participant completed 8
\textbf{filler trials} drawn from separate topic domains (train travel,
reading light, cooking, coffee, vacation, hydration, fresh air, noise).
Filler items used the same trial format but were excluded from the
main analysis; they serve to dilute the critical items and reduce
potential demand characteristics.

\subsubsection{Dependent variable}

The dependent variable was the marker chosen from a drop-down menu
displaying all five options in fixed order (as listed in
Section~\ref{sub:markers}).
No response time threshold was imposed, but response times were logged.
Recorded per trial: \texttt{trial\_id}, \texttt{topic}, \texttt{is\_filler},
\texttt{condition\_index}, \texttt{pc\_prag}, \texttt{g},
\texttt{selected\_marker}, \texttt{rt} (ms); and once per participant:
\texttt{list\_num}.
The continuous \texttt{pc\_prop\_rating} covariate (topic-level norming mean)
is joined to the trial data during analysis.

\subsection{Procedure}

Participants completed the experiment online via the
magpie platform \citep{kochetkova2024magpie} in a browser.
After reading instructions in German, participants were assigned to one
of 4 counterbalancing lists (via URL parameter or randomly).
They then proceeded through 16 trials---8 critical and 8 fillers---in a
fully randomised order, followed by a post-test questionnaire eliciting
native language, age, gender, education level, and optional comments.
Each trial displayed (a) a context paragraph, (b) a goal instruction,
and (c) the sentence frame with the marker drop-down embedded at the
Mittelfeld position.
The ``Weiter'' (Next) button was disabled until a marker had been
selected, ensuring a response on every trial.

\end{document}
   % (inside a \section{Experiment} environment)

\documentclass[12pt]{article}
\usepackage{CSPArticleStyle}
\usepackage{csquotes}
\usepackage[margin=1in]{geometry}
\usepackage{booktabs}
\usepackage{parskip}
\usepackage{amsmath}
\usepackage{times}
\usepackage[T1]{fontenc}
\usepackage[utf8]{inputenc}
\usepackage{xcolor}
\usepackage{tabularray}   % robust tables; fall back to tabular if unavailable
\usepackage{gb4e}         % linguistic examples
\noautomath

\title{Production Experiment: Eliciting Consensus-Marker Choice\\
       Across Perceived-Controversy and Goal-Strength Conditions}
\author{}
\date{}

\begin{document}

%% ============================================================
%% When embedding in main.tex, delete from \documentclass to
%% \begin{document} and delete \end{document} at the bottom.
%% ============================================================

\section{Experiment: Consensus-Marker Production}

\subsection{Overview}

The production experiment empirically estimates the decision thresholds
that govern the use of five German consensus markers by placing
participants in the speaker role.
Participants read short vignettes that independently varied
\emph{pragmatic perceived controversy} (\textsc{pc\textsubscript{prag}})
and the speaker's \emph{persuasive-goal strength} ($g$), then selected
which consensus marker they would most naturally use in the described
situation.
\emph{Propositional perceived controversy} (\textsc{pc\textsubscript{prop}})
is held constant within the experiment and operationalised as a continuous
topic-level covariate obtained from a separate norming study: for each
of the eight topic items, participants rated the degree to which the
proposition $p$ is considered established/accepted (0--100 scale);
the resulting means enter the regression analyses as a fixed covariate.
The resulting distribution of marker choices across the
$\textsc{pc}_\text{prag} \times g$ design, combined with the
continuous \textsc{pc\textsubscript{prop}} covariate, provides the
empirical basis for estimating the ordinal thresholds
$\mu_1 < \cdots < \mu_4$ of the noisy-threshold RSA model.

\subsection{Participants}

Participants will be recruited via Prolific and will be required to be
native speakers of German (language-at-home filter).
Data collection is ongoing; a target of $N = 80$ participants is planned
($n = 20$ per list across 4 lists), each completing 16 trials (8 critical + 8 filler).

\subsection{Materials}

\subsubsection{Markers}

The five consensus markers under investigation are, from weakest to
strongest presumed common-ground commitment:

\begin{enumerate}
  \item \textit{sofern ich wei\ss{}} (`as far as I know')
  \item \textit{wie du ja wei\ss t} (`as you (prt) know')
  \item \textit{wie wir wissen} (`as we know')
  \item \textit{ja} (modal particle; roughly `as you know / of course')
  \item \textit{bekanntlich} (`as is well known / notoriously')
\end{enumerate}

These markers were embedded in the sentential \emph{Mittelfeld}
(middle field), immediately after the finite verb, as illustrated in
example~(\ref{ex:frame}):

\begin{exe}
  \ex \label{ex:frame}
  Digitale Kompetenzen sind \textbf{[MARKER]}
  in der heutigen Arbeitswelt unverzichtbar.\\
  `Digital skills are \textbf{[MARKER]} indispensable in today's
  working world.'
\end{exe}

This syntactic position is licit for all five markers: modal particles
such as \textit{ja} canonically occupy the left edge of the Mittelfeld
\citep{thurmair1989modalpartikeln}, while the adverb \textit{bekanntlich}
and parenthetical clauses (\textit{wie wir wissen}, etc.) can appear
there as well (see \citealt{jacobs_semantics_1991} for discussion).
Each target sentence was followed by a consequent clause with the
recommended action $q$: \textit{Deshalb solltest du [q].}
(`Therefore, you should [q].')

\subsubsection{Topic Items}

Eight topic domains were used, each associated with a fixed proposition
$p$ (sentence frame) and consequent action $q$ (Table~\ref{tab:items}).
The first four items were used in earlier pilot work; items 5--8 were
added to reach the minimum of eight required for a complete Latin square.

\begin{table}[ht]
\centering
\caption{Topic items: sentence frame for $p$ (with marker slot [M]) and consequent $q$.}
\label{tab:items}
\small
\begin{tabular}{clll}
\toprule
\textbf{\#} & \textbf{Topic} & \textbf{Sentence frame} & \textbf{Consequent $q$} \\
\midrule
1 & Klimawandel &
  \textit{Der Klimawandel wird [M]} &
  auf Flugreisen verzichten \\
  & & \textit{durch menschliche Aktivit\"aten verursacht.} & \\[3pt]
2 & Ern\"ahrung &
  \textit{Eine \"uberwiegend pflanzliche Ern\"ahrung} &
  mehr pflanzliche Lebens- \\
  & & \textit{f\"ordert [M] die Gesundheit.} &
  mittel einbauen \\[3pt]
3 & Stadtverkehr &
  \textit{Ein ausgebautes Nahverkehrsnetz} &
  h\"aufiger auf \"offentliche \\
  & & \textit{entlastet [M] den Stadtverkehr deutlich.} &
  Verkehrsmittel umsteigen \\[3pt]
4 & Digitalisierung &
  \textit{Digitale Kompetenzen sind [M] in der} &
  regelm\"a\ss{}ig Weiterbildungs- \\
  & & \textit{heutigen Arbeitswelt unverzichtbar.} &
  angebote nutzen \\[3pt]
5 & Schlaf &
  \textit{Ausreichend Schlaf ist [M] entscheidend} &
  auf feste Schlafzeiten \\
  & & \textit{f\"ur die kognitive Leistungsf\"ahigkeit.} &
  achten \\[3pt]
6 & Plastik &
  \textit{Einwegplastik schadet [M] den} &
  konsequent auf Einweg- \\
  & & \textit{Meeres\"okosystemen erheblich.} &
  plastik verzichten \\[3pt]
7 & Sport &
  \textit{Regelm\"a\ss{}ige k\"orperliche Bewegung} &
  regelm\"a\ss{}ig Sport in \\
  & & \textit{senkt [M] das Risiko f\"ur} &
  den Alltag integrieren \\
  & & \textit{Herz-Kreislauf-Erkrankungen deutlich.} & \\[3pt]
8 & Lokalkauf &
  \textit{Das Kaufen bei lokalen H\"andlern} &
  \"ofter bei lokalen \\
  & & \textit{st\"arkt [M] die regionale Wirtschaft} &
  Gesch\"aften einkaufen \\
  & & \textit{nachhaltig.} & \\
\bottomrule
\end{tabular}
\end{table}

\noindent [M] denotes the marker slot (dropdown).

\subsubsection{Vignettes}

\paragraph{Propositional perceived controversy (\textsc{pc\textsubscript{prop}}).}
\textsc{pc\textsubscript{prop}} was \emph{not} manipulated within the
production experiment.
Instead, it was operationalised as a continuous, topic-level covariate
obtained from a separate norming study conducted prior to the production
experiment.
In that norming study, a different sample of native German speakers rated,
for each topic proposition $p$, the degree to which they considered $p$
accepted/established in (German) public discourse, on a 0--100 scale
($0 = $ `completely contested', $100 = $ `universally accepted').
The resulting per-topic means (e.g., \textit{Digitalisierung}: 85/100;
\textit{Lokalkauf}: 58/100) enter the regression analyses as a
continuous, centred predictor \textsc{pc\textsubscript{prop}}$_c$
(mean-centred, SD $\approx$ 1).
Context paragraphs in the production experiment therefore present
the proposition $p$ without any editorial evaluation, so that
participants' experience of propositional controversy is shaped
solely by their own prior knowledge and the norming-derived covariate
captures this variation in analysis.

\paragraph{Pragmatic perceived controversy (\textsc{pc\textsubscript{prag}}).}
\textsc{pc\textsubscript{prag}} was manipulated between instances of each item
by varying whether the speaker's social circle \emph{largely agreed} (low)
or \emph{held divergent views} (high) about $p$.
The two context versions for the \textit{Digitalisierung} item are:

\begin{itemize}
  \item $\textsc{pc}_\text{prag} = \text{low}$: \textit{Du sprichst mit einer
        Kollegin {\"u}ber Berufsausbildung und Weiterbildung.
        In eurem Team teilen die meisten die Ansicht,
        dass digitale Kompetenzen in der heutigen Arbeitswelt unverzichtbar sind.}
  \item $\textsc{pc}_\text{prag} = \text{high}$: \textit{Du sprichst mit einer
        Kollegin {\"u}ber Berufsausbildung und Weiterbildung.
        In eurem Team gehen die Meinungen dazu stark auseinander.}
\end{itemize}

\paragraph{Speaker goal strength ($g$).}
Goal strength was manipulated via a standardised goal instruction
appended after the context paragraph:

\begin{itemize}
  \item $g_\text{low}$: \textit{Du m\"ochtest das nur kurz anmerken.}
        (`You just want to make a brief note of this.')
  \item $g_\text{high}$: \textit{Es ist dir sehr wichtig, dass dein
        Gespr\"achspartner das akzeptiert.}
        (`It is very important to you that your interlocutor accepts this.')
\end{itemize}

\paragraph{Complete sample item (all four conditions).}
Examples (\ref{ex:cond1})--(\ref{ex:cond4}) show the full vignette for
the topic \textit{Digitalisierung}
(norming $\textsc{pc}_\text{prop}$ rating: 85/100;
sentence frame: \textit{Digitale Kompetenzen sind [M] in der heutigen
Arbeitswelt unverzichtbar.
Deshalb solltest du regelm\"a\ss{}ig digitale Weiterbildungsangebote nutzen.})
across all four conditions of the $2\times2$ design.

\begin{exe}
  \ex \label{ex:cond1}
  \textbf{Condition 0: $\textsc{pc}_\text{prag} = \text{low}$, $g = \text{low}$}\\
  \textit{Du sprichst mit einer Kollegin {\"u}ber Berufsausbildung und Weiterbildung.
  In eurem Team teilen die meisten die Ansicht,
  dass digitale Kompetenzen in der heutigen Arbeitswelt unverzichtbar sind.}\\
  \textit{Du m\"ochtest das nur kurz anmerken.}\\
  Digitale Kompetenzen sind \underline{\hspace{2.2cm}}
  in der heutigen Arbeitswelt unverzichtbar.

  \ex \label{ex:cond2}
  \textbf{Condition 1: $\textsc{pc}_\text{prag} = \text{low}$, $g = \text{high}$}\\
  \textit{Du sprichst mit einer Kollegin {\"u}ber Berufsausbildung und Weiterbildung.
  In eurem Team teilen die meisten die Ansicht,
  dass digitale Kompetenzen in der heutigen Arbeitswelt unverzichtbar sind.}\\
  \textit{Es ist dir sehr wichtig, dass dein Gespr\"achspartner das akzeptiert.}\\
  Digitale Kompetenzen sind \underline{\hspace{2.2cm}}
  in der heutigen Arbeitswelt unverzichtbar.

  \ex \label{ex:cond3}
  \textbf{Condition 2: $\textsc{pc}_\text{prag} = \text{high}$, $g = \text{low}$}\\
  \textit{Du sprichst mit einer Kollegin {\"u}ber Berufsausbildung und Weiterbildung.
  In eurem Team gehen die Meinungen dazu stark auseinander.}\\
  \textit{Du m\"ochtest das nur kurz anmerken.}\\
  Digitale Kompetenzen sind \underline{\hspace{2.2cm}}
  in der heutigen Arbeitswelt unverzichtbar.

  \ex \label{ex:cond4}
  \textbf{Condition 3: $\textsc{pc}_\text{prag} = \text{high}$, $g = \text{high}$}\\
  \textit{Du sprichst mit einer Kollegin {\"u}ber Berufsausbildung und Weiterbildung.
  In eurem Team gehen die Meinungen dazu stark auseinander.}\\
  \textit{Es ist dir sehr wichtig, dass dein Gespr\"achspartner das akzeptiert.}\\
  Digitale Kompetenzen sind \underline{\hspace{2.2cm}}
  in der heutigen Arbeitswelt unverzichtbar.
\end{exe}

\noindent In all four vignettes the sentence frame and consequent clause
(\textit{Deshalb solltest du regelm\"a\ss{}ig digitale
Weiterbildungsangebote nutzen.}) are identical; only the context
paragraph and goal instruction vary across conditions.
The marker slot [M] / \underline{\hspace{2.2cm}} was presented as a
drop-down menu in the actual experiment.

\subsection{Design}

\subsubsection{Factorial structure}

The experiment employed a \textbf{2 (pc\textsubscript{prag}: low vs.\ high)
$\times$ 2 ($g$: low vs.\ high) fully crossed design},
yielding 4 conditions.
\textsc{pc\textsubscript{prop}} is not a within-experiment factor; it enters
as a continuous topic-level covariate from the norming study.
With 8 topic items and a Latin square rotation across 4 lists, each
participant saw exactly \textbf{8 critical trials}---two items per condition---plus
8 filler trials, for \textbf{16 trials} in total
(Table~\ref{tab:design}).

\begin{table}[ht]
\centering
\caption{%
  The $2\times2$ design (4 conditions).
  \textsc{pc\textsubscript{prop}} is a continuous norming covariate, not a
  within-experiment factor.
  In each list, the 8 items are assigned to conditions via the Latin square
  rotation described in Section~\ref{subsub:counterbalancing}.}
\label{tab:design}
\small
\begin{tabular}{ccl}
\toprule
\textbf{pc\textsubscript{prag}} &
\textbf{$g$} &
\textbf{Condition label} \\
\midrule
low  & low  & (0) community agrees, casual \\
low  & high & (1) community agrees, pressing \\
high & low  & (2) community divided, casual \\
high & high & (3) community divided, pressing \\
\bottomrule
\end{tabular}
\end{table}

\subsubsection{Counterbalancing}
\label{subsub:counterbalancing}

The experiment used a \textbf{Latin square} to counterbalance item--condition
pairings across participants.
With $k = 4$ conditions and 8 items, a balanced assignment is achieved:
each of the 4 lists assigns item $i$ ($i = 0,\ldots,7$) to condition
$(i + l)\bmod 4$, where $l \in \{0,1,2,3\}$ is the list number.
The resulting assignment matrix (Table~\ref{tab:ls}) guarantees that
\begin{itemize}
  \item each participant sees every item exactly once;
  \item each participant sees each condition exactly twice (two items per
        condition), ensuring balanced coverage of the 2$\times$2 design;
  \item across all 4 lists, each item appears equally often in each
        condition (twice per condition over the full participant sample).
\end{itemize}

\begin{table}[ht]
\centering
\caption{Latin square assignment of conditions to items across 4 lists.
  Cell entries are condition indices (0--3) as defined in
  Table~\ref{tab:design}.
  Rows = lists ($l$); columns = items ($i$).}
\label{tab:ls}
\small
\begin{tabular}{ccccccccc}
\toprule
\textbf{List} &
\textbf{Item 1} & \textbf{Item 2} & \textbf{Item 3} & \textbf{Item 4} &
\textbf{Item 5} & \textbf{Item 6} & \textbf{Item 7} & \textbf{Item 8} \\
\midrule
0 & 0 & 1 & 2 & 3 & 0 & 1 & 2 & 3 \\
1 & 1 & 2 & 3 & 0 & 1 & 2 & 3 & 0 \\
2 & 2 & 3 & 0 & 1 & 2 & 3 & 0 & 1 \\
3 & 3 & 0 & 1 & 2 & 3 & 0 & 1 & 2 \\
\bottomrule
\end{tabular}
\end{table}

\noindent List assignment is determined automatically from the URL query
parameter \texttt{?list=N} (0--3); if absent, a random list is selected.
Each participant is thus within-subjects across items and sees all four
conditions (twice each). Condition effects are estimated from within-participant
variation across the 2$\times$2 design, controlling for by-participant
random intercepts.

In addition to the 8 critical trials, each participant completed 8
\textbf{filler trials} drawn from separate topic domains (train travel,
reading light, cooking, coffee, vacation, hydration, fresh air, noise).
Filler items used the same trial format but were excluded from the
main analysis; they serve to dilute the critical items and reduce
potential demand characteristics.

\subsubsection{Dependent variable}

The dependent variable was the marker chosen from a drop-down menu
displaying all five options in fixed order (as listed in
Section~\ref{sub:markers}).
No response time threshold was imposed, but response times were logged.
Recorded per trial: \texttt{trial\_id}, \texttt{topic}, \texttt{is\_filler},
\texttt{condition\_index}, \texttt{pc\_prag}, \texttt{g},
\texttt{selected\_marker}, \texttt{rt} (ms); and once per participant:
\texttt{list\_num}.
The continuous \texttt{pc\_prop\_rating} covariate (topic-level norming mean)
is joined to the trial data during analysis.

\subsection{Procedure}

Participants completed the experiment online via the
magpie platform \citep{kochetkova2024magpie} in a browser.
After reading instructions in German, participants were assigned to one
of 4 counterbalancing lists (via URL parameter or randomly).
They then proceeded through 16 trials---8 critical and 8 fillers---in a
fully randomised order, followed by a post-test questionnaire eliciting
native language, age, gender, education level, and optional comments.
Each trial displayed (a) a context paragraph, (b) a goal instruction,
and (c) the sentence frame with the marker drop-down embedded at the
Mittelfeld position.
The ``Weiter'' (Next) button was disabled until a marker had been
selected, ensuring a response on every trial.

\end{document}
.
%
% Compile stand-alone:
%   pdflatex production_exp.tex
%
% Embed in main.tex:
%   % production_exp.tex
% Stand-alone description of the production experiment.
% Can be embedded in main.tex via % production_exp.tex
% Stand-alone description of the production experiment.
% Can be embedded in main.tex via % production_exp.tex
% Stand-alone description of the production experiment.
% Can be embedded in main.tex via \input{production_exp}.
%
% Compile stand-alone:
%   pdflatex production_exp.tex
%
% Embed in main.tex:
%   \input{production_exp}   % (inside a \section{Experiment} environment)

\documentclass[12pt]{article}
\usepackage{CSPArticleStyle}
\usepackage{csquotes}
\usepackage[margin=1in]{geometry}
\usepackage{booktabs}
\usepackage{parskip}
\usepackage{amsmath}
\usepackage{times}
\usepackage[T1]{fontenc}
\usepackage[utf8]{inputenc}
\usepackage{xcolor}
\usepackage{tabularray}   % robust tables; fall back to tabular if unavailable
\usepackage{gb4e}         % linguistic examples
\noautomath

\title{Production Experiment: Eliciting Consensus-Marker Choice\\
       Across Perceived-Controversy and Goal-Strength Conditions}
\author{}
\date{}

\begin{document}

%% ============================================================
%% When embedding in main.tex, delete from \documentclass to
%% \begin{document} and delete \end{document} at the bottom.
%% ============================================================

\section{Experiment: Consensus-Marker Production}

\subsection{Overview}

The production experiment empirically estimates the decision thresholds
that govern the use of five German consensus markers by placing
participants in the speaker role.
Participants read short vignettes that independently varied
\emph{pragmatic perceived controversy} (\textsc{pc\textsubscript{prag}})
and the speaker's \emph{persuasive-goal strength} ($g$), then selected
which consensus marker they would most naturally use in the described
situation.
\emph{Propositional perceived controversy} (\textsc{pc\textsubscript{prop}})
is held constant within the experiment and operationalised as a continuous
topic-level covariate obtained from a separate norming study: for each
of the eight topic items, participants rated the degree to which the
proposition $p$ is considered established/accepted (0--100 scale);
the resulting means enter the regression analyses as a fixed covariate.
The resulting distribution of marker choices across the
$\textsc{pc}_\text{prag} \times g$ design, combined with the
continuous \textsc{pc\textsubscript{prop}} covariate, provides the
empirical basis for estimating the ordinal thresholds
$\mu_1 < \cdots < \mu_4$ of the noisy-threshold RSA model.

\subsection{Participants}

Participants will be recruited via Prolific and will be required to be
native speakers of German (language-at-home filter).
Data collection is ongoing; a target of $N = 80$ participants is planned
($n = 20$ per list across 4 lists), each completing 16 trials (8 critical + 8 filler).

\subsection{Materials}

\subsubsection{Markers}

The five consensus markers under investigation are, from weakest to
strongest presumed common-ground commitment:

\begin{enumerate}
  \item \textit{sofern ich wei\ss{}} (`as far as I know')
  \item \textit{wie du ja wei\ss t} (`as you (prt) know')
  \item \textit{wie wir wissen} (`as we know')
  \item \textit{ja} (modal particle; roughly `as you know / of course')
  \item \textit{bekanntlich} (`as is well known / notoriously')
\end{enumerate}

These markers were embedded in the sentential \emph{Mittelfeld}
(middle field), immediately after the finite verb, as illustrated in
example~(\ref{ex:frame}):

\begin{exe}
  \ex \label{ex:frame}
  Digitale Kompetenzen sind \textbf{[MARKER]}
  in der heutigen Arbeitswelt unverzichtbar.\\
  `Digital skills are \textbf{[MARKER]} indispensable in today's
  working world.'
\end{exe}

This syntactic position is licit for all five markers: modal particles
such as \textit{ja} canonically occupy the left edge of the Mittelfeld
\citep{thurmair1989modalpartikeln}, while the adverb \textit{bekanntlich}
and parenthetical clauses (\textit{wie wir wissen}, etc.) can appear
there as well (see \citealt{jacobs_semantics_1991} for discussion).
Each target sentence was followed by a consequent clause with the
recommended action $q$: \textit{Deshalb solltest du [q].}
(`Therefore, you should [q].')

\subsubsection{Topic Items}

Eight topic domains were used, each associated with a fixed proposition
$p$ (sentence frame) and consequent action $q$ (Table~\ref{tab:items}).
The first four items were used in earlier pilot work; items 5--8 were
added to reach the minimum of eight required for a complete Latin square.

\begin{table}[ht]
\centering
\caption{Topic items: sentence frame for $p$ (with marker slot [M]) and consequent $q$.}
\label{tab:items}
\small
\begin{tabular}{clll}
\toprule
\textbf{\#} & \textbf{Topic} & \textbf{Sentence frame} & \textbf{Consequent $q$} \\
\midrule
1 & Klimawandel &
  \textit{Der Klimawandel wird [M]} &
  auf Flugreisen verzichten \\
  & & \textit{durch menschliche Aktivit\"aten verursacht.} & \\[3pt]
2 & Ern\"ahrung &
  \textit{Eine \"uberwiegend pflanzliche Ern\"ahrung} &
  mehr pflanzliche Lebens- \\
  & & \textit{f\"ordert [M] die Gesundheit.} &
  mittel einbauen \\[3pt]
3 & Stadtverkehr &
  \textit{Ein ausgebautes Nahverkehrsnetz} &
  h\"aufiger auf \"offentliche \\
  & & \textit{entlastet [M] den Stadtverkehr deutlich.} &
  Verkehrsmittel umsteigen \\[3pt]
4 & Digitalisierung &
  \textit{Digitale Kompetenzen sind [M] in der} &
  regelm\"a\ss{}ig Weiterbildungs- \\
  & & \textit{heutigen Arbeitswelt unverzichtbar.} &
  angebote nutzen \\[3pt]
5 & Schlaf &
  \textit{Ausreichend Schlaf ist [M] entscheidend} &
  auf feste Schlafzeiten \\
  & & \textit{f\"ur die kognitive Leistungsf\"ahigkeit.} &
  achten \\[3pt]
6 & Plastik &
  \textit{Einwegplastik schadet [M] den} &
  konsequent auf Einweg- \\
  & & \textit{Meeres\"okosystemen erheblich.} &
  plastik verzichten \\[3pt]
7 & Sport &
  \textit{Regelm\"a\ss{}ige k\"orperliche Bewegung} &
  regelm\"a\ss{}ig Sport in \\
  & & \textit{senkt [M] das Risiko f\"ur} &
  den Alltag integrieren \\
  & & \textit{Herz-Kreislauf-Erkrankungen deutlich.} & \\[3pt]
8 & Lokalkauf &
  \textit{Das Kaufen bei lokalen H\"andlern} &
  \"ofter bei lokalen \\
  & & \textit{st\"arkt [M] die regionale Wirtschaft} &
  Gesch\"aften einkaufen \\
  & & \textit{nachhaltig.} & \\
\bottomrule
\end{tabular}
\end{table}

\noindent [M] denotes the marker slot (dropdown).

\subsubsection{Vignettes}

\paragraph{Propositional perceived controversy (\textsc{pc\textsubscript{prop}}).}
\textsc{pc\textsubscript{prop}} was \emph{not} manipulated within the
production experiment.
Instead, it was operationalised as a continuous, topic-level covariate
obtained from a separate norming study conducted prior to the production
experiment.
In that norming study, a different sample of native German speakers rated,
for each topic proposition $p$, the degree to which they considered $p$
accepted/established in (German) public discourse, on a 0--100 scale
($0 = $ `completely contested', $100 = $ `universally accepted').
The resulting per-topic means (e.g., \textit{Digitalisierung}: 85/100;
\textit{Lokalkauf}: 58/100) enter the regression analyses as a
continuous, centred predictor \textsc{pc\textsubscript{prop}}$_c$
(mean-centred, SD $\approx$ 1).
Context paragraphs in the production experiment therefore present
the proposition $p$ without any editorial evaluation, so that
participants' experience of propositional controversy is shaped
solely by their own prior knowledge and the norming-derived covariate
captures this variation in analysis.

\paragraph{Pragmatic perceived controversy (\textsc{pc\textsubscript{prag}}).}
\textsc{pc\textsubscript{prag}} was manipulated between instances of each item
by varying whether the speaker's social circle \emph{largely agreed} (low)
or \emph{held divergent views} (high) about $p$.
The two context versions for the \textit{Digitalisierung} item are:

\begin{itemize}
  \item $\textsc{pc}_\text{prag} = \text{low}$: \textit{Du sprichst mit einer
        Kollegin {\"u}ber Berufsausbildung und Weiterbildung.
        In eurem Team teilen die meisten die Ansicht,
        dass digitale Kompetenzen in der heutigen Arbeitswelt unverzichtbar sind.}
  \item $\textsc{pc}_\text{prag} = \text{high}$: \textit{Du sprichst mit einer
        Kollegin {\"u}ber Berufsausbildung und Weiterbildung.
        In eurem Team gehen die Meinungen dazu stark auseinander.}
\end{itemize}

\paragraph{Speaker goal strength ($g$).}
Goal strength was manipulated via a standardised goal instruction
appended after the context paragraph:

\begin{itemize}
  \item $g_\text{low}$: \textit{Du m\"ochtest das nur kurz anmerken.}
        (`You just want to make a brief note of this.')
  \item $g_\text{high}$: \textit{Es ist dir sehr wichtig, dass dein
        Gespr\"achspartner das akzeptiert.}
        (`It is very important to you that your interlocutor accepts this.')
\end{itemize}

\paragraph{Complete sample item (all four conditions).}
Examples (\ref{ex:cond1})--(\ref{ex:cond4}) show the full vignette for
the topic \textit{Digitalisierung}
(norming $\textsc{pc}_\text{prop}$ rating: 85/100;
sentence frame: \textit{Digitale Kompetenzen sind [M] in der heutigen
Arbeitswelt unverzichtbar.
Deshalb solltest du regelm\"a\ss{}ig digitale Weiterbildungsangebote nutzen.})
across all four conditions of the $2\times2$ design.

\begin{exe}
  \ex \label{ex:cond1}
  \textbf{Condition 0: $\textsc{pc}_\text{prag} = \text{low}$, $g = \text{low}$}\\
  \textit{Du sprichst mit einer Kollegin {\"u}ber Berufsausbildung und Weiterbildung.
  In eurem Team teilen die meisten die Ansicht,
  dass digitale Kompetenzen in der heutigen Arbeitswelt unverzichtbar sind.}\\
  \textit{Du m\"ochtest das nur kurz anmerken.}\\
  Digitale Kompetenzen sind \underline{\hspace{2.2cm}}
  in der heutigen Arbeitswelt unverzichtbar.

  \ex \label{ex:cond2}
  \textbf{Condition 1: $\textsc{pc}_\text{prag} = \text{low}$, $g = \text{high}$}\\
  \textit{Du sprichst mit einer Kollegin {\"u}ber Berufsausbildung und Weiterbildung.
  In eurem Team teilen die meisten die Ansicht,
  dass digitale Kompetenzen in der heutigen Arbeitswelt unverzichtbar sind.}\\
  \textit{Es ist dir sehr wichtig, dass dein Gespr\"achspartner das akzeptiert.}\\
  Digitale Kompetenzen sind \underline{\hspace{2.2cm}}
  in der heutigen Arbeitswelt unverzichtbar.

  \ex \label{ex:cond3}
  \textbf{Condition 2: $\textsc{pc}_\text{prag} = \text{high}$, $g = \text{low}$}\\
  \textit{Du sprichst mit einer Kollegin {\"u}ber Berufsausbildung und Weiterbildung.
  In eurem Team gehen die Meinungen dazu stark auseinander.}\\
  \textit{Du m\"ochtest das nur kurz anmerken.}\\
  Digitale Kompetenzen sind \underline{\hspace{2.2cm}}
  in der heutigen Arbeitswelt unverzichtbar.

  \ex \label{ex:cond4}
  \textbf{Condition 3: $\textsc{pc}_\text{prag} = \text{high}$, $g = \text{high}$}\\
  \textit{Du sprichst mit einer Kollegin {\"u}ber Berufsausbildung und Weiterbildung.
  In eurem Team gehen die Meinungen dazu stark auseinander.}\\
  \textit{Es ist dir sehr wichtig, dass dein Gespr\"achspartner das akzeptiert.}\\
  Digitale Kompetenzen sind \underline{\hspace{2.2cm}}
  in der heutigen Arbeitswelt unverzichtbar.
\end{exe}

\noindent In all four vignettes the sentence frame and consequent clause
(\textit{Deshalb solltest du regelm\"a\ss{}ig digitale
Weiterbildungsangebote nutzen.}) are identical; only the context
paragraph and goal instruction vary across conditions.
The marker slot [M] / \underline{\hspace{2.2cm}} was presented as a
drop-down menu in the actual experiment.

\subsection{Design}

\subsubsection{Factorial structure}

The experiment employed a \textbf{2 (pc\textsubscript{prag}: low vs.\ high)
$\times$ 2 ($g$: low vs.\ high) fully crossed design},
yielding 4 conditions.
\textsc{pc\textsubscript{prop}} is not a within-experiment factor; it enters
as a continuous topic-level covariate from the norming study.
With 8 topic items and a Latin square rotation across 4 lists, each
participant saw exactly \textbf{8 critical trials}---two items per condition---plus
8 filler trials, for \textbf{16 trials} in total
(Table~\ref{tab:design}).

\begin{table}[ht]
\centering
\caption{%
  The $2\times2$ design (4 conditions).
  \textsc{pc\textsubscript{prop}} is a continuous norming covariate, not a
  within-experiment factor.
  In each list, the 8 items are assigned to conditions via the Latin square
  rotation described in Section~\ref{subsub:counterbalancing}.}
\label{tab:design}
\small
\begin{tabular}{ccl}
\toprule
\textbf{pc\textsubscript{prag}} &
\textbf{$g$} &
\textbf{Condition label} \\
\midrule
low  & low  & (0) community agrees, casual \\
low  & high & (1) community agrees, pressing \\
high & low  & (2) community divided, casual \\
high & high & (3) community divided, pressing \\
\bottomrule
\end{tabular}
\end{table}

\subsubsection{Counterbalancing}
\label{subsub:counterbalancing}

The experiment used a \textbf{Latin square} to counterbalance item--condition
pairings across participants.
With $k = 4$ conditions and 8 items, a balanced assignment is achieved:
each of the 4 lists assigns item $i$ ($i = 0,\ldots,7$) to condition
$(i + l)\bmod 4$, where $l \in \{0,1,2,3\}$ is the list number.
The resulting assignment matrix (Table~\ref{tab:ls}) guarantees that
\begin{itemize}
  \item each participant sees every item exactly once;
  \item each participant sees each condition exactly twice (two items per
        condition), ensuring balanced coverage of the 2$\times$2 design;
  \item across all 4 lists, each item appears equally often in each
        condition (twice per condition over the full participant sample).
\end{itemize}

\begin{table}[ht]
\centering
\caption{Latin square assignment of conditions to items across 4 lists.
  Cell entries are condition indices (0--3) as defined in
  Table~\ref{tab:design}.
  Rows = lists ($l$); columns = items ($i$).}
\label{tab:ls}
\small
\begin{tabular}{ccccccccc}
\toprule
\textbf{List} &
\textbf{Item 1} & \textbf{Item 2} & \textbf{Item 3} & \textbf{Item 4} &
\textbf{Item 5} & \textbf{Item 6} & \textbf{Item 7} & \textbf{Item 8} \\
\midrule
0 & 0 & 1 & 2 & 3 & 0 & 1 & 2 & 3 \\
1 & 1 & 2 & 3 & 0 & 1 & 2 & 3 & 0 \\
2 & 2 & 3 & 0 & 1 & 2 & 3 & 0 & 1 \\
3 & 3 & 0 & 1 & 2 & 3 & 0 & 1 & 2 \\
\bottomrule
\end{tabular}
\end{table}

\noindent List assignment is determined automatically from the URL query
parameter \texttt{?list=N} (0--3); if absent, a random list is selected.
Each participant is thus within-subjects across items and sees all four
conditions (twice each). Condition effects are estimated from within-participant
variation across the 2$\times$2 design, controlling for by-participant
random intercepts.

In addition to the 8 critical trials, each participant completed 8
\textbf{filler trials} drawn from separate topic domains (train travel,
reading light, cooking, coffee, vacation, hydration, fresh air, noise).
Filler items used the same trial format but were excluded from the
main analysis; they serve to dilute the critical items and reduce
potential demand characteristics.

\subsubsection{Dependent variable}

The dependent variable was the marker chosen from a drop-down menu
displaying all five options in fixed order (as listed in
Section~\ref{sub:markers}).
No response time threshold was imposed, but response times were logged.
Recorded per trial: \texttt{trial\_id}, \texttt{topic}, \texttt{is\_filler},
\texttt{condition\_index}, \texttt{pc\_prag}, \texttt{g},
\texttt{selected\_marker}, \texttt{rt} (ms); and once per participant:
\texttt{list\_num}.
The continuous \texttt{pc\_prop\_rating} covariate (topic-level norming mean)
is joined to the trial data during analysis.

\subsection{Procedure}

Participants completed the experiment online via the
magpie platform \citep{kochetkova2024magpie} in a browser.
After reading instructions in German, participants were assigned to one
of 4 counterbalancing lists (via URL parameter or randomly).
They then proceeded through 16 trials---8 critical and 8 fillers---in a
fully randomised order, followed by a post-test questionnaire eliciting
native language, age, gender, education level, and optional comments.
Each trial displayed (a) a context paragraph, (b) a goal instruction,
and (c) the sentence frame with the marker drop-down embedded at the
Mittelfeld position.
The ``Weiter'' (Next) button was disabled until a marker had been
selected, ensuring a response on every trial.

\end{document}
.
%
% Compile stand-alone:
%   pdflatex production_exp.tex
%
% Embed in main.tex:
%   % production_exp.tex
% Stand-alone description of the production experiment.
% Can be embedded in main.tex via \input{production_exp}.
%
% Compile stand-alone:
%   pdflatex production_exp.tex
%
% Embed in main.tex:
%   \input{production_exp}   % (inside a \section{Experiment} environment)

\documentclass[12pt]{article}
\usepackage{CSPArticleStyle}
\usepackage{csquotes}
\usepackage[margin=1in]{geometry}
\usepackage{booktabs}
\usepackage{parskip}
\usepackage{amsmath}
\usepackage{times}
\usepackage[T1]{fontenc}
\usepackage[utf8]{inputenc}
\usepackage{xcolor}
\usepackage{tabularray}   % robust tables; fall back to tabular if unavailable
\usepackage{gb4e}         % linguistic examples
\noautomath

\title{Production Experiment: Eliciting Consensus-Marker Choice\\
       Across Perceived-Controversy and Goal-Strength Conditions}
\author{}
\date{}

\begin{document}

%% ============================================================
%% When embedding in main.tex, delete from \documentclass to
%% \begin{document} and delete \end{document} at the bottom.
%% ============================================================

\section{Experiment: Consensus-Marker Production}

\subsection{Overview}

The production experiment empirically estimates the decision thresholds
that govern the use of five German consensus markers by placing
participants in the speaker role.
Participants read short vignettes that independently varied
\emph{pragmatic perceived controversy} (\textsc{pc\textsubscript{prag}})
and the speaker's \emph{persuasive-goal strength} ($g$), then selected
which consensus marker they would most naturally use in the described
situation.
\emph{Propositional perceived controversy} (\textsc{pc\textsubscript{prop}})
is held constant within the experiment and operationalised as a continuous
topic-level covariate obtained from a separate norming study: for each
of the eight topic items, participants rated the degree to which the
proposition $p$ is considered established/accepted (0--100 scale);
the resulting means enter the regression analyses as a fixed covariate.
The resulting distribution of marker choices across the
$\textsc{pc}_\text{prag} \times g$ design, combined with the
continuous \textsc{pc\textsubscript{prop}} covariate, provides the
empirical basis for estimating the ordinal thresholds
$\mu_1 < \cdots < \mu_4$ of the noisy-threshold RSA model.

\subsection{Participants}

Participants will be recruited via Prolific and will be required to be
native speakers of German (language-at-home filter).
Data collection is ongoing; a target of $N = 80$ participants is planned
($n = 20$ per list across 4 lists), each completing 16 trials (8 critical + 8 filler).

\subsection{Materials}

\subsubsection{Markers}

The five consensus markers under investigation are, from weakest to
strongest presumed common-ground commitment:

\begin{enumerate}
  \item \textit{sofern ich wei\ss{}} (`as far as I know')
  \item \textit{wie du ja wei\ss t} (`as you (prt) know')
  \item \textit{wie wir wissen} (`as we know')
  \item \textit{ja} (modal particle; roughly `as you know / of course')
  \item \textit{bekanntlich} (`as is well known / notoriously')
\end{enumerate}

These markers were embedded in the sentential \emph{Mittelfeld}
(middle field), immediately after the finite verb, as illustrated in
example~(\ref{ex:frame}):

\begin{exe}
  \ex \label{ex:frame}
  Digitale Kompetenzen sind \textbf{[MARKER]}
  in der heutigen Arbeitswelt unverzichtbar.\\
  `Digital skills are \textbf{[MARKER]} indispensable in today's
  working world.'
\end{exe}

This syntactic position is licit for all five markers: modal particles
such as \textit{ja} canonically occupy the left edge of the Mittelfeld
\citep{thurmair1989modalpartikeln}, while the adverb \textit{bekanntlich}
and parenthetical clauses (\textit{wie wir wissen}, etc.) can appear
there as well (see \citealt{jacobs_semantics_1991} for discussion).
Each target sentence was followed by a consequent clause with the
recommended action $q$: \textit{Deshalb solltest du [q].}
(`Therefore, you should [q].')

\subsubsection{Topic Items}

Eight topic domains were used, each associated with a fixed proposition
$p$ (sentence frame) and consequent action $q$ (Table~\ref{tab:items}).
The first four items were used in earlier pilot work; items 5--8 were
added to reach the minimum of eight required for a complete Latin square.

\begin{table}[ht]
\centering
\caption{Topic items: sentence frame for $p$ (with marker slot [M]) and consequent $q$.}
\label{tab:items}
\small
\begin{tabular}{clll}
\toprule
\textbf{\#} & \textbf{Topic} & \textbf{Sentence frame} & \textbf{Consequent $q$} \\
\midrule
1 & Klimawandel &
  \textit{Der Klimawandel wird [M]} &
  auf Flugreisen verzichten \\
  & & \textit{durch menschliche Aktivit\"aten verursacht.} & \\[3pt]
2 & Ern\"ahrung &
  \textit{Eine \"uberwiegend pflanzliche Ern\"ahrung} &
  mehr pflanzliche Lebens- \\
  & & \textit{f\"ordert [M] die Gesundheit.} &
  mittel einbauen \\[3pt]
3 & Stadtverkehr &
  \textit{Ein ausgebautes Nahverkehrsnetz} &
  h\"aufiger auf \"offentliche \\
  & & \textit{entlastet [M] den Stadtverkehr deutlich.} &
  Verkehrsmittel umsteigen \\[3pt]
4 & Digitalisierung &
  \textit{Digitale Kompetenzen sind [M] in der} &
  regelm\"a\ss{}ig Weiterbildungs- \\
  & & \textit{heutigen Arbeitswelt unverzichtbar.} &
  angebote nutzen \\[3pt]
5 & Schlaf &
  \textit{Ausreichend Schlaf ist [M] entscheidend} &
  auf feste Schlafzeiten \\
  & & \textit{f\"ur die kognitive Leistungsf\"ahigkeit.} &
  achten \\[3pt]
6 & Plastik &
  \textit{Einwegplastik schadet [M] den} &
  konsequent auf Einweg- \\
  & & \textit{Meeres\"okosystemen erheblich.} &
  plastik verzichten \\[3pt]
7 & Sport &
  \textit{Regelm\"a\ss{}ige k\"orperliche Bewegung} &
  regelm\"a\ss{}ig Sport in \\
  & & \textit{senkt [M] das Risiko f\"ur} &
  den Alltag integrieren \\
  & & \textit{Herz-Kreislauf-Erkrankungen deutlich.} & \\[3pt]
8 & Lokalkauf &
  \textit{Das Kaufen bei lokalen H\"andlern} &
  \"ofter bei lokalen \\
  & & \textit{st\"arkt [M] die regionale Wirtschaft} &
  Gesch\"aften einkaufen \\
  & & \textit{nachhaltig.} & \\
\bottomrule
\end{tabular}
\end{table}

\noindent [M] denotes the marker slot (dropdown).

\subsubsection{Vignettes}

\paragraph{Propositional perceived controversy (\textsc{pc\textsubscript{prop}}).}
\textsc{pc\textsubscript{prop}} was \emph{not} manipulated within the
production experiment.
Instead, it was operationalised as a continuous, topic-level covariate
obtained from a separate norming study conducted prior to the production
experiment.
In that norming study, a different sample of native German speakers rated,
for each topic proposition $p$, the degree to which they considered $p$
accepted/established in (German) public discourse, on a 0--100 scale
($0 = $ `completely contested', $100 = $ `universally accepted').
The resulting per-topic means (e.g., \textit{Digitalisierung}: 85/100;
\textit{Lokalkauf}: 58/100) enter the regression analyses as a
continuous, centred predictor \textsc{pc\textsubscript{prop}}$_c$
(mean-centred, SD $\approx$ 1).
Context paragraphs in the production experiment therefore present
the proposition $p$ without any editorial evaluation, so that
participants' experience of propositional controversy is shaped
solely by their own prior knowledge and the norming-derived covariate
captures this variation in analysis.

\paragraph{Pragmatic perceived controversy (\textsc{pc\textsubscript{prag}}).}
\textsc{pc\textsubscript{prag}} was manipulated between instances of each item
by varying whether the speaker's social circle \emph{largely agreed} (low)
or \emph{held divergent views} (high) about $p$.
The two context versions for the \textit{Digitalisierung} item are:

\begin{itemize}
  \item $\textsc{pc}_\text{prag} = \text{low}$: \textit{Du sprichst mit einer
        Kollegin {\"u}ber Berufsausbildung und Weiterbildung.
        In eurem Team teilen die meisten die Ansicht,
        dass digitale Kompetenzen in der heutigen Arbeitswelt unverzichtbar sind.}
  \item $\textsc{pc}_\text{prag} = \text{high}$: \textit{Du sprichst mit einer
        Kollegin {\"u}ber Berufsausbildung und Weiterbildung.
        In eurem Team gehen die Meinungen dazu stark auseinander.}
\end{itemize}

\paragraph{Speaker goal strength ($g$).}
Goal strength was manipulated via a standardised goal instruction
appended after the context paragraph:

\begin{itemize}
  \item $g_\text{low}$: \textit{Du m\"ochtest das nur kurz anmerken.}
        (`You just want to make a brief note of this.')
  \item $g_\text{high}$: \textit{Es ist dir sehr wichtig, dass dein
        Gespr\"achspartner das akzeptiert.}
        (`It is very important to you that your interlocutor accepts this.')
\end{itemize}

\paragraph{Complete sample item (all four conditions).}
Examples (\ref{ex:cond1})--(\ref{ex:cond4}) show the full vignette for
the topic \textit{Digitalisierung}
(norming $\textsc{pc}_\text{prop}$ rating: 85/100;
sentence frame: \textit{Digitale Kompetenzen sind [M] in der heutigen
Arbeitswelt unverzichtbar.
Deshalb solltest du regelm\"a\ss{}ig digitale Weiterbildungsangebote nutzen.})
across all four conditions of the $2\times2$ design.

\begin{exe}
  \ex \label{ex:cond1}
  \textbf{Condition 0: $\textsc{pc}_\text{prag} = \text{low}$, $g = \text{low}$}\\
  \textit{Du sprichst mit einer Kollegin {\"u}ber Berufsausbildung und Weiterbildung.
  In eurem Team teilen die meisten die Ansicht,
  dass digitale Kompetenzen in der heutigen Arbeitswelt unverzichtbar sind.}\\
  \textit{Du m\"ochtest das nur kurz anmerken.}\\
  Digitale Kompetenzen sind \underline{\hspace{2.2cm}}
  in der heutigen Arbeitswelt unverzichtbar.

  \ex \label{ex:cond2}
  \textbf{Condition 1: $\textsc{pc}_\text{prag} = \text{low}$, $g = \text{high}$}\\
  \textit{Du sprichst mit einer Kollegin {\"u}ber Berufsausbildung und Weiterbildung.
  In eurem Team teilen die meisten die Ansicht,
  dass digitale Kompetenzen in der heutigen Arbeitswelt unverzichtbar sind.}\\
  \textit{Es ist dir sehr wichtig, dass dein Gespr\"achspartner das akzeptiert.}\\
  Digitale Kompetenzen sind \underline{\hspace{2.2cm}}
  in der heutigen Arbeitswelt unverzichtbar.

  \ex \label{ex:cond3}
  \textbf{Condition 2: $\textsc{pc}_\text{prag} = \text{high}$, $g = \text{low}$}\\
  \textit{Du sprichst mit einer Kollegin {\"u}ber Berufsausbildung und Weiterbildung.
  In eurem Team gehen die Meinungen dazu stark auseinander.}\\
  \textit{Du m\"ochtest das nur kurz anmerken.}\\
  Digitale Kompetenzen sind \underline{\hspace{2.2cm}}
  in der heutigen Arbeitswelt unverzichtbar.

  \ex \label{ex:cond4}
  \textbf{Condition 3: $\textsc{pc}_\text{prag} = \text{high}$, $g = \text{high}$}\\
  \textit{Du sprichst mit einer Kollegin {\"u}ber Berufsausbildung und Weiterbildung.
  In eurem Team gehen die Meinungen dazu stark auseinander.}\\
  \textit{Es ist dir sehr wichtig, dass dein Gespr\"achspartner das akzeptiert.}\\
  Digitale Kompetenzen sind \underline{\hspace{2.2cm}}
  in der heutigen Arbeitswelt unverzichtbar.
\end{exe}

\noindent In all four vignettes the sentence frame and consequent clause
(\textit{Deshalb solltest du regelm\"a\ss{}ig digitale
Weiterbildungsangebote nutzen.}) are identical; only the context
paragraph and goal instruction vary across conditions.
The marker slot [M] / \underline{\hspace{2.2cm}} was presented as a
drop-down menu in the actual experiment.

\subsection{Design}

\subsubsection{Factorial structure}

The experiment employed a \textbf{2 (pc\textsubscript{prag}: low vs.\ high)
$\times$ 2 ($g$: low vs.\ high) fully crossed design},
yielding 4 conditions.
\textsc{pc\textsubscript{prop}} is not a within-experiment factor; it enters
as a continuous topic-level covariate from the norming study.
With 8 topic items and a Latin square rotation across 4 lists, each
participant saw exactly \textbf{8 critical trials}---two items per condition---plus
8 filler trials, for \textbf{16 trials} in total
(Table~\ref{tab:design}).

\begin{table}[ht]
\centering
\caption{%
  The $2\times2$ design (4 conditions).
  \textsc{pc\textsubscript{prop}} is a continuous norming covariate, not a
  within-experiment factor.
  In each list, the 8 items are assigned to conditions via the Latin square
  rotation described in Section~\ref{subsub:counterbalancing}.}
\label{tab:design}
\small
\begin{tabular}{ccl}
\toprule
\textbf{pc\textsubscript{prag}} &
\textbf{$g$} &
\textbf{Condition label} \\
\midrule
low  & low  & (0) community agrees, casual \\
low  & high & (1) community agrees, pressing \\
high & low  & (2) community divided, casual \\
high & high & (3) community divided, pressing \\
\bottomrule
\end{tabular}
\end{table}

\subsubsection{Counterbalancing}
\label{subsub:counterbalancing}

The experiment used a \textbf{Latin square} to counterbalance item--condition
pairings across participants.
With $k = 4$ conditions and 8 items, a balanced assignment is achieved:
each of the 4 lists assigns item $i$ ($i = 0,\ldots,7$) to condition
$(i + l)\bmod 4$, where $l \in \{0,1,2,3\}$ is the list number.
The resulting assignment matrix (Table~\ref{tab:ls}) guarantees that
\begin{itemize}
  \item each participant sees every item exactly once;
  \item each participant sees each condition exactly twice (two items per
        condition), ensuring balanced coverage of the 2$\times$2 design;
  \item across all 4 lists, each item appears equally often in each
        condition (twice per condition over the full participant sample).
\end{itemize}

\begin{table}[ht]
\centering
\caption{Latin square assignment of conditions to items across 4 lists.
  Cell entries are condition indices (0--3) as defined in
  Table~\ref{tab:design}.
  Rows = lists ($l$); columns = items ($i$).}
\label{tab:ls}
\small
\begin{tabular}{ccccccccc}
\toprule
\textbf{List} &
\textbf{Item 1} & \textbf{Item 2} & \textbf{Item 3} & \textbf{Item 4} &
\textbf{Item 5} & \textbf{Item 6} & \textbf{Item 7} & \textbf{Item 8} \\
\midrule
0 & 0 & 1 & 2 & 3 & 0 & 1 & 2 & 3 \\
1 & 1 & 2 & 3 & 0 & 1 & 2 & 3 & 0 \\
2 & 2 & 3 & 0 & 1 & 2 & 3 & 0 & 1 \\
3 & 3 & 0 & 1 & 2 & 3 & 0 & 1 & 2 \\
\bottomrule
\end{tabular}
\end{table}

\noindent List assignment is determined automatically from the URL query
parameter \texttt{?list=N} (0--3); if absent, a random list is selected.
Each participant is thus within-subjects across items and sees all four
conditions (twice each). Condition effects are estimated from within-participant
variation across the 2$\times$2 design, controlling for by-participant
random intercepts.

In addition to the 8 critical trials, each participant completed 8
\textbf{filler trials} drawn from separate topic domains (train travel,
reading light, cooking, coffee, vacation, hydration, fresh air, noise).
Filler items used the same trial format but were excluded from the
main analysis; they serve to dilute the critical items and reduce
potential demand characteristics.

\subsubsection{Dependent variable}

The dependent variable was the marker chosen from a drop-down menu
displaying all five options in fixed order (as listed in
Section~\ref{sub:markers}).
No response time threshold was imposed, but response times were logged.
Recorded per trial: \texttt{trial\_id}, \texttt{topic}, \texttt{is\_filler},
\texttt{condition\_index}, \texttt{pc\_prag}, \texttt{g},
\texttt{selected\_marker}, \texttt{rt} (ms); and once per participant:
\texttt{list\_num}.
The continuous \texttt{pc\_prop\_rating} covariate (topic-level norming mean)
is joined to the trial data during analysis.

\subsection{Procedure}

Participants completed the experiment online via the
magpie platform \citep{kochetkova2024magpie} in a browser.
After reading instructions in German, participants were assigned to one
of 4 counterbalancing lists (via URL parameter or randomly).
They then proceeded through 16 trials---8 critical and 8 fillers---in a
fully randomised order, followed by a post-test questionnaire eliciting
native language, age, gender, education level, and optional comments.
Each trial displayed (a) a context paragraph, (b) a goal instruction,
and (c) the sentence frame with the marker drop-down embedded at the
Mittelfeld position.
The ``Weiter'' (Next) button was disabled until a marker had been
selected, ensuring a response on every trial.

\end{document}
   % (inside a \section{Experiment} environment)

\documentclass[12pt]{article}
\usepackage{CSPArticleStyle}
\usepackage{csquotes}
\usepackage[margin=1in]{geometry}
\usepackage{booktabs}
\usepackage{parskip}
\usepackage{amsmath}
\usepackage{times}
\usepackage[T1]{fontenc}
\usepackage[utf8]{inputenc}
\usepackage{xcolor}
\usepackage{tabularray}   % robust tables; fall back to tabular if unavailable
\usepackage{gb4e}         % linguistic examples
\noautomath

\title{Production Experiment: Eliciting Consensus-Marker Choice\\
       Across Perceived-Controversy and Goal-Strength Conditions}
\author{}
\date{}

\begin{document}

%% ============================================================
%% When embedding in main.tex, delete from \documentclass to
%% \begin{document} and delete \end{document} at the bottom.
%% ============================================================

\section{Experiment: Consensus-Marker Production}

\subsection{Overview}

The production experiment empirically estimates the decision thresholds
that govern the use of five German consensus markers by placing
participants in the speaker role.
Participants read short vignettes that independently varied
\emph{pragmatic perceived controversy} (\textsc{pc\textsubscript{prag}})
and the speaker's \emph{persuasive-goal strength} ($g$), then selected
which consensus marker they would most naturally use in the described
situation.
\emph{Propositional perceived controversy} (\textsc{pc\textsubscript{prop}})
is held constant within the experiment and operationalised as a continuous
topic-level covariate obtained from a separate norming study: for each
of the eight topic items, participants rated the degree to which the
proposition $p$ is considered established/accepted (0--100 scale);
the resulting means enter the regression analyses as a fixed covariate.
The resulting distribution of marker choices across the
$\textsc{pc}_\text{prag} \times g$ design, combined with the
continuous \textsc{pc\textsubscript{prop}} covariate, provides the
empirical basis for estimating the ordinal thresholds
$\mu_1 < \cdots < \mu_4$ of the noisy-threshold RSA model.

\subsection{Participants}

Participants will be recruited via Prolific and will be required to be
native speakers of German (language-at-home filter).
Data collection is ongoing; a target of $N = 80$ participants is planned
($n = 20$ per list across 4 lists), each completing 16 trials (8 critical + 8 filler).

\subsection{Materials}

\subsubsection{Markers}

The five consensus markers under investigation are, from weakest to
strongest presumed common-ground commitment:

\begin{enumerate}
  \item \textit{sofern ich wei\ss{}} (`as far as I know')
  \item \textit{wie du ja wei\ss t} (`as you (prt) know')
  \item \textit{wie wir wissen} (`as we know')
  \item \textit{ja} (modal particle; roughly `as you know / of course')
  \item \textit{bekanntlich} (`as is well known / notoriously')
\end{enumerate}

These markers were embedded in the sentential \emph{Mittelfeld}
(middle field), immediately after the finite verb, as illustrated in
example~(\ref{ex:frame}):

\begin{exe}
  \ex \label{ex:frame}
  Digitale Kompetenzen sind \textbf{[MARKER]}
  in der heutigen Arbeitswelt unverzichtbar.\\
  `Digital skills are \textbf{[MARKER]} indispensable in today's
  working world.'
\end{exe}

This syntactic position is licit for all five markers: modal particles
such as \textit{ja} canonically occupy the left edge of the Mittelfeld
\citep{thurmair1989modalpartikeln}, while the adverb \textit{bekanntlich}
and parenthetical clauses (\textit{wie wir wissen}, etc.) can appear
there as well (see \citealt{jacobs_semantics_1991} for discussion).
Each target sentence was followed by a consequent clause with the
recommended action $q$: \textit{Deshalb solltest du [q].}
(`Therefore, you should [q].')

\subsubsection{Topic Items}

Eight topic domains were used, each associated with a fixed proposition
$p$ (sentence frame) and consequent action $q$ (Table~\ref{tab:items}).
The first four items were used in earlier pilot work; items 5--8 were
added to reach the minimum of eight required for a complete Latin square.

\begin{table}[ht]
\centering
\caption{Topic items: sentence frame for $p$ (with marker slot [M]) and consequent $q$.}
\label{tab:items}
\small
\begin{tabular}{clll}
\toprule
\textbf{\#} & \textbf{Topic} & \textbf{Sentence frame} & \textbf{Consequent $q$} \\
\midrule
1 & Klimawandel &
  \textit{Der Klimawandel wird [M]} &
  auf Flugreisen verzichten \\
  & & \textit{durch menschliche Aktivit\"aten verursacht.} & \\[3pt]
2 & Ern\"ahrung &
  \textit{Eine \"uberwiegend pflanzliche Ern\"ahrung} &
  mehr pflanzliche Lebens- \\
  & & \textit{f\"ordert [M] die Gesundheit.} &
  mittel einbauen \\[3pt]
3 & Stadtverkehr &
  \textit{Ein ausgebautes Nahverkehrsnetz} &
  h\"aufiger auf \"offentliche \\
  & & \textit{entlastet [M] den Stadtverkehr deutlich.} &
  Verkehrsmittel umsteigen \\[3pt]
4 & Digitalisierung &
  \textit{Digitale Kompetenzen sind [M] in der} &
  regelm\"a\ss{}ig Weiterbildungs- \\
  & & \textit{heutigen Arbeitswelt unverzichtbar.} &
  angebote nutzen \\[3pt]
5 & Schlaf &
  \textit{Ausreichend Schlaf ist [M] entscheidend} &
  auf feste Schlafzeiten \\
  & & \textit{f\"ur die kognitive Leistungsf\"ahigkeit.} &
  achten \\[3pt]
6 & Plastik &
  \textit{Einwegplastik schadet [M] den} &
  konsequent auf Einweg- \\
  & & \textit{Meeres\"okosystemen erheblich.} &
  plastik verzichten \\[3pt]
7 & Sport &
  \textit{Regelm\"a\ss{}ige k\"orperliche Bewegung} &
  regelm\"a\ss{}ig Sport in \\
  & & \textit{senkt [M] das Risiko f\"ur} &
  den Alltag integrieren \\
  & & \textit{Herz-Kreislauf-Erkrankungen deutlich.} & \\[3pt]
8 & Lokalkauf &
  \textit{Das Kaufen bei lokalen H\"andlern} &
  \"ofter bei lokalen \\
  & & \textit{st\"arkt [M] die regionale Wirtschaft} &
  Gesch\"aften einkaufen \\
  & & \textit{nachhaltig.} & \\
\bottomrule
\end{tabular}
\end{table}

\noindent [M] denotes the marker slot (dropdown).

\subsubsection{Vignettes}

\paragraph{Propositional perceived controversy (\textsc{pc\textsubscript{prop}}).}
\textsc{pc\textsubscript{prop}} was \emph{not} manipulated within the
production experiment.
Instead, it was operationalised as a continuous, topic-level covariate
obtained from a separate norming study conducted prior to the production
experiment.
In that norming study, a different sample of native German speakers rated,
for each topic proposition $p$, the degree to which they considered $p$
accepted/established in (German) public discourse, on a 0--100 scale
($0 = $ `completely contested', $100 = $ `universally accepted').
The resulting per-topic means (e.g., \textit{Digitalisierung}: 85/100;
\textit{Lokalkauf}: 58/100) enter the regression analyses as a
continuous, centred predictor \textsc{pc\textsubscript{prop}}$_c$
(mean-centred, SD $\approx$ 1).
Context paragraphs in the production experiment therefore present
the proposition $p$ without any editorial evaluation, so that
participants' experience of propositional controversy is shaped
solely by their own prior knowledge and the norming-derived covariate
captures this variation in analysis.

\paragraph{Pragmatic perceived controversy (\textsc{pc\textsubscript{prag}}).}
\textsc{pc\textsubscript{prag}} was manipulated between instances of each item
by varying whether the speaker's social circle \emph{largely agreed} (low)
or \emph{held divergent views} (high) about $p$.
The two context versions for the \textit{Digitalisierung} item are:

\begin{itemize}
  \item $\textsc{pc}_\text{prag} = \text{low}$: \textit{Du sprichst mit einer
        Kollegin {\"u}ber Berufsausbildung und Weiterbildung.
        In eurem Team teilen die meisten die Ansicht,
        dass digitale Kompetenzen in der heutigen Arbeitswelt unverzichtbar sind.}
  \item $\textsc{pc}_\text{prag} = \text{high}$: \textit{Du sprichst mit einer
        Kollegin {\"u}ber Berufsausbildung und Weiterbildung.
        In eurem Team gehen die Meinungen dazu stark auseinander.}
\end{itemize}

\paragraph{Speaker goal strength ($g$).}
Goal strength was manipulated via a standardised goal instruction
appended after the context paragraph:

\begin{itemize}
  \item $g_\text{low}$: \textit{Du m\"ochtest das nur kurz anmerken.}
        (`You just want to make a brief note of this.')
  \item $g_\text{high}$: \textit{Es ist dir sehr wichtig, dass dein
        Gespr\"achspartner das akzeptiert.}
        (`It is very important to you that your interlocutor accepts this.')
\end{itemize}

\paragraph{Complete sample item (all four conditions).}
Examples (\ref{ex:cond1})--(\ref{ex:cond4}) show the full vignette for
the topic \textit{Digitalisierung}
(norming $\textsc{pc}_\text{prop}$ rating: 85/100;
sentence frame: \textit{Digitale Kompetenzen sind [M] in der heutigen
Arbeitswelt unverzichtbar.
Deshalb solltest du regelm\"a\ss{}ig digitale Weiterbildungsangebote nutzen.})
across all four conditions of the $2\times2$ design.

\begin{exe}
  \ex \label{ex:cond1}
  \textbf{Condition 0: $\textsc{pc}_\text{prag} = \text{low}$, $g = \text{low}$}\\
  \textit{Du sprichst mit einer Kollegin {\"u}ber Berufsausbildung und Weiterbildung.
  In eurem Team teilen die meisten die Ansicht,
  dass digitale Kompetenzen in der heutigen Arbeitswelt unverzichtbar sind.}\\
  \textit{Du m\"ochtest das nur kurz anmerken.}\\
  Digitale Kompetenzen sind \underline{\hspace{2.2cm}}
  in der heutigen Arbeitswelt unverzichtbar.

  \ex \label{ex:cond2}
  \textbf{Condition 1: $\textsc{pc}_\text{prag} = \text{low}$, $g = \text{high}$}\\
  \textit{Du sprichst mit einer Kollegin {\"u}ber Berufsausbildung und Weiterbildung.
  In eurem Team teilen die meisten die Ansicht,
  dass digitale Kompetenzen in der heutigen Arbeitswelt unverzichtbar sind.}\\
  \textit{Es ist dir sehr wichtig, dass dein Gespr\"achspartner das akzeptiert.}\\
  Digitale Kompetenzen sind \underline{\hspace{2.2cm}}
  in der heutigen Arbeitswelt unverzichtbar.

  \ex \label{ex:cond3}
  \textbf{Condition 2: $\textsc{pc}_\text{prag} = \text{high}$, $g = \text{low}$}\\
  \textit{Du sprichst mit einer Kollegin {\"u}ber Berufsausbildung und Weiterbildung.
  In eurem Team gehen die Meinungen dazu stark auseinander.}\\
  \textit{Du m\"ochtest das nur kurz anmerken.}\\
  Digitale Kompetenzen sind \underline{\hspace{2.2cm}}
  in der heutigen Arbeitswelt unverzichtbar.

  \ex \label{ex:cond4}
  \textbf{Condition 3: $\textsc{pc}_\text{prag} = \text{high}$, $g = \text{high}$}\\
  \textit{Du sprichst mit einer Kollegin {\"u}ber Berufsausbildung und Weiterbildung.
  In eurem Team gehen die Meinungen dazu stark auseinander.}\\
  \textit{Es ist dir sehr wichtig, dass dein Gespr\"achspartner das akzeptiert.}\\
  Digitale Kompetenzen sind \underline{\hspace{2.2cm}}
  in der heutigen Arbeitswelt unverzichtbar.
\end{exe}

\noindent In all four vignettes the sentence frame and consequent clause
(\textit{Deshalb solltest du regelm\"a\ss{}ig digitale
Weiterbildungsangebote nutzen.}) are identical; only the context
paragraph and goal instruction vary across conditions.
The marker slot [M] / \underline{\hspace{2.2cm}} was presented as a
drop-down menu in the actual experiment.

\subsection{Design}

\subsubsection{Factorial structure}

The experiment employed a \textbf{2 (pc\textsubscript{prag}: low vs.\ high)
$\times$ 2 ($g$: low vs.\ high) fully crossed design},
yielding 4 conditions.
\textsc{pc\textsubscript{prop}} is not a within-experiment factor; it enters
as a continuous topic-level covariate from the norming study.
With 8 topic items and a Latin square rotation across 4 lists, each
participant saw exactly \textbf{8 critical trials}---two items per condition---plus
8 filler trials, for \textbf{16 trials} in total
(Table~\ref{tab:design}).

\begin{table}[ht]
\centering
\caption{%
  The $2\times2$ design (4 conditions).
  \textsc{pc\textsubscript{prop}} is a continuous norming covariate, not a
  within-experiment factor.
  In each list, the 8 items are assigned to conditions via the Latin square
  rotation described in Section~\ref{subsub:counterbalancing}.}
\label{tab:design}
\small
\begin{tabular}{ccl}
\toprule
\textbf{pc\textsubscript{prag}} &
\textbf{$g$} &
\textbf{Condition label} \\
\midrule
low  & low  & (0) community agrees, casual \\
low  & high & (1) community agrees, pressing \\
high & low  & (2) community divided, casual \\
high & high & (3) community divided, pressing \\
\bottomrule
\end{tabular}
\end{table}

\subsubsection{Counterbalancing}
\label{subsub:counterbalancing}

The experiment used a \textbf{Latin square} to counterbalance item--condition
pairings across participants.
With $k = 4$ conditions and 8 items, a balanced assignment is achieved:
each of the 4 lists assigns item $i$ ($i = 0,\ldots,7$) to condition
$(i + l)\bmod 4$, where $l \in \{0,1,2,3\}$ is the list number.
The resulting assignment matrix (Table~\ref{tab:ls}) guarantees that
\begin{itemize}
  \item each participant sees every item exactly once;
  \item each participant sees each condition exactly twice (two items per
        condition), ensuring balanced coverage of the 2$\times$2 design;
  \item across all 4 lists, each item appears equally often in each
        condition (twice per condition over the full participant sample).
\end{itemize}

\begin{table}[ht]
\centering
\caption{Latin square assignment of conditions to items across 4 lists.
  Cell entries are condition indices (0--3) as defined in
  Table~\ref{tab:design}.
  Rows = lists ($l$); columns = items ($i$).}
\label{tab:ls}
\small
\begin{tabular}{ccccccccc}
\toprule
\textbf{List} &
\textbf{Item 1} & \textbf{Item 2} & \textbf{Item 3} & \textbf{Item 4} &
\textbf{Item 5} & \textbf{Item 6} & \textbf{Item 7} & \textbf{Item 8} \\
\midrule
0 & 0 & 1 & 2 & 3 & 0 & 1 & 2 & 3 \\
1 & 1 & 2 & 3 & 0 & 1 & 2 & 3 & 0 \\
2 & 2 & 3 & 0 & 1 & 2 & 3 & 0 & 1 \\
3 & 3 & 0 & 1 & 2 & 3 & 0 & 1 & 2 \\
\bottomrule
\end{tabular}
\end{table}

\noindent List assignment is determined automatically from the URL query
parameter \texttt{?list=N} (0--3); if absent, a random list is selected.
Each participant is thus within-subjects across items and sees all four
conditions (twice each). Condition effects are estimated from within-participant
variation across the 2$\times$2 design, controlling for by-participant
random intercepts.

In addition to the 8 critical trials, each participant completed 8
\textbf{filler trials} drawn from separate topic domains (train travel,
reading light, cooking, coffee, vacation, hydration, fresh air, noise).
Filler items used the same trial format but were excluded from the
main analysis; they serve to dilute the critical items and reduce
potential demand characteristics.

\subsubsection{Dependent variable}

The dependent variable was the marker chosen from a drop-down menu
displaying all five options in fixed order (as listed in
Section~\ref{sub:markers}).
No response time threshold was imposed, but response times were logged.
Recorded per trial: \texttt{trial\_id}, \texttt{topic}, \texttt{is\_filler},
\texttt{condition\_index}, \texttt{pc\_prag}, \texttt{g},
\texttt{selected\_marker}, \texttt{rt} (ms); and once per participant:
\texttt{list\_num}.
The continuous \texttt{pc\_prop\_rating} covariate (topic-level norming mean)
is joined to the trial data during analysis.

\subsection{Procedure}

Participants completed the experiment online via the
magpie platform \citep{kochetkova2024magpie} in a browser.
After reading instructions in German, participants were assigned to one
of 4 counterbalancing lists (via URL parameter or randomly).
They then proceeded through 16 trials---8 critical and 8 fillers---in a
fully randomised order, followed by a post-test questionnaire eliciting
native language, age, gender, education level, and optional comments.
Each trial displayed (a) a context paragraph, (b) a goal instruction,
and (c) the sentence frame with the marker drop-down embedded at the
Mittelfeld position.
The ``Weiter'' (Next) button was disabled until a marker had been
selected, ensuring a response on every trial.

\end{document}
.
%
% Compile stand-alone:
%   pdflatex production_exp.tex
%
% Embed in main.tex:
%   % production_exp.tex
% Stand-alone description of the production experiment.
% Can be embedded in main.tex via % production_exp.tex
% Stand-alone description of the production experiment.
% Can be embedded in main.tex via \input{production_exp}.
%
% Compile stand-alone:
%   pdflatex production_exp.tex
%
% Embed in main.tex:
%   \input{production_exp}   % (inside a \section{Experiment} environment)

\documentclass[12pt]{article}
\usepackage{CSPArticleStyle}
\usepackage{csquotes}
\usepackage[margin=1in]{geometry}
\usepackage{booktabs}
\usepackage{parskip}
\usepackage{amsmath}
\usepackage{times}
\usepackage[T1]{fontenc}
\usepackage[utf8]{inputenc}
\usepackage{xcolor}
\usepackage{tabularray}   % robust tables; fall back to tabular if unavailable
\usepackage{gb4e}         % linguistic examples
\noautomath

\title{Production Experiment: Eliciting Consensus-Marker Choice\\
       Across Perceived-Controversy and Goal-Strength Conditions}
\author{}
\date{}

\begin{document}

%% ============================================================
%% When embedding in main.tex, delete from \documentclass to
%% \begin{document} and delete \end{document} at the bottom.
%% ============================================================

\section{Experiment: Consensus-Marker Production}

\subsection{Overview}

The production experiment empirically estimates the decision thresholds
that govern the use of five German consensus markers by placing
participants in the speaker role.
Participants read short vignettes that independently varied
\emph{pragmatic perceived controversy} (\textsc{pc\textsubscript{prag}})
and the speaker's \emph{persuasive-goal strength} ($g$), then selected
which consensus marker they would most naturally use in the described
situation.
\emph{Propositional perceived controversy} (\textsc{pc\textsubscript{prop}})
is held constant within the experiment and operationalised as a continuous
topic-level covariate obtained from a separate norming study: for each
of the eight topic items, participants rated the degree to which the
proposition $p$ is considered established/accepted (0--100 scale);
the resulting means enter the regression analyses as a fixed covariate.
The resulting distribution of marker choices across the
$\textsc{pc}_\text{prag} \times g$ design, combined with the
continuous \textsc{pc\textsubscript{prop}} covariate, provides the
empirical basis for estimating the ordinal thresholds
$\mu_1 < \cdots < \mu_4$ of the noisy-threshold RSA model.

\subsection{Participants}

Participants will be recruited via Prolific and will be required to be
native speakers of German (language-at-home filter).
Data collection is ongoing; a target of $N = 80$ participants is planned
($n = 20$ per list across 4 lists), each completing 16 trials (8 critical + 8 filler).

\subsection{Materials}

\subsubsection{Markers}

The five consensus markers under investigation are, from weakest to
strongest presumed common-ground commitment:

\begin{enumerate}
  \item \textit{sofern ich wei\ss{}} (`as far as I know')
  \item \textit{wie du ja wei\ss t} (`as you (prt) know')
  \item \textit{wie wir wissen} (`as we know')
  \item \textit{ja} (modal particle; roughly `as you know / of course')
  \item \textit{bekanntlich} (`as is well known / notoriously')
\end{enumerate}

These markers were embedded in the sentential \emph{Mittelfeld}
(middle field), immediately after the finite verb, as illustrated in
example~(\ref{ex:frame}):

\begin{exe}
  \ex \label{ex:frame}
  Digitale Kompetenzen sind \textbf{[MARKER]}
  in der heutigen Arbeitswelt unverzichtbar.\\
  `Digital skills are \textbf{[MARKER]} indispensable in today's
  working world.'
\end{exe}

This syntactic position is licit for all five markers: modal particles
such as \textit{ja} canonically occupy the left edge of the Mittelfeld
\citep{thurmair1989modalpartikeln}, while the adverb \textit{bekanntlich}
and parenthetical clauses (\textit{wie wir wissen}, etc.) can appear
there as well (see \citealt{jacobs_semantics_1991} for discussion).
Each target sentence was followed by a consequent clause with the
recommended action $q$: \textit{Deshalb solltest du [q].}
(`Therefore, you should [q].')

\subsubsection{Topic Items}

Eight topic domains were used, each associated with a fixed proposition
$p$ (sentence frame) and consequent action $q$ (Table~\ref{tab:items}).
The first four items were used in earlier pilot work; items 5--8 were
added to reach the minimum of eight required for a complete Latin square.

\begin{table}[ht]
\centering
\caption{Topic items: sentence frame for $p$ (with marker slot [M]) and consequent $q$.}
\label{tab:items}
\small
\begin{tabular}{clll}
\toprule
\textbf{\#} & \textbf{Topic} & \textbf{Sentence frame} & \textbf{Consequent $q$} \\
\midrule
1 & Klimawandel &
  \textit{Der Klimawandel wird [M]} &
  auf Flugreisen verzichten \\
  & & \textit{durch menschliche Aktivit\"aten verursacht.} & \\[3pt]
2 & Ern\"ahrung &
  \textit{Eine \"uberwiegend pflanzliche Ern\"ahrung} &
  mehr pflanzliche Lebens- \\
  & & \textit{f\"ordert [M] die Gesundheit.} &
  mittel einbauen \\[3pt]
3 & Stadtverkehr &
  \textit{Ein ausgebautes Nahverkehrsnetz} &
  h\"aufiger auf \"offentliche \\
  & & \textit{entlastet [M] den Stadtverkehr deutlich.} &
  Verkehrsmittel umsteigen \\[3pt]
4 & Digitalisierung &
  \textit{Digitale Kompetenzen sind [M] in der} &
  regelm\"a\ss{}ig Weiterbildungs- \\
  & & \textit{heutigen Arbeitswelt unverzichtbar.} &
  angebote nutzen \\[3pt]
5 & Schlaf &
  \textit{Ausreichend Schlaf ist [M] entscheidend} &
  auf feste Schlafzeiten \\
  & & \textit{f\"ur die kognitive Leistungsf\"ahigkeit.} &
  achten \\[3pt]
6 & Plastik &
  \textit{Einwegplastik schadet [M] den} &
  konsequent auf Einweg- \\
  & & \textit{Meeres\"okosystemen erheblich.} &
  plastik verzichten \\[3pt]
7 & Sport &
  \textit{Regelm\"a\ss{}ige k\"orperliche Bewegung} &
  regelm\"a\ss{}ig Sport in \\
  & & \textit{senkt [M] das Risiko f\"ur} &
  den Alltag integrieren \\
  & & \textit{Herz-Kreislauf-Erkrankungen deutlich.} & \\[3pt]
8 & Lokalkauf &
  \textit{Das Kaufen bei lokalen H\"andlern} &
  \"ofter bei lokalen \\
  & & \textit{st\"arkt [M] die regionale Wirtschaft} &
  Gesch\"aften einkaufen \\
  & & \textit{nachhaltig.} & \\
\bottomrule
\end{tabular}
\end{table}

\noindent [M] denotes the marker slot (dropdown).

\subsubsection{Vignettes}

\paragraph{Propositional perceived controversy (\textsc{pc\textsubscript{prop}}).}
\textsc{pc\textsubscript{prop}} was \emph{not} manipulated within the
production experiment.
Instead, it was operationalised as a continuous, topic-level covariate
obtained from a separate norming study conducted prior to the production
experiment.
In that norming study, a different sample of native German speakers rated,
for each topic proposition $p$, the degree to which they considered $p$
accepted/established in (German) public discourse, on a 0--100 scale
($0 = $ `completely contested', $100 = $ `universally accepted').
The resulting per-topic means (e.g., \textit{Digitalisierung}: 85/100;
\textit{Lokalkauf}: 58/100) enter the regression analyses as a
continuous, centred predictor \textsc{pc\textsubscript{prop}}$_c$
(mean-centred, SD $\approx$ 1).
Context paragraphs in the production experiment therefore present
the proposition $p$ without any editorial evaluation, so that
participants' experience of propositional controversy is shaped
solely by their own prior knowledge and the norming-derived covariate
captures this variation in analysis.

\paragraph{Pragmatic perceived controversy (\textsc{pc\textsubscript{prag}}).}
\textsc{pc\textsubscript{prag}} was manipulated between instances of each item
by varying whether the speaker's social circle \emph{largely agreed} (low)
or \emph{held divergent views} (high) about $p$.
The two context versions for the \textit{Digitalisierung} item are:

\begin{itemize}
  \item $\textsc{pc}_\text{prag} = \text{low}$: \textit{Du sprichst mit einer
        Kollegin {\"u}ber Berufsausbildung und Weiterbildung.
        In eurem Team teilen die meisten die Ansicht,
        dass digitale Kompetenzen in der heutigen Arbeitswelt unverzichtbar sind.}
  \item $\textsc{pc}_\text{prag} = \text{high}$: \textit{Du sprichst mit einer
        Kollegin {\"u}ber Berufsausbildung und Weiterbildung.
        In eurem Team gehen die Meinungen dazu stark auseinander.}
\end{itemize}

\paragraph{Speaker goal strength ($g$).}
Goal strength was manipulated via a standardised goal instruction
appended after the context paragraph:

\begin{itemize}
  \item $g_\text{low}$: \textit{Du m\"ochtest das nur kurz anmerken.}
        (`You just want to make a brief note of this.')
  \item $g_\text{high}$: \textit{Es ist dir sehr wichtig, dass dein
        Gespr\"achspartner das akzeptiert.}
        (`It is very important to you that your interlocutor accepts this.')
\end{itemize}

\paragraph{Complete sample item (all four conditions).}
Examples (\ref{ex:cond1})--(\ref{ex:cond4}) show the full vignette for
the topic \textit{Digitalisierung}
(norming $\textsc{pc}_\text{prop}$ rating: 85/100;
sentence frame: \textit{Digitale Kompetenzen sind [M] in der heutigen
Arbeitswelt unverzichtbar.
Deshalb solltest du regelm\"a\ss{}ig digitale Weiterbildungsangebote nutzen.})
across all four conditions of the $2\times2$ design.

\begin{exe}
  \ex \label{ex:cond1}
  \textbf{Condition 0: $\textsc{pc}_\text{prag} = \text{low}$, $g = \text{low}$}\\
  \textit{Du sprichst mit einer Kollegin {\"u}ber Berufsausbildung und Weiterbildung.
  In eurem Team teilen die meisten die Ansicht,
  dass digitale Kompetenzen in der heutigen Arbeitswelt unverzichtbar sind.}\\
  \textit{Du m\"ochtest das nur kurz anmerken.}\\
  Digitale Kompetenzen sind \underline{\hspace{2.2cm}}
  in der heutigen Arbeitswelt unverzichtbar.

  \ex \label{ex:cond2}
  \textbf{Condition 1: $\textsc{pc}_\text{prag} = \text{low}$, $g = \text{high}$}\\
  \textit{Du sprichst mit einer Kollegin {\"u}ber Berufsausbildung und Weiterbildung.
  In eurem Team teilen die meisten die Ansicht,
  dass digitale Kompetenzen in der heutigen Arbeitswelt unverzichtbar sind.}\\
  \textit{Es ist dir sehr wichtig, dass dein Gespr\"achspartner das akzeptiert.}\\
  Digitale Kompetenzen sind \underline{\hspace{2.2cm}}
  in der heutigen Arbeitswelt unverzichtbar.

  \ex \label{ex:cond3}
  \textbf{Condition 2: $\textsc{pc}_\text{prag} = \text{high}$, $g = \text{low}$}\\
  \textit{Du sprichst mit einer Kollegin {\"u}ber Berufsausbildung und Weiterbildung.
  In eurem Team gehen die Meinungen dazu stark auseinander.}\\
  \textit{Du m\"ochtest das nur kurz anmerken.}\\
  Digitale Kompetenzen sind \underline{\hspace{2.2cm}}
  in der heutigen Arbeitswelt unverzichtbar.

  \ex \label{ex:cond4}
  \textbf{Condition 3: $\textsc{pc}_\text{prag} = \text{high}$, $g = \text{high}$}\\
  \textit{Du sprichst mit einer Kollegin {\"u}ber Berufsausbildung und Weiterbildung.
  In eurem Team gehen die Meinungen dazu stark auseinander.}\\
  \textit{Es ist dir sehr wichtig, dass dein Gespr\"achspartner das akzeptiert.}\\
  Digitale Kompetenzen sind \underline{\hspace{2.2cm}}
  in der heutigen Arbeitswelt unverzichtbar.
\end{exe}

\noindent In all four vignettes the sentence frame and consequent clause
(\textit{Deshalb solltest du regelm\"a\ss{}ig digitale
Weiterbildungsangebote nutzen.}) are identical; only the context
paragraph and goal instruction vary across conditions.
The marker slot [M] / \underline{\hspace{2.2cm}} was presented as a
drop-down menu in the actual experiment.

\subsection{Design}

\subsubsection{Factorial structure}

The experiment employed a \textbf{2 (pc\textsubscript{prag}: low vs.\ high)
$\times$ 2 ($g$: low vs.\ high) fully crossed design},
yielding 4 conditions.
\textsc{pc\textsubscript{prop}} is not a within-experiment factor; it enters
as a continuous topic-level covariate from the norming study.
With 8 topic items and a Latin square rotation across 4 lists, each
participant saw exactly \textbf{8 critical trials}---two items per condition---plus
8 filler trials, for \textbf{16 trials} in total
(Table~\ref{tab:design}).

\begin{table}[ht]
\centering
\caption{%
  The $2\times2$ design (4 conditions).
  \textsc{pc\textsubscript{prop}} is a continuous norming covariate, not a
  within-experiment factor.
  In each list, the 8 items are assigned to conditions via the Latin square
  rotation described in Section~\ref{subsub:counterbalancing}.}
\label{tab:design}
\small
\begin{tabular}{ccl}
\toprule
\textbf{pc\textsubscript{prag}} &
\textbf{$g$} &
\textbf{Condition label} \\
\midrule
low  & low  & (0) community agrees, casual \\
low  & high & (1) community agrees, pressing \\
high & low  & (2) community divided, casual \\
high & high & (3) community divided, pressing \\
\bottomrule
\end{tabular}
\end{table}

\subsubsection{Counterbalancing}
\label{subsub:counterbalancing}

The experiment used a \textbf{Latin square} to counterbalance item--condition
pairings across participants.
With $k = 4$ conditions and 8 items, a balanced assignment is achieved:
each of the 4 lists assigns item $i$ ($i = 0,\ldots,7$) to condition
$(i + l)\bmod 4$, where $l \in \{0,1,2,3\}$ is the list number.
The resulting assignment matrix (Table~\ref{tab:ls}) guarantees that
\begin{itemize}
  \item each participant sees every item exactly once;
  \item each participant sees each condition exactly twice (two items per
        condition), ensuring balanced coverage of the 2$\times$2 design;
  \item across all 4 lists, each item appears equally often in each
        condition (twice per condition over the full participant sample).
\end{itemize}

\begin{table}[ht]
\centering
\caption{Latin square assignment of conditions to items across 4 lists.
  Cell entries are condition indices (0--3) as defined in
  Table~\ref{tab:design}.
  Rows = lists ($l$); columns = items ($i$).}
\label{tab:ls}
\small
\begin{tabular}{ccccccccc}
\toprule
\textbf{List} &
\textbf{Item 1} & \textbf{Item 2} & \textbf{Item 3} & \textbf{Item 4} &
\textbf{Item 5} & \textbf{Item 6} & \textbf{Item 7} & \textbf{Item 8} \\
\midrule
0 & 0 & 1 & 2 & 3 & 0 & 1 & 2 & 3 \\
1 & 1 & 2 & 3 & 0 & 1 & 2 & 3 & 0 \\
2 & 2 & 3 & 0 & 1 & 2 & 3 & 0 & 1 \\
3 & 3 & 0 & 1 & 2 & 3 & 0 & 1 & 2 \\
\bottomrule
\end{tabular}
\end{table}

\noindent List assignment is determined automatically from the URL query
parameter \texttt{?list=N} (0--3); if absent, a random list is selected.
Each participant is thus within-subjects across items and sees all four
conditions (twice each). Condition effects are estimated from within-participant
variation across the 2$\times$2 design, controlling for by-participant
random intercepts.

In addition to the 8 critical trials, each participant completed 8
\textbf{filler trials} drawn from separate topic domains (train travel,
reading light, cooking, coffee, vacation, hydration, fresh air, noise).
Filler items used the same trial format but were excluded from the
main analysis; they serve to dilute the critical items and reduce
potential demand characteristics.

\subsubsection{Dependent variable}

The dependent variable was the marker chosen from a drop-down menu
displaying all five options in fixed order (as listed in
Section~\ref{sub:markers}).
No response time threshold was imposed, but response times were logged.
Recorded per trial: \texttt{trial\_id}, \texttt{topic}, \texttt{is\_filler},
\texttt{condition\_index}, \texttt{pc\_prag}, \texttt{g},
\texttt{selected\_marker}, \texttt{rt} (ms); and once per participant:
\texttt{list\_num}.
The continuous \texttt{pc\_prop\_rating} covariate (topic-level norming mean)
is joined to the trial data during analysis.

\subsection{Procedure}

Participants completed the experiment online via the
magpie platform \citep{kochetkova2024magpie} in a browser.
After reading instructions in German, participants were assigned to one
of 4 counterbalancing lists (via URL parameter or randomly).
They then proceeded through 16 trials---8 critical and 8 fillers---in a
fully randomised order, followed by a post-test questionnaire eliciting
native language, age, gender, education level, and optional comments.
Each trial displayed (a) a context paragraph, (b) a goal instruction,
and (c) the sentence frame with the marker drop-down embedded at the
Mittelfeld position.
The ``Weiter'' (Next) button was disabled until a marker had been
selected, ensuring a response on every trial.

\end{document}
.
%
% Compile stand-alone:
%   pdflatex production_exp.tex
%
% Embed in main.tex:
%   % production_exp.tex
% Stand-alone description of the production experiment.
% Can be embedded in main.tex via \input{production_exp}.
%
% Compile stand-alone:
%   pdflatex production_exp.tex
%
% Embed in main.tex:
%   \input{production_exp}   % (inside a \section{Experiment} environment)

\documentclass[12pt]{article}
\usepackage{CSPArticleStyle}
\usepackage{csquotes}
\usepackage[margin=1in]{geometry}
\usepackage{booktabs}
\usepackage{parskip}
\usepackage{amsmath}
\usepackage{times}
\usepackage[T1]{fontenc}
\usepackage[utf8]{inputenc}
\usepackage{xcolor}
\usepackage{tabularray}   % robust tables; fall back to tabular if unavailable
\usepackage{gb4e}         % linguistic examples
\noautomath

\title{Production Experiment: Eliciting Consensus-Marker Choice\\
       Across Perceived-Controversy and Goal-Strength Conditions}
\author{}
\date{}

\begin{document}

%% ============================================================
%% When embedding in main.tex, delete from \documentclass to
%% \begin{document} and delete \end{document} at the bottom.
%% ============================================================

\section{Experiment: Consensus-Marker Production}

\subsection{Overview}

The production experiment empirically estimates the decision thresholds
that govern the use of five German consensus markers by placing
participants in the speaker role.
Participants read short vignettes that independently varied
\emph{pragmatic perceived controversy} (\textsc{pc\textsubscript{prag}})
and the speaker's \emph{persuasive-goal strength} ($g$), then selected
which consensus marker they would most naturally use in the described
situation.
\emph{Propositional perceived controversy} (\textsc{pc\textsubscript{prop}})
is held constant within the experiment and operationalised as a continuous
topic-level covariate obtained from a separate norming study: for each
of the eight topic items, participants rated the degree to which the
proposition $p$ is considered established/accepted (0--100 scale);
the resulting means enter the regression analyses as a fixed covariate.
The resulting distribution of marker choices across the
$\textsc{pc}_\text{prag} \times g$ design, combined with the
continuous \textsc{pc\textsubscript{prop}} covariate, provides the
empirical basis for estimating the ordinal thresholds
$\mu_1 < \cdots < \mu_4$ of the noisy-threshold RSA model.

\subsection{Participants}

Participants will be recruited via Prolific and will be required to be
native speakers of German (language-at-home filter).
Data collection is ongoing; a target of $N = 80$ participants is planned
($n = 20$ per list across 4 lists), each completing 16 trials (8 critical + 8 filler).

\subsection{Materials}

\subsubsection{Markers}

The five consensus markers under investigation are, from weakest to
strongest presumed common-ground commitment:

\begin{enumerate}
  \item \textit{sofern ich wei\ss{}} (`as far as I know')
  \item \textit{wie du ja wei\ss t} (`as you (prt) know')
  \item \textit{wie wir wissen} (`as we know')
  \item \textit{ja} (modal particle; roughly `as you know / of course')
  \item \textit{bekanntlich} (`as is well known / notoriously')
\end{enumerate}

These markers were embedded in the sentential \emph{Mittelfeld}
(middle field), immediately after the finite verb, as illustrated in
example~(\ref{ex:frame}):

\begin{exe}
  \ex \label{ex:frame}
  Digitale Kompetenzen sind \textbf{[MARKER]}
  in der heutigen Arbeitswelt unverzichtbar.\\
  `Digital skills are \textbf{[MARKER]} indispensable in today's
  working world.'
\end{exe}

This syntactic position is licit for all five markers: modal particles
such as \textit{ja} canonically occupy the left edge of the Mittelfeld
\citep{thurmair1989modalpartikeln}, while the adverb \textit{bekanntlich}
and parenthetical clauses (\textit{wie wir wissen}, etc.) can appear
there as well (see \citealt{jacobs_semantics_1991} for discussion).
Each target sentence was followed by a consequent clause with the
recommended action $q$: \textit{Deshalb solltest du [q].}
(`Therefore, you should [q].')

\subsubsection{Topic Items}

Eight topic domains were used, each associated with a fixed proposition
$p$ (sentence frame) and consequent action $q$ (Table~\ref{tab:items}).
The first four items were used in earlier pilot work; items 5--8 were
added to reach the minimum of eight required for a complete Latin square.

\begin{table}[ht]
\centering
\caption{Topic items: sentence frame for $p$ (with marker slot [M]) and consequent $q$.}
\label{tab:items}
\small
\begin{tabular}{clll}
\toprule
\textbf{\#} & \textbf{Topic} & \textbf{Sentence frame} & \textbf{Consequent $q$} \\
\midrule
1 & Klimawandel &
  \textit{Der Klimawandel wird [M]} &
  auf Flugreisen verzichten \\
  & & \textit{durch menschliche Aktivit\"aten verursacht.} & \\[3pt]
2 & Ern\"ahrung &
  \textit{Eine \"uberwiegend pflanzliche Ern\"ahrung} &
  mehr pflanzliche Lebens- \\
  & & \textit{f\"ordert [M] die Gesundheit.} &
  mittel einbauen \\[3pt]
3 & Stadtverkehr &
  \textit{Ein ausgebautes Nahverkehrsnetz} &
  h\"aufiger auf \"offentliche \\
  & & \textit{entlastet [M] den Stadtverkehr deutlich.} &
  Verkehrsmittel umsteigen \\[3pt]
4 & Digitalisierung &
  \textit{Digitale Kompetenzen sind [M] in der} &
  regelm\"a\ss{}ig Weiterbildungs- \\
  & & \textit{heutigen Arbeitswelt unverzichtbar.} &
  angebote nutzen \\[3pt]
5 & Schlaf &
  \textit{Ausreichend Schlaf ist [M] entscheidend} &
  auf feste Schlafzeiten \\
  & & \textit{f\"ur die kognitive Leistungsf\"ahigkeit.} &
  achten \\[3pt]
6 & Plastik &
  \textit{Einwegplastik schadet [M] den} &
  konsequent auf Einweg- \\
  & & \textit{Meeres\"okosystemen erheblich.} &
  plastik verzichten \\[3pt]
7 & Sport &
  \textit{Regelm\"a\ss{}ige k\"orperliche Bewegung} &
  regelm\"a\ss{}ig Sport in \\
  & & \textit{senkt [M] das Risiko f\"ur} &
  den Alltag integrieren \\
  & & \textit{Herz-Kreislauf-Erkrankungen deutlich.} & \\[3pt]
8 & Lokalkauf &
  \textit{Das Kaufen bei lokalen H\"andlern} &
  \"ofter bei lokalen \\
  & & \textit{st\"arkt [M] die regionale Wirtschaft} &
  Gesch\"aften einkaufen \\
  & & \textit{nachhaltig.} & \\
\bottomrule
\end{tabular}
\end{table}

\noindent [M] denotes the marker slot (dropdown).

\subsubsection{Vignettes}

\paragraph{Propositional perceived controversy (\textsc{pc\textsubscript{prop}}).}
\textsc{pc\textsubscript{prop}} was \emph{not} manipulated within the
production experiment.
Instead, it was operationalised as a continuous, topic-level covariate
obtained from a separate norming study conducted prior to the production
experiment.
In that norming study, a different sample of native German speakers rated,
for each topic proposition $p$, the degree to which they considered $p$
accepted/established in (German) public discourse, on a 0--100 scale
($0 = $ `completely contested', $100 = $ `universally accepted').
The resulting per-topic means (e.g., \textit{Digitalisierung}: 85/100;
\textit{Lokalkauf}: 58/100) enter the regression analyses as a
continuous, centred predictor \textsc{pc\textsubscript{prop}}$_c$
(mean-centred, SD $\approx$ 1).
Context paragraphs in the production experiment therefore present
the proposition $p$ without any editorial evaluation, so that
participants' experience of propositional controversy is shaped
solely by their own prior knowledge and the norming-derived covariate
captures this variation in analysis.

\paragraph{Pragmatic perceived controversy (\textsc{pc\textsubscript{prag}}).}
\textsc{pc\textsubscript{prag}} was manipulated between instances of each item
by varying whether the speaker's social circle \emph{largely agreed} (low)
or \emph{held divergent views} (high) about $p$.
The two context versions for the \textit{Digitalisierung} item are:

\begin{itemize}
  \item $\textsc{pc}_\text{prag} = \text{low}$: \textit{Du sprichst mit einer
        Kollegin {\"u}ber Berufsausbildung und Weiterbildung.
        In eurem Team teilen die meisten die Ansicht,
        dass digitale Kompetenzen in der heutigen Arbeitswelt unverzichtbar sind.}
  \item $\textsc{pc}_\text{prag} = \text{high}$: \textit{Du sprichst mit einer
        Kollegin {\"u}ber Berufsausbildung und Weiterbildung.
        In eurem Team gehen die Meinungen dazu stark auseinander.}
\end{itemize}

\paragraph{Speaker goal strength ($g$).}
Goal strength was manipulated via a standardised goal instruction
appended after the context paragraph:

\begin{itemize}
  \item $g_\text{low}$: \textit{Du m\"ochtest das nur kurz anmerken.}
        (`You just want to make a brief note of this.')
  \item $g_\text{high}$: \textit{Es ist dir sehr wichtig, dass dein
        Gespr\"achspartner das akzeptiert.}
        (`It is very important to you that your interlocutor accepts this.')
\end{itemize}

\paragraph{Complete sample item (all four conditions).}
Examples (\ref{ex:cond1})--(\ref{ex:cond4}) show the full vignette for
the topic \textit{Digitalisierung}
(norming $\textsc{pc}_\text{prop}$ rating: 85/100;
sentence frame: \textit{Digitale Kompetenzen sind [M] in der heutigen
Arbeitswelt unverzichtbar.
Deshalb solltest du regelm\"a\ss{}ig digitale Weiterbildungsangebote nutzen.})
across all four conditions of the $2\times2$ design.

\begin{exe}
  \ex \label{ex:cond1}
  \textbf{Condition 0: $\textsc{pc}_\text{prag} = \text{low}$, $g = \text{low}$}\\
  \textit{Du sprichst mit einer Kollegin {\"u}ber Berufsausbildung und Weiterbildung.
  In eurem Team teilen die meisten die Ansicht,
  dass digitale Kompetenzen in der heutigen Arbeitswelt unverzichtbar sind.}\\
  \textit{Du m\"ochtest das nur kurz anmerken.}\\
  Digitale Kompetenzen sind \underline{\hspace{2.2cm}}
  in der heutigen Arbeitswelt unverzichtbar.

  \ex \label{ex:cond2}
  \textbf{Condition 1: $\textsc{pc}_\text{prag} = \text{low}$, $g = \text{high}$}\\
  \textit{Du sprichst mit einer Kollegin {\"u}ber Berufsausbildung und Weiterbildung.
  In eurem Team teilen die meisten die Ansicht,
  dass digitale Kompetenzen in der heutigen Arbeitswelt unverzichtbar sind.}\\
  \textit{Es ist dir sehr wichtig, dass dein Gespr\"achspartner das akzeptiert.}\\
  Digitale Kompetenzen sind \underline{\hspace{2.2cm}}
  in der heutigen Arbeitswelt unverzichtbar.

  \ex \label{ex:cond3}
  \textbf{Condition 2: $\textsc{pc}_\text{prag} = \text{high}$, $g = \text{low}$}\\
  \textit{Du sprichst mit einer Kollegin {\"u}ber Berufsausbildung und Weiterbildung.
  In eurem Team gehen die Meinungen dazu stark auseinander.}\\
  \textit{Du m\"ochtest das nur kurz anmerken.}\\
  Digitale Kompetenzen sind \underline{\hspace{2.2cm}}
  in der heutigen Arbeitswelt unverzichtbar.

  \ex \label{ex:cond4}
  \textbf{Condition 3: $\textsc{pc}_\text{prag} = \text{high}$, $g = \text{high}$}\\
  \textit{Du sprichst mit einer Kollegin {\"u}ber Berufsausbildung und Weiterbildung.
  In eurem Team gehen die Meinungen dazu stark auseinander.}\\
  \textit{Es ist dir sehr wichtig, dass dein Gespr\"achspartner das akzeptiert.}\\
  Digitale Kompetenzen sind \underline{\hspace{2.2cm}}
  in der heutigen Arbeitswelt unverzichtbar.
\end{exe}

\noindent In all four vignettes the sentence frame and consequent clause
(\textit{Deshalb solltest du regelm\"a\ss{}ig digitale
Weiterbildungsangebote nutzen.}) are identical; only the context
paragraph and goal instruction vary across conditions.
The marker slot [M] / \underline{\hspace{2.2cm}} was presented as a
drop-down menu in the actual experiment.

\subsection{Design}

\subsubsection{Factorial structure}

The experiment employed a \textbf{2 (pc\textsubscript{prag}: low vs.\ high)
$\times$ 2 ($g$: low vs.\ high) fully crossed design},
yielding 4 conditions.
\textsc{pc\textsubscript{prop}} is not a within-experiment factor; it enters
as a continuous topic-level covariate from the norming study.
With 8 topic items and a Latin square rotation across 4 lists, each
participant saw exactly \textbf{8 critical trials}---two items per condition---plus
8 filler trials, for \textbf{16 trials} in total
(Table~\ref{tab:design}).

\begin{table}[ht]
\centering
\caption{%
  The $2\times2$ design (4 conditions).
  \textsc{pc\textsubscript{prop}} is a continuous norming covariate, not a
  within-experiment factor.
  In each list, the 8 items are assigned to conditions via the Latin square
  rotation described in Section~\ref{subsub:counterbalancing}.}
\label{tab:design}
\small
\begin{tabular}{ccl}
\toprule
\textbf{pc\textsubscript{prag}} &
\textbf{$g$} &
\textbf{Condition label} \\
\midrule
low  & low  & (0) community agrees, casual \\
low  & high & (1) community agrees, pressing \\
high & low  & (2) community divided, casual \\
high & high & (3) community divided, pressing \\
\bottomrule
\end{tabular}
\end{table}

\subsubsection{Counterbalancing}
\label{subsub:counterbalancing}

The experiment used a \textbf{Latin square} to counterbalance item--condition
pairings across participants.
With $k = 4$ conditions and 8 items, a balanced assignment is achieved:
each of the 4 lists assigns item $i$ ($i = 0,\ldots,7$) to condition
$(i + l)\bmod 4$, where $l \in \{0,1,2,3\}$ is the list number.
The resulting assignment matrix (Table~\ref{tab:ls}) guarantees that
\begin{itemize}
  \item each participant sees every item exactly once;
  \item each participant sees each condition exactly twice (two items per
        condition), ensuring balanced coverage of the 2$\times$2 design;
  \item across all 4 lists, each item appears equally often in each
        condition (twice per condition over the full participant sample).
\end{itemize}

\begin{table}[ht]
\centering
\caption{Latin square assignment of conditions to items across 4 lists.
  Cell entries are condition indices (0--3) as defined in
  Table~\ref{tab:design}.
  Rows = lists ($l$); columns = items ($i$).}
\label{tab:ls}
\small
\begin{tabular}{ccccccccc}
\toprule
\textbf{List} &
\textbf{Item 1} & \textbf{Item 2} & \textbf{Item 3} & \textbf{Item 4} &
\textbf{Item 5} & \textbf{Item 6} & \textbf{Item 7} & \textbf{Item 8} \\
\midrule
0 & 0 & 1 & 2 & 3 & 0 & 1 & 2 & 3 \\
1 & 1 & 2 & 3 & 0 & 1 & 2 & 3 & 0 \\
2 & 2 & 3 & 0 & 1 & 2 & 3 & 0 & 1 \\
3 & 3 & 0 & 1 & 2 & 3 & 0 & 1 & 2 \\
\bottomrule
\end{tabular}
\end{table}

\noindent List assignment is determined automatically from the URL query
parameter \texttt{?list=N} (0--3); if absent, a random list is selected.
Each participant is thus within-subjects across items and sees all four
conditions (twice each). Condition effects are estimated from within-participant
variation across the 2$\times$2 design, controlling for by-participant
random intercepts.

In addition to the 8 critical trials, each participant completed 8
\textbf{filler trials} drawn from separate topic domains (train travel,
reading light, cooking, coffee, vacation, hydration, fresh air, noise).
Filler items used the same trial format but were excluded from the
main analysis; they serve to dilute the critical items and reduce
potential demand characteristics.

\subsubsection{Dependent variable}

The dependent variable was the marker chosen from a drop-down menu
displaying all five options in fixed order (as listed in
Section~\ref{sub:markers}).
No response time threshold was imposed, but response times were logged.
Recorded per trial: \texttt{trial\_id}, \texttt{topic}, \texttt{is\_filler},
\texttt{condition\_index}, \texttt{pc\_prag}, \texttt{g},
\texttt{selected\_marker}, \texttt{rt} (ms); and once per participant:
\texttt{list\_num}.
The continuous \texttt{pc\_prop\_rating} covariate (topic-level norming mean)
is joined to the trial data during analysis.

\subsection{Procedure}

Participants completed the experiment online via the
magpie platform \citep{kochetkova2024magpie} in a browser.
After reading instructions in German, participants were assigned to one
of 4 counterbalancing lists (via URL parameter or randomly).
They then proceeded through 16 trials---8 critical and 8 fillers---in a
fully randomised order, followed by a post-test questionnaire eliciting
native language, age, gender, education level, and optional comments.
Each trial displayed (a) a context paragraph, (b) a goal instruction,
and (c) the sentence frame with the marker drop-down embedded at the
Mittelfeld position.
The ``Weiter'' (Next) button was disabled until a marker had been
selected, ensuring a response on every trial.

\end{document}
   % (inside a \section{Experiment} environment)

\documentclass[12pt]{article}
\usepackage{CSPArticleStyle}
\usepackage{csquotes}
\usepackage[margin=1in]{geometry}
\usepackage{booktabs}
\usepackage{parskip}
\usepackage{amsmath}
\usepackage{times}
\usepackage[T1]{fontenc}
\usepackage[utf8]{inputenc}
\usepackage{xcolor}
\usepackage{tabularray}   % robust tables; fall back to tabular if unavailable
\usepackage{gb4e}         % linguistic examples
\noautomath

\title{Production Experiment: Eliciting Consensus-Marker Choice\\
       Across Perceived-Controversy and Goal-Strength Conditions}
\author{}
\date{}

\begin{document}

%% ============================================================
%% When embedding in main.tex, delete from \documentclass to
%% \begin{document} and delete \end{document} at the bottom.
%% ============================================================

\section{Experiment: Consensus-Marker Production}

\subsection{Overview}

The production experiment empirically estimates the decision thresholds
that govern the use of five German consensus markers by placing
participants in the speaker role.
Participants read short vignettes that independently varied
\emph{pragmatic perceived controversy} (\textsc{pc\textsubscript{prag}})
and the speaker's \emph{persuasive-goal strength} ($g$), then selected
which consensus marker they would most naturally use in the described
situation.
\emph{Propositional perceived controversy} (\textsc{pc\textsubscript{prop}})
is held constant within the experiment and operationalised as a continuous
topic-level covariate obtained from a separate norming study: for each
of the eight topic items, participants rated the degree to which the
proposition $p$ is considered established/accepted (0--100 scale);
the resulting means enter the regression analyses as a fixed covariate.
The resulting distribution of marker choices across the
$\textsc{pc}_\text{prag} \times g$ design, combined with the
continuous \textsc{pc\textsubscript{prop}} covariate, provides the
empirical basis for estimating the ordinal thresholds
$\mu_1 < \cdots < \mu_4$ of the noisy-threshold RSA model.

\subsection{Participants}

Participants will be recruited via Prolific and will be required to be
native speakers of German (language-at-home filter).
Data collection is ongoing; a target of $N = 80$ participants is planned
($n = 20$ per list across 4 lists), each completing 16 trials (8 critical + 8 filler).

\subsection{Materials}

\subsubsection{Markers}

The five consensus markers under investigation are, from weakest to
strongest presumed common-ground commitment:

\begin{enumerate}
  \item \textit{sofern ich wei\ss{}} (`as far as I know')
  \item \textit{wie du ja wei\ss t} (`as you (prt) know')
  \item \textit{wie wir wissen} (`as we know')
  \item \textit{ja} (modal particle; roughly `as you know / of course')
  \item \textit{bekanntlich} (`as is well known / notoriously')
\end{enumerate}

These markers were embedded in the sentential \emph{Mittelfeld}
(middle field), immediately after the finite verb, as illustrated in
example~(\ref{ex:frame}):

\begin{exe}
  \ex \label{ex:frame}
  Digitale Kompetenzen sind \textbf{[MARKER]}
  in der heutigen Arbeitswelt unverzichtbar.\\
  `Digital skills are \textbf{[MARKER]} indispensable in today's
  working world.'
\end{exe}

This syntactic position is licit for all five markers: modal particles
such as \textit{ja} canonically occupy the left edge of the Mittelfeld
\citep{thurmair1989modalpartikeln}, while the adverb \textit{bekanntlich}
and parenthetical clauses (\textit{wie wir wissen}, etc.) can appear
there as well (see \citealt{jacobs_semantics_1991} for discussion).
Each target sentence was followed by a consequent clause with the
recommended action $q$: \textit{Deshalb solltest du [q].}
(`Therefore, you should [q].')

\subsubsection{Topic Items}

Eight topic domains were used, each associated with a fixed proposition
$p$ (sentence frame) and consequent action $q$ (Table~\ref{tab:items}).
The first four items were used in earlier pilot work; items 5--8 were
added to reach the minimum of eight required for a complete Latin square.

\begin{table}[ht]
\centering
\caption{Topic items: sentence frame for $p$ (with marker slot [M]) and consequent $q$.}
\label{tab:items}
\small
\begin{tabular}{clll}
\toprule
\textbf{\#} & \textbf{Topic} & \textbf{Sentence frame} & \textbf{Consequent $q$} \\
\midrule
1 & Klimawandel &
  \textit{Der Klimawandel wird [M]} &
  auf Flugreisen verzichten \\
  & & \textit{durch menschliche Aktivit\"aten verursacht.} & \\[3pt]
2 & Ern\"ahrung &
  \textit{Eine \"uberwiegend pflanzliche Ern\"ahrung} &
  mehr pflanzliche Lebens- \\
  & & \textit{f\"ordert [M] die Gesundheit.} &
  mittel einbauen \\[3pt]
3 & Stadtverkehr &
  \textit{Ein ausgebautes Nahverkehrsnetz} &
  h\"aufiger auf \"offentliche \\
  & & \textit{entlastet [M] den Stadtverkehr deutlich.} &
  Verkehrsmittel umsteigen \\[3pt]
4 & Digitalisierung &
  \textit{Digitale Kompetenzen sind [M] in der} &
  regelm\"a\ss{}ig Weiterbildungs- \\
  & & \textit{heutigen Arbeitswelt unverzichtbar.} &
  angebote nutzen \\[3pt]
5 & Schlaf &
  \textit{Ausreichend Schlaf ist [M] entscheidend} &
  auf feste Schlafzeiten \\
  & & \textit{f\"ur die kognitive Leistungsf\"ahigkeit.} &
  achten \\[3pt]
6 & Plastik &
  \textit{Einwegplastik schadet [M] den} &
  konsequent auf Einweg- \\
  & & \textit{Meeres\"okosystemen erheblich.} &
  plastik verzichten \\[3pt]
7 & Sport &
  \textit{Regelm\"a\ss{}ige k\"orperliche Bewegung} &
  regelm\"a\ss{}ig Sport in \\
  & & \textit{senkt [M] das Risiko f\"ur} &
  den Alltag integrieren \\
  & & \textit{Herz-Kreislauf-Erkrankungen deutlich.} & \\[3pt]
8 & Lokalkauf &
  \textit{Das Kaufen bei lokalen H\"andlern} &
  \"ofter bei lokalen \\
  & & \textit{st\"arkt [M] die regionale Wirtschaft} &
  Gesch\"aften einkaufen \\
  & & \textit{nachhaltig.} & \\
\bottomrule
\end{tabular}
\end{table}

\noindent [M] denotes the marker slot (dropdown).

\subsubsection{Vignettes}

\paragraph{Propositional perceived controversy (\textsc{pc\textsubscript{prop}}).}
\textsc{pc\textsubscript{prop}} was \emph{not} manipulated within the
production experiment.
Instead, it was operationalised as a continuous, topic-level covariate
obtained from a separate norming study conducted prior to the production
experiment.
In that norming study, a different sample of native German speakers rated,
for each topic proposition $p$, the degree to which they considered $p$
accepted/established in (German) public discourse, on a 0--100 scale
($0 = $ `completely contested', $100 = $ `universally accepted').
The resulting per-topic means (e.g., \textit{Digitalisierung}: 85/100;
\textit{Lokalkauf}: 58/100) enter the regression analyses as a
continuous, centred predictor \textsc{pc\textsubscript{prop}}$_c$
(mean-centred, SD $\approx$ 1).
Context paragraphs in the production experiment therefore present
the proposition $p$ without any editorial evaluation, so that
participants' experience of propositional controversy is shaped
solely by their own prior knowledge and the norming-derived covariate
captures this variation in analysis.

\paragraph{Pragmatic perceived controversy (\textsc{pc\textsubscript{prag}}).}
\textsc{pc\textsubscript{prag}} was manipulated between instances of each item
by varying whether the speaker's social circle \emph{largely agreed} (low)
or \emph{held divergent views} (high) about $p$.
The two context versions for the \textit{Digitalisierung} item are:

\begin{itemize}
  \item $\textsc{pc}_\text{prag} = \text{low}$: \textit{Du sprichst mit einer
        Kollegin {\"u}ber Berufsausbildung und Weiterbildung.
        In eurem Team teilen die meisten die Ansicht,
        dass digitale Kompetenzen in der heutigen Arbeitswelt unverzichtbar sind.}
  \item $\textsc{pc}_\text{prag} = \text{high}$: \textit{Du sprichst mit einer
        Kollegin {\"u}ber Berufsausbildung und Weiterbildung.
        In eurem Team gehen die Meinungen dazu stark auseinander.}
\end{itemize}

\paragraph{Speaker goal strength ($g$).}
Goal strength was manipulated via a standardised goal instruction
appended after the context paragraph:

\begin{itemize}
  \item $g_\text{low}$: \textit{Du m\"ochtest das nur kurz anmerken.}
        (`You just want to make a brief note of this.')
  \item $g_\text{high}$: \textit{Es ist dir sehr wichtig, dass dein
        Gespr\"achspartner das akzeptiert.}
        (`It is very important to you that your interlocutor accepts this.')
\end{itemize}

\paragraph{Complete sample item (all four conditions).}
Examples (\ref{ex:cond1})--(\ref{ex:cond4}) show the full vignette for
the topic \textit{Digitalisierung}
(norming $\textsc{pc}_\text{prop}$ rating: 85/100;
sentence frame: \textit{Digitale Kompetenzen sind [M] in der heutigen
Arbeitswelt unverzichtbar.
Deshalb solltest du regelm\"a\ss{}ig digitale Weiterbildungsangebote nutzen.})
across all four conditions of the $2\times2$ design.

\begin{exe}
  \ex \label{ex:cond1}
  \textbf{Condition 0: $\textsc{pc}_\text{prag} = \text{low}$, $g = \text{low}$}\\
  \textit{Du sprichst mit einer Kollegin {\"u}ber Berufsausbildung und Weiterbildung.
  In eurem Team teilen die meisten die Ansicht,
  dass digitale Kompetenzen in der heutigen Arbeitswelt unverzichtbar sind.}\\
  \textit{Du m\"ochtest das nur kurz anmerken.}\\
  Digitale Kompetenzen sind \underline{\hspace{2.2cm}}
  in der heutigen Arbeitswelt unverzichtbar.

  \ex \label{ex:cond2}
  \textbf{Condition 1: $\textsc{pc}_\text{prag} = \text{low}$, $g = \text{high}$}\\
  \textit{Du sprichst mit einer Kollegin {\"u}ber Berufsausbildung und Weiterbildung.
  In eurem Team teilen die meisten die Ansicht,
  dass digitale Kompetenzen in der heutigen Arbeitswelt unverzichtbar sind.}\\
  \textit{Es ist dir sehr wichtig, dass dein Gespr\"achspartner das akzeptiert.}\\
  Digitale Kompetenzen sind \underline{\hspace{2.2cm}}
  in der heutigen Arbeitswelt unverzichtbar.

  \ex \label{ex:cond3}
  \textbf{Condition 2: $\textsc{pc}_\text{prag} = \text{high}$, $g = \text{low}$}\\
  \textit{Du sprichst mit einer Kollegin {\"u}ber Berufsausbildung und Weiterbildung.
  In eurem Team gehen die Meinungen dazu stark auseinander.}\\
  \textit{Du m\"ochtest das nur kurz anmerken.}\\
  Digitale Kompetenzen sind \underline{\hspace{2.2cm}}
  in der heutigen Arbeitswelt unverzichtbar.

  \ex \label{ex:cond4}
  \textbf{Condition 3: $\textsc{pc}_\text{prag} = \text{high}$, $g = \text{high}$}\\
  \textit{Du sprichst mit einer Kollegin {\"u}ber Berufsausbildung und Weiterbildung.
  In eurem Team gehen die Meinungen dazu stark auseinander.}\\
  \textit{Es ist dir sehr wichtig, dass dein Gespr\"achspartner das akzeptiert.}\\
  Digitale Kompetenzen sind \underline{\hspace{2.2cm}}
  in der heutigen Arbeitswelt unverzichtbar.
\end{exe}

\noindent In all four vignettes the sentence frame and consequent clause
(\textit{Deshalb solltest du regelm\"a\ss{}ig digitale
Weiterbildungsangebote nutzen.}) are identical; only the context
paragraph and goal instruction vary across conditions.
The marker slot [M] / \underline{\hspace{2.2cm}} was presented as a
drop-down menu in the actual experiment.

\subsection{Design}

\subsubsection{Factorial structure}

The experiment employed a \textbf{2 (pc\textsubscript{prag}: low vs.\ high)
$\times$ 2 ($g$: low vs.\ high) fully crossed design},
yielding 4 conditions.
\textsc{pc\textsubscript{prop}} is not a within-experiment factor; it enters
as a continuous topic-level covariate from the norming study.
With 8 topic items and a Latin square rotation across 4 lists, each
participant saw exactly \textbf{8 critical trials}---two items per condition---plus
8 filler trials, for \textbf{16 trials} in total
(Table~\ref{tab:design}).

\begin{table}[ht]
\centering
\caption{%
  The $2\times2$ design (4 conditions).
  \textsc{pc\textsubscript{prop}} is a continuous norming covariate, not a
  within-experiment factor.
  In each list, the 8 items are assigned to conditions via the Latin square
  rotation described in Section~\ref{subsub:counterbalancing}.}
\label{tab:design}
\small
\begin{tabular}{ccl}
\toprule
\textbf{pc\textsubscript{prag}} &
\textbf{$g$} &
\textbf{Condition label} \\
\midrule
low  & low  & (0) community agrees, casual \\
low  & high & (1) community agrees, pressing \\
high & low  & (2) community divided, casual \\
high & high & (3) community divided, pressing \\
\bottomrule
\end{tabular}
\end{table}

\subsubsection{Counterbalancing}
\label{subsub:counterbalancing}

The experiment used a \textbf{Latin square} to counterbalance item--condition
pairings across participants.
With $k = 4$ conditions and 8 items, a balanced assignment is achieved:
each of the 4 lists assigns item $i$ ($i = 0,\ldots,7$) to condition
$(i + l)\bmod 4$, where $l \in \{0,1,2,3\}$ is the list number.
The resulting assignment matrix (Table~\ref{tab:ls}) guarantees that
\begin{itemize}
  \item each participant sees every item exactly once;
  \item each participant sees each condition exactly twice (two items per
        condition), ensuring balanced coverage of the 2$\times$2 design;
  \item across all 4 lists, each item appears equally often in each
        condition (twice per condition over the full participant sample).
\end{itemize}

\begin{table}[ht]
\centering
\caption{Latin square assignment of conditions to items across 4 lists.
  Cell entries are condition indices (0--3) as defined in
  Table~\ref{tab:design}.
  Rows = lists ($l$); columns = items ($i$).}
\label{tab:ls}
\small
\begin{tabular}{ccccccccc}
\toprule
\textbf{List} &
\textbf{Item 1} & \textbf{Item 2} & \textbf{Item 3} & \textbf{Item 4} &
\textbf{Item 5} & \textbf{Item 6} & \textbf{Item 7} & \textbf{Item 8} \\
\midrule
0 & 0 & 1 & 2 & 3 & 0 & 1 & 2 & 3 \\
1 & 1 & 2 & 3 & 0 & 1 & 2 & 3 & 0 \\
2 & 2 & 3 & 0 & 1 & 2 & 3 & 0 & 1 \\
3 & 3 & 0 & 1 & 2 & 3 & 0 & 1 & 2 \\
\bottomrule
\end{tabular}
\end{table}

\noindent List assignment is determined automatically from the URL query
parameter \texttt{?list=N} (0--3); if absent, a random list is selected.
Each participant is thus within-subjects across items and sees all four
conditions (twice each). Condition effects are estimated from within-participant
variation across the 2$\times$2 design, controlling for by-participant
random intercepts.

In addition to the 8 critical trials, each participant completed 8
\textbf{filler trials} drawn from separate topic domains (train travel,
reading light, cooking, coffee, vacation, hydration, fresh air, noise).
Filler items used the same trial format but were excluded from the
main analysis; they serve to dilute the critical items and reduce
potential demand characteristics.

\subsubsection{Dependent variable}

The dependent variable was the marker chosen from a drop-down menu
displaying all five options in fixed order (as listed in
Section~\ref{sub:markers}).
No response time threshold was imposed, but response times were logged.
Recorded per trial: \texttt{trial\_id}, \texttt{topic}, \texttt{is\_filler},
\texttt{condition\_index}, \texttt{pc\_prag}, \texttt{g},
\texttt{selected\_marker}, \texttt{rt} (ms); and once per participant:
\texttt{list\_num}.
The continuous \texttt{pc\_prop\_rating} covariate (topic-level norming mean)
is joined to the trial data during analysis.

\subsection{Procedure}

Participants completed the experiment online via the
magpie platform \citep{kochetkova2024magpie} in a browser.
After reading instructions in German, participants were assigned to one
of 4 counterbalancing lists (via URL parameter or randomly).
They then proceeded through 16 trials---8 critical and 8 fillers---in a
fully randomised order, followed by a post-test questionnaire eliciting
native language, age, gender, education level, and optional comments.
Each trial displayed (a) a context paragraph, (b) a goal instruction,
and (c) the sentence frame with the marker drop-down embedded at the
Mittelfeld position.
The ``Weiter'' (Next) button was disabled until a marker had been
selected, ensuring a response on every trial.

\end{document}
   % (inside a \section{Experiment} environment)

\documentclass[12pt]{article}
\usepackage{CSPArticleStyle}
\usepackage{csquotes}
\usepackage[margin=1in]{geometry}
\usepackage{booktabs}
\usepackage{parskip}
\usepackage{amsmath}
\usepackage{times}
\usepackage[T1]{fontenc}
\usepackage[utf8]{inputenc}
\usepackage{xcolor}
\usepackage{tabularray}   % robust tables; fall back to tabular if unavailable
\usepackage{gb4e}         % linguistic examples
\noautomath

\title{Production Experiment: Eliciting Consensus-Marker Choice\\
       Across Perceived-Controversy and Goal-Strength Conditions}
\author{}
\date{}

\begin{document}

%% ============================================================
%% When embedding in main.tex, delete from \documentclass to
%% \begin{document} and delete \end{document} at the bottom.
%% ============================================================

\section{Experiment: Consensus-Marker Production}

\subsection{Overview}

The production experiment empirically estimates the decision thresholds
that govern the use of five German consensus markers by placing
participants in the speaker role.
Participants read short vignettes that independently varied
\emph{pragmatic perceived controversy} (\textsc{pc\textsubscript{prag}})
and the speaker's \emph{persuasive-goal strength} ($g$), then selected
which consensus marker they would most naturally use in the described
situation.
\emph{Propositional perceived controversy} (\textsc{pc\textsubscript{prop}})
is held constant within the experiment and operationalised as a continuous
topic-level covariate obtained from a separate norming study: for each
of the eight topic items, participants rated the degree to which the
proposition $p$ is considered established/accepted (0--100 scale);
the resulting means enter the regression analyses as a fixed covariate.
The resulting distribution of marker choices across the
$\textsc{pc}_\text{prag} \times g$ design, combined with the
continuous \textsc{pc\textsubscript{prop}} covariate, provides the
empirical basis for estimating the ordinal thresholds
$\mu_1 < \cdots < \mu_4$ of the noisy-threshold RSA model.

\subsection{Participants}

Participants will be recruited via Prolific and will be required to be
native speakers of German (language-at-home filter).
Data collection is ongoing; a target of $N = 80$ participants is planned
($n = 20$ per list across 4 lists), each completing 16 trials (8 critical + 8 filler).

\subsection{Materials}

\subsubsection{Markers}

The five consensus markers under investigation are, from weakest to
strongest presumed common-ground commitment:

\begin{enumerate}
  \item \textit{sofern ich wei\ss{}} (`as far as I know')
  \item \textit{wie du ja wei\ss t} (`as you (prt) know')
  \item \textit{wie wir wissen} (`as we know')
  \item \textit{ja} (modal particle; roughly `as you know / of course')
  \item \textit{bekanntlich} (`as is well known / notoriously')
\end{enumerate}

These markers were embedded in the sentential \emph{Mittelfeld}
(middle field), immediately after the finite verb, as illustrated in
example~(\ref{ex:frame}):

\begin{exe}
  \ex \label{ex:frame}
  Digitale Kompetenzen sind \textbf{[MARKER]}
  in der heutigen Arbeitswelt unverzichtbar.\\
  `Digital skills are \textbf{[MARKER]} indispensable in today's
  working world.'
\end{exe}

This syntactic position is licit for all five markers: modal particles
such as \textit{ja} canonically occupy the left edge of the Mittelfeld
\citep{thurmair1989modalpartikeln}, while the adverb \textit{bekanntlich}
and parenthetical clauses (\textit{wie wir wissen}, etc.) can appear
there as well (see \citealt{jacobs_semantics_1991} for discussion).
Each target sentence was followed by a consequent clause with the
recommended action $q$: \textit{Deshalb solltest du [q].}
(`Therefore, you should [q].')

\subsubsection{Topic Items}

Eight topic domains were used, each associated with a fixed proposition
$p$ (sentence frame) and consequent action $q$ (Table~\ref{tab:items}).
The first four items were used in earlier pilot work; items 5--8 were
added to reach the minimum of eight required for a complete Latin square.

\begin{table}[ht]
\centering
\caption{Topic items: sentence frame for $p$ (with marker slot [M]) and consequent $q$.}
\label{tab:items}
\small
\begin{tabular}{clll}
\toprule
\textbf{\#} & \textbf{Topic} & \textbf{Sentence frame} & \textbf{Consequent $q$} \\
\midrule
1 & Klimawandel &
  \textit{Der Klimawandel wird [M]} &
  auf Flugreisen verzichten \\
  & & \textit{durch menschliche Aktivit\"aten verursacht.} & \\[3pt]
2 & Ern\"ahrung &
  \textit{Eine \"uberwiegend pflanzliche Ern\"ahrung} &
  mehr pflanzliche Lebens- \\
  & & \textit{f\"ordert [M] die Gesundheit.} &
  mittel einbauen \\[3pt]
3 & Stadtverkehr &
  \textit{Ein ausgebautes Nahverkehrsnetz} &
  h\"aufiger auf \"offentliche \\
  & & \textit{entlastet [M] den Stadtverkehr deutlich.} &
  Verkehrsmittel umsteigen \\[3pt]
4 & Digitalisierung &
  \textit{Digitale Kompetenzen sind [M] in der} &
  regelm\"a\ss{}ig Weiterbildungs- \\
  & & \textit{heutigen Arbeitswelt unverzichtbar.} &
  angebote nutzen \\[3pt]
5 & Schlaf &
  \textit{Ausreichend Schlaf ist [M] entscheidend} &
  auf feste Schlafzeiten \\
  & & \textit{f\"ur die kognitive Leistungsf\"ahigkeit.} &
  achten \\[3pt]
6 & Plastik &
  \textit{Einwegplastik schadet [M] den} &
  konsequent auf Einweg- \\
  & & \textit{Meeres\"okosystemen erheblich.} &
  plastik verzichten \\[3pt]
7 & Sport &
  \textit{Regelm\"a\ss{}ige k\"orperliche Bewegung} &
  regelm\"a\ss{}ig Sport in \\
  & & \textit{senkt [M] das Risiko f\"ur} &
  den Alltag integrieren \\
  & & \textit{Herz-Kreislauf-Erkrankungen deutlich.} & \\[3pt]
8 & Lokalkauf &
  \textit{Das Kaufen bei lokalen H\"andlern} &
  \"ofter bei lokalen \\
  & & \textit{st\"arkt [M] die regionale Wirtschaft} &
  Gesch\"aften einkaufen \\
  & & \textit{nachhaltig.} & \\
\bottomrule
\end{tabular}
\end{table}

\noindent [M] denotes the marker slot (dropdown).

\subsubsection{Vignettes}

\paragraph{Propositional perceived controversy (\textsc{pc\textsubscript{prop}}).}
\textsc{pc\textsubscript{prop}} was \emph{not} manipulated within the
production experiment.
Instead, it was operationalised as a continuous, topic-level covariate
obtained from a separate norming study conducted prior to the production
experiment.
In that norming study, a different sample of native German speakers rated,
for each topic proposition $p$, the degree to which they considered $p$
accepted/established in (German) public discourse, on a 0--100 scale
($0 = $ `completely contested', $100 = $ `universally accepted').
The resulting per-topic means (e.g., \textit{Digitalisierung}: 85/100;
\textit{Lokalkauf}: 58/100) enter the regression analyses as a
continuous, centred predictor \textsc{pc\textsubscript{prop}}$_c$
(mean-centred, SD $\approx$ 1).
Context paragraphs in the production experiment therefore present
the proposition $p$ without any editorial evaluation, so that
participants' experience of propositional controversy is shaped
solely by their own prior knowledge and the norming-derived covariate
captures this variation in analysis.

\paragraph{Pragmatic perceived controversy (\textsc{pc\textsubscript{prag}}).}
\textsc{pc\textsubscript{prag}} was manipulated between instances of each item
by varying whether the speaker's social circle \emph{largely agreed} (low)
or \emph{held divergent views} (high) about $p$.
The two context versions for the \textit{Digitalisierung} item are:

\begin{itemize}
  \item $\textsc{pc}_\text{prag} = \text{low}$: \textit{Du sprichst mit einer
        Kollegin {\"u}ber Berufsausbildung und Weiterbildung.
        In eurem Team teilen die meisten die Ansicht,
        dass digitale Kompetenzen in der heutigen Arbeitswelt unverzichtbar sind.}
  \item $\textsc{pc}_\text{prag} = \text{high}$: \textit{Du sprichst mit einer
        Kollegin {\"u}ber Berufsausbildung und Weiterbildung.
        In eurem Team gehen die Meinungen dazu stark auseinander.}
\end{itemize}

\paragraph{Speaker goal strength ($g$).}
Goal strength was manipulated via a standardised goal instruction
appended after the context paragraph:

\begin{itemize}
  \item $g_\text{low}$: \textit{Du m\"ochtest das nur kurz anmerken.}
        (`You just want to make a brief note of this.')
  \item $g_\text{high}$: \textit{Es ist dir sehr wichtig, dass dein
        Gespr\"achspartner das akzeptiert.}
        (`It is very important to you that your interlocutor accepts this.')
\end{itemize}

\paragraph{Complete sample item (all four conditions).}
Examples (\ref{ex:cond1})--(\ref{ex:cond4}) show the full vignette for
the topic \textit{Digitalisierung}
(norming $\textsc{pc}_\text{prop}$ rating: 85/100;
sentence frame: \textit{Digitale Kompetenzen sind [M] in der heutigen
Arbeitswelt unverzichtbar.
Deshalb solltest du regelm\"a\ss{}ig digitale Weiterbildungsangebote nutzen.})
across all four conditions of the $2\times2$ design.

\begin{exe}
  \ex \label{ex:cond1}
  \textbf{Condition 0: $\textsc{pc}_\text{prag} = \text{low}$, $g = \text{low}$}\\
  \textit{Du sprichst mit einer Kollegin {\"u}ber Berufsausbildung und Weiterbildung.
  In eurem Team teilen die meisten die Ansicht,
  dass digitale Kompetenzen in der heutigen Arbeitswelt unverzichtbar sind.}\\
  \textit{Du m\"ochtest das nur kurz anmerken.}\\
  Digitale Kompetenzen sind \underline{\hspace{2.2cm}}
  in der heutigen Arbeitswelt unverzichtbar.

  \ex \label{ex:cond2}
  \textbf{Condition 1: $\textsc{pc}_\text{prag} = \text{low}$, $g = \text{high}$}\\
  \textit{Du sprichst mit einer Kollegin {\"u}ber Berufsausbildung und Weiterbildung.
  In eurem Team teilen die meisten die Ansicht,
  dass digitale Kompetenzen in der heutigen Arbeitswelt unverzichtbar sind.}\\
  \textit{Es ist dir sehr wichtig, dass dein Gespr\"achspartner das akzeptiert.}\\
  Digitale Kompetenzen sind \underline{\hspace{2.2cm}}
  in der heutigen Arbeitswelt unverzichtbar.

  \ex \label{ex:cond3}
  \textbf{Condition 2: $\textsc{pc}_\text{prag} = \text{high}$, $g = \text{low}$}\\
  \textit{Du sprichst mit einer Kollegin {\"u}ber Berufsausbildung und Weiterbildung.
  In eurem Team gehen die Meinungen dazu stark auseinander.}\\
  \textit{Du m\"ochtest das nur kurz anmerken.}\\
  Digitale Kompetenzen sind \underline{\hspace{2.2cm}}
  in der heutigen Arbeitswelt unverzichtbar.

  \ex \label{ex:cond4}
  \textbf{Condition 3: $\textsc{pc}_\text{prag} = \text{high}$, $g = \text{high}$}\\
  \textit{Du sprichst mit einer Kollegin {\"u}ber Berufsausbildung und Weiterbildung.
  In eurem Team gehen die Meinungen dazu stark auseinander.}\\
  \textit{Es ist dir sehr wichtig, dass dein Gespr\"achspartner das akzeptiert.}\\
  Digitale Kompetenzen sind \underline{\hspace{2.2cm}}
  in der heutigen Arbeitswelt unverzichtbar.
\end{exe}

\noindent In all four vignettes the sentence frame and consequent clause
(\textit{Deshalb solltest du regelm\"a\ss{}ig digitale
Weiterbildungsangebote nutzen.}) are identical; only the context
paragraph and goal instruction vary across conditions.
The marker slot [M] / \underline{\hspace{2.2cm}} was presented as a
drop-down menu in the actual experiment.

\subsection{Design}

\subsubsection{Factorial structure}

The experiment employed a \textbf{2 (pc\textsubscript{prag}: low vs.\ high)
$\times$ 2 ($g$: low vs.\ high) fully crossed design},
yielding 4 conditions.
\textsc{pc\textsubscript{prop}} is not a within-experiment factor; it enters
as a continuous topic-level covariate from the norming study.
With 8 topic items and a Latin square rotation across 4 lists, each
participant saw exactly \textbf{8 critical trials}---two items per condition---plus
8 filler trials, for \textbf{16 trials} in total
(Table~\ref{tab:design}).

\begin{table}[ht]
\centering
\caption{%
  The $2\times2$ design (4 conditions).
  \textsc{pc\textsubscript{prop}} is a continuous norming covariate, not a
  within-experiment factor.
  In each list, the 8 items are assigned to conditions via the Latin square
  rotation described in Section~\ref{subsub:counterbalancing}.}
\label{tab:design}
\small
\begin{tabular}{ccl}
\toprule
\textbf{pc\textsubscript{prag}} &
\textbf{$g$} &
\textbf{Condition label} \\
\midrule
low  & low  & (0) community agrees, casual \\
low  & high & (1) community agrees, pressing \\
high & low  & (2) community divided, casual \\
high & high & (3) community divided, pressing \\
\bottomrule
\end{tabular}
\end{table}

\subsubsection{Counterbalancing}
\label{subsub:counterbalancing}

The experiment used a \textbf{Latin square} to counterbalance item--condition
pairings across participants.
With $k = 4$ conditions and 8 items, a balanced assignment is achieved:
each of the 4 lists assigns item $i$ ($i = 0,\ldots,7$) to condition
$(i + l)\bmod 4$, where $l \in \{0,1,2,3\}$ is the list number.
The resulting assignment matrix (Table~\ref{tab:ls}) guarantees that
\begin{itemize}
  \item each participant sees every item exactly once;
  \item each participant sees each condition exactly twice (two items per
        condition), ensuring balanced coverage of the 2$\times$2 design;
  \item across all 4 lists, each item appears equally often in each
        condition (twice per condition over the full participant sample).
\end{itemize}

\begin{table}[ht]
\centering
\caption{Latin square assignment of conditions to items across 4 lists.
  Cell entries are condition indices (0--3) as defined in
  Table~\ref{tab:design}.
  Rows = lists ($l$); columns = items ($i$).}
\label{tab:ls}
\small
\begin{tabular}{ccccccccc}
\toprule
\textbf{List} &
\textbf{Item 1} & \textbf{Item 2} & \textbf{Item 3} & \textbf{Item 4} &
\textbf{Item 5} & \textbf{Item 6} & \textbf{Item 7} & \textbf{Item 8} \\
\midrule
0 & 0 & 1 & 2 & 3 & 0 & 1 & 2 & 3 \\
1 & 1 & 2 & 3 & 0 & 1 & 2 & 3 & 0 \\
2 & 2 & 3 & 0 & 1 & 2 & 3 & 0 & 1 \\
3 & 3 & 0 & 1 & 2 & 3 & 0 & 1 & 2 \\
\bottomrule
\end{tabular}
\end{table}

\noindent List assignment is determined automatically from the URL query
parameter \texttt{?list=N} (0--3); if absent, a random list is selected.
Each participant is thus within-subjects across items and sees all four
conditions (twice each). Condition effects are estimated from within-participant
variation across the 2$\times$2 design, controlling for by-participant
random intercepts.

In addition to the 8 critical trials, each participant completed 8
\textbf{filler trials} drawn from separate topic domains (train travel,
reading light, cooking, coffee, vacation, hydration, fresh air, noise).
Filler items used the same trial format but were excluded from the
main analysis; they serve to dilute the critical items and reduce
potential demand characteristics.

\subsubsection{Dependent variable}

The dependent variable was the marker chosen from a drop-down menu
displaying all five options in fixed order (as listed in
Section~\ref{sub:markers}).
No response time threshold was imposed, but response times were logged.
Recorded per trial: \texttt{trial\_id}, \texttt{topic}, \texttt{is\_filler},
\texttt{condition\_index}, \texttt{pc\_prag}, \texttt{g},
\texttt{selected\_marker}, \texttt{rt} (ms); and once per participant:
\texttt{list\_num}.
The continuous \texttt{pc\_prop\_rating} covariate (topic-level norming mean)
is joined to the trial data during analysis.

\subsection{Procedure}

Participants completed the experiment online via the
magpie platform \citep{kochetkova2024magpie} in a browser.
After reading instructions in German, participants were assigned to one
of 4 counterbalancing lists (via URL parameter or randomly).
They then proceeded through 16 trials---8 critical and 8 fillers---in a
fully randomised order, followed by a post-test questionnaire eliciting
native language, age, gender, education level, and optional comments.
Each trial displayed (a) a context paragraph, (b) a goal instruction,
and (c) the sentence frame with the marker drop-down embedded at the
Mittelfeld position.
The ``Weiter'' (Next) button was disabled until a marker had been
selected, ensuring a response on every trial.

\end{document}
   % (inside a \section{Experiment} environment)

\documentclass[12pt]{article}
\usepackage{CSPArticleStyle}
\usepackage{csquotes}
\usepackage[margin=1in]{geometry}
\usepackage{booktabs}
\usepackage{parskip}
\usepackage{amsmath}
\usepackage{times}
\usepackage[T1]{fontenc}
\usepackage[utf8]{inputenc}
\usepackage{xcolor}
\usepackage{tabularray}   % robust tables; fall back to tabular if unavailable
\usepackage{gb4e}         % linguistic examples
\noautomath

\title{Production Experiment: Eliciting Consensus-Marker Choice\\
       Across Perceived-Controversy and Goal-Strength Conditions}
\author{}
\date{}

\begin{document}

%% ============================================================
%% When embedding in main.tex, delete from \documentclass to
%% \begin{document} and delete \end{document} at the bottom.
%% ============================================================

\section{Experiment: Consensus-Marker Production}

\subsection{Overview}

The production experiment empirically estimates the decision thresholds
that govern the use of five German consensus markers by placing
participants in the speaker role.
Participants read short vignettes that independently varied
\emph{pragmatic perceived controversy} (\textsc{pc\textsubscript{prag}})
and the speaker's \emph{persuasive-goal strength} ($g$), then selected
which consensus marker they would most naturally use in the described
situation.
\emph{Propositional perceived controversy} (\textsc{pc\textsubscript{prop}})
is held constant within the experiment and operationalised as a continuous
topic-level covariate obtained from a separate norming study: for each
of the eight topic items, participants rated the degree to which the
proposition $p$ is considered established/accepted (0--100 scale);
the resulting means enter the regression analyses as a fixed covariate.
The resulting distribution of marker choices across the
$\textsc{pc}_\text{prag} \times g$ design, combined with the
continuous \textsc{pc\textsubscript{prop}} covariate, provides the
empirical basis for estimating the ordinal thresholds
$\mu_1 < \cdots < \mu_4$ of the noisy-threshold RSA model.

\subsection{Participants}

Participants will be recruited via Prolific and will be required to be
native speakers of German (language-at-home filter).
Data collection is ongoing; a target of $N = 80$ participants is planned
($n = 20$ per list across 4 lists), each completing 16 trials (8 critical + 8 filler).

\subsection{Materials}

\subsubsection{Markers}

The five consensus markers under investigation are, from weakest to
strongest presumed common-ground commitment:

\begin{enumerate}
  \item \textit{sofern ich wei\ss{}} (`as far as I know')
  \item \textit{wie du ja wei\ss t} (`as you (prt) know')
  \item \textit{wie wir wissen} (`as we know')
  \item \textit{ja} (modal particle; roughly `as you know / of course')
  \item \textit{bekanntlich} (`as is well known / notoriously')
\end{enumerate}

These markers were embedded in the sentential \emph{Mittelfeld}
(middle field), immediately after the finite verb, as illustrated in
example~(\ref{ex:frame}):

\begin{exe}
  \ex \label{ex:frame}
  Digitale Kompetenzen sind \textbf{[MARKER]}
  in der heutigen Arbeitswelt unverzichtbar.\\
  `Digital skills are \textbf{[MARKER]} indispensable in today's
  working world.'
\end{exe}

This syntactic position is licit for all five markers: modal particles
such as \textit{ja} canonically occupy the left edge of the Mittelfeld
\citep{thurmair1989modalpartikeln}, while the adverb \textit{bekanntlich}
and parenthetical clauses (\textit{wie wir wissen}, etc.) can appear
there as well (see \citealt{jacobs_semantics_1991} for discussion).
Each target sentence was followed by a consequent clause with the
recommended action $q$: \textit{Deshalb solltest du [q].}
(`Therefore, you should [q].')

\subsubsection{Topic Items}

Eight topic domains were used, each associated with a fixed proposition
$p$ (sentence frame) and consequent action $q$ (Table~\ref{tab:items}).
The first four items were used in earlier pilot work; items 5--8 were
added to reach the minimum of eight required for a complete Latin square.

\begin{table}[ht]
\centering
\caption{Topic items: sentence frame for $p$ (with marker slot [M]) and consequent $q$.}
\label{tab:items}
\small
\begin{tabular}{clll}
\toprule
\textbf{\#} & \textbf{Topic} & \textbf{Sentence frame} & \textbf{Consequent $q$} \\
\midrule
1 & Klimawandel &
  \textit{Der Klimawandel wird [M]} &
  auf Flugreisen verzichten \\
  & & \textit{durch menschliche Aktivit\"aten verursacht.} & \\[3pt]
2 & Ern\"ahrung &
  \textit{Eine \"uberwiegend pflanzliche Ern\"ahrung} &
  mehr pflanzliche Lebens- \\
  & & \textit{f\"ordert [M] die Gesundheit.} &
  mittel einbauen \\[3pt]
3 & Stadtverkehr &
  \textit{Ein ausgebautes Nahverkehrsnetz} &
  h\"aufiger auf \"offentliche \\
  & & \textit{entlastet [M] den Stadtverkehr deutlich.} &
  Verkehrsmittel umsteigen \\[3pt]
4 & Digitalisierung &
  \textit{Digitale Kompetenzen sind [M] in der} &
  regelm\"a\ss{}ig Weiterbildungs- \\
  & & \textit{heutigen Arbeitswelt unverzichtbar.} &
  angebote nutzen \\[3pt]
5 & Schlaf &
  \textit{Ausreichend Schlaf ist [M] entscheidend} &
  auf feste Schlafzeiten \\
  & & \textit{f\"ur die kognitive Leistungsf\"ahigkeit.} &
  achten \\[3pt]
6 & Plastik &
  \textit{Einwegplastik schadet [M] den} &
  konsequent auf Einweg- \\
  & & \textit{Meeres\"okosystemen erheblich.} &
  plastik verzichten \\[3pt]
7 & Sport &
  \textit{Regelm\"a\ss{}ige k\"orperliche Bewegung} &
  regelm\"a\ss{}ig Sport in \\
  & & \textit{senkt [M] das Risiko f\"ur} &
  den Alltag integrieren \\
  & & \textit{Herz-Kreislauf-Erkrankungen deutlich.} & \\[3pt]
8 & Lokalkauf &
  \textit{Das Kaufen bei lokalen H\"andlern} &
  \"ofter bei lokalen \\
  & & \textit{st\"arkt [M] die regionale Wirtschaft} &
  Gesch\"aften einkaufen \\
  & & \textit{nachhaltig.} & \\
\bottomrule
\end{tabular}
\end{table}

\noindent [M] denotes the marker slot (dropdown).

\subsubsection{Vignettes}

\paragraph{Propositional perceived controversy (\textsc{pc\textsubscript{prop}}).}
\textsc{pc\textsubscript{prop}} was \emph{not} manipulated within the
production experiment.
Instead, it was operationalised as a continuous, topic-level covariate
obtained from a separate norming study conducted prior to the production
experiment.
In that norming study, a different sample of native German speakers rated,
for each topic proposition $p$, the degree to which they considered $p$
accepted/established in (German) public discourse, on a 0--100 scale
($0 = $ `completely contested', $100 = $ `universally accepted').
The resulting per-topic means (e.g., \textit{Digitalisierung}: 85/100;
\textit{Lokalkauf}: 58/100) enter the regression analyses as a
continuous, centred predictor \textsc{pc\textsubscript{prop}}$_c$
(mean-centred, SD $\approx$ 1).
Context paragraphs in the production experiment therefore present
the proposition $p$ without any editorial evaluation, so that
participants' experience of propositional controversy is shaped
solely by their own prior knowledge and the norming-derived covariate
captures this variation in analysis.

\paragraph{Pragmatic perceived controversy (\textsc{pc\textsubscript{prag}}).}
\textsc{pc\textsubscript{prag}} was manipulated between instances of each item
by varying whether the speaker's social circle \emph{largely agreed} (low)
or \emph{held divergent views} (high) about $p$.
The two context versions for the \textit{Digitalisierung} item are:

\begin{itemize}
  \item $\textsc{pc}_\text{prag} = \text{low}$: \textit{Du sprichst mit einer
        Kollegin {\"u}ber Berufsausbildung und Weiterbildung.
        In eurem Team teilen die meisten die Ansicht,
        dass digitale Kompetenzen in der heutigen Arbeitswelt unverzichtbar sind.}
  \item $\textsc{pc}_\text{prag} = \text{high}$: \textit{Du sprichst mit einer
        Kollegin {\"u}ber Berufsausbildung und Weiterbildung.
        In eurem Team gehen die Meinungen dazu stark auseinander.}
\end{itemize}

\paragraph{Speaker goal strength ($g$).}
Goal strength was manipulated via a standardised goal instruction
appended after the context paragraph:

\begin{itemize}
  \item $g_\text{low}$: \textit{Du m\"ochtest das nur kurz anmerken.}
        (`You just want to make a brief note of this.')
  \item $g_\text{high}$: \textit{Es ist dir sehr wichtig, dass dein
        Gespr\"achspartner das akzeptiert.}
        (`It is very important to you that your interlocutor accepts this.')
\end{itemize}

\paragraph{Complete sample item (all four conditions).}
Examples (\ref{ex:cond1})--(\ref{ex:cond4}) show the full vignette for
the topic \textit{Digitalisierung}
(norming $\textsc{pc}_\text{prop}$ rating: 85/100;
sentence frame: \textit{Digitale Kompetenzen sind [M] in der heutigen
Arbeitswelt unverzichtbar.
Deshalb solltest du regelm\"a\ss{}ig digitale Weiterbildungsangebote nutzen.})
across all four conditions of the $2\times2$ design.

\begin{exe}
  \ex \label{ex:cond1}
  \textbf{Condition 0: $\textsc{pc}_\text{prag} = \text{low}$, $g = \text{low}$}\\
  \textit{Du sprichst mit einer Kollegin {\"u}ber Berufsausbildung und Weiterbildung.
  In eurem Team teilen die meisten die Ansicht,
  dass digitale Kompetenzen in der heutigen Arbeitswelt unverzichtbar sind.}\\
  \textit{Du m\"ochtest das nur kurz anmerken.}\\
  Digitale Kompetenzen sind \underline{\hspace{2.2cm}}
  in der heutigen Arbeitswelt unverzichtbar.

  \ex \label{ex:cond2}
  \textbf{Condition 1: $\textsc{pc}_\text{prag} = \text{low}$, $g = \text{high}$}\\
  \textit{Du sprichst mit einer Kollegin {\"u}ber Berufsausbildung und Weiterbildung.
  In eurem Team teilen die meisten die Ansicht,
  dass digitale Kompetenzen in der heutigen Arbeitswelt unverzichtbar sind.}\\
  \textit{Es ist dir sehr wichtig, dass dein Gespr\"achspartner das akzeptiert.}\\
  Digitale Kompetenzen sind \underline{\hspace{2.2cm}}
  in der heutigen Arbeitswelt unverzichtbar.

  \ex \label{ex:cond3}
  \textbf{Condition 2: $\textsc{pc}_\text{prag} = \text{high}$, $g = \text{low}$}\\
  \textit{Du sprichst mit einer Kollegin {\"u}ber Berufsausbildung und Weiterbildung.
  In eurem Team gehen die Meinungen dazu stark auseinander.}\\
  \textit{Du m\"ochtest das nur kurz anmerken.}\\
  Digitale Kompetenzen sind \underline{\hspace{2.2cm}}
  in der heutigen Arbeitswelt unverzichtbar.

  \ex \label{ex:cond4}
  \textbf{Condition 3: $\textsc{pc}_\text{prag} = \text{high}$, $g = \text{high}$}\\
  \textit{Du sprichst mit einer Kollegin {\"u}ber Berufsausbildung und Weiterbildung.
  In eurem Team gehen die Meinungen dazu stark auseinander.}\\
  \textit{Es ist dir sehr wichtig, dass dein Gespr\"achspartner das akzeptiert.}\\
  Digitale Kompetenzen sind \underline{\hspace{2.2cm}}
  in der heutigen Arbeitswelt unverzichtbar.
\end{exe}

\noindent In all four vignettes the sentence frame and consequent clause
(\textit{Deshalb solltest du regelm\"a\ss{}ig digitale
Weiterbildungsangebote nutzen.}) are identical; only the context
paragraph and goal instruction vary across conditions.
The marker slot [M] / \underline{\hspace{2.2cm}} was presented as a
drop-down menu in the actual experiment.

\subsection{Design}

\subsubsection{Factorial structure}

The experiment employed a \textbf{2 (pc\textsubscript{prag}: low vs.\ high)
$\times$ 2 ($g$: low vs.\ high) fully crossed design},
yielding 4 conditions.
\textsc{pc\textsubscript{prop}} is not a within-experiment factor; it enters
as a continuous topic-level covariate from the norming study.
With 8 topic items and a Latin square rotation across 4 lists, each
participant saw exactly \textbf{8 critical trials}---two items per condition---plus
8 filler trials, for \textbf{16 trials} in total
(Table~\ref{tab:design}).

\begin{table}[ht]
\centering
\caption{%
  The $2\times2$ design (4 conditions).
  \textsc{pc\textsubscript{prop}} is a continuous norming covariate, not a
  within-experiment factor.
  In each list, the 8 items are assigned to conditions via the Latin square
  rotation described in Section~\ref{subsub:counterbalancing}.}
\label{tab:design}
\small
\begin{tabular}{ccl}
\toprule
\textbf{pc\textsubscript{prag}} &
\textbf{$g$} &
\textbf{Condition label} \\
\midrule
low  & low  & (0) community agrees, casual \\
low  & high & (1) community agrees, pressing \\
high & low  & (2) community divided, casual \\
high & high & (3) community divided, pressing \\
\bottomrule
\end{tabular}
\end{table}

\subsubsection{Counterbalancing}
\label{subsub:counterbalancing}

The experiment used a \textbf{Latin square} to counterbalance item--condition
pairings across participants.
With $k = 4$ conditions and 8 items, a balanced assignment is achieved:
each of the 4 lists assigns item $i$ ($i = 0,\ldots,7$) to condition
$(i + l)\bmod 4$, where $l \in \{0,1,2,3\}$ is the list number.
The resulting assignment matrix (Table~\ref{tab:ls}) guarantees that
\begin{itemize}
  \item each participant sees every item exactly once;
  \item each participant sees each condition exactly twice (two items per
        condition), ensuring balanced coverage of the 2$\times$2 design;
  \item across all 4 lists, each item appears equally often in each
        condition (twice per condition over the full participant sample).
\end{itemize}

\begin{table}[ht]
\centering
\caption{Latin square assignment of conditions to items across 4 lists.
  Cell entries are condition indices (0--3) as defined in
  Table~\ref{tab:design}.
  Rows = lists ($l$); columns = items ($i$).}
\label{tab:ls}
\small
\begin{tabular}{ccccccccc}
\toprule
\textbf{List} &
\textbf{Item 1} & \textbf{Item 2} & \textbf{Item 3} & \textbf{Item 4} &
\textbf{Item 5} & \textbf{Item 6} & \textbf{Item 7} & \textbf{Item 8} \\
\midrule
0 & 0 & 1 & 2 & 3 & 0 & 1 & 2 & 3 \\
1 & 1 & 2 & 3 & 0 & 1 & 2 & 3 & 0 \\
2 & 2 & 3 & 0 & 1 & 2 & 3 & 0 & 1 \\
3 & 3 & 0 & 1 & 2 & 3 & 0 & 1 & 2 \\
\bottomrule
\end{tabular}
\end{table}

\noindent List assignment is determined automatically from the URL query
parameter \texttt{?list=N} (0--3); if absent, a random list is selected.
Each participant is thus within-subjects across items and sees all four
conditions (twice each). Condition effects are estimated from within-participant
variation across the 2$\times$2 design, controlling for by-participant
random intercepts.

In addition to the 8 critical trials, each participant completed 8
\textbf{filler trials} drawn from separate topic domains (train travel,
reading light, cooking, coffee, vacation, hydration, fresh air, noise).
Filler items used the same trial format but were excluded from the
main analysis; they serve to dilute the critical items and reduce
potential demand characteristics.

\subsubsection{Dependent variable}

The dependent variable was the marker chosen from a drop-down menu
displaying all five options in fixed order (as listed in
Section~\ref{sub:markers}).
No response time threshold was imposed, but response times were logged.
Recorded per trial: \texttt{trial\_id}, \texttt{topic}, \texttt{is\_filler},
\texttt{condition\_index}, \texttt{pc\_prag}, \texttt{g},
\texttt{selected\_marker}, \texttt{rt} (ms); and once per participant:
\texttt{list\_num}.
The continuous \texttt{pc\_prop\_rating} covariate (topic-level norming mean)
is joined to the trial data during analysis.

\subsection{Procedure}

Participants completed the experiment online via the
magpie platform \citep{kochetkova2024magpie} in a browser.
After reading instructions in German, participants were assigned to one
of 4 counterbalancing lists (via URL parameter or randomly).
They then proceeded through 16 trials---8 critical and 8 fillers---in a
fully randomised order, followed by a post-test questionnaire eliciting
native language, age, gender, education level, and optional comments.
Each trial displayed (a) a context paragraph, (b) a goal instruction,
and (c) the sentence frame with the marker drop-down embedded at the
Mittelfeld position.
The ``Weiter'' (Next) button was disabled until a marker had been
selected, ensuring a response on every trial.

\end{document}
